\section{Simultaneous primes with coprime adjacent norms}
\label{sec:v5-sun-proof}

This appendix gives the prime-pattern argument needed for the construction.
Its strategy is due to Sun~\cite{v5sun}. We supply the local calculations,
the two ranges of the sieve, the normalization of the ray-class densities,
and the order in which the parameters are chosen. For the analytic input
we prove the special case for fixed ray class fields of $\Q(i)$ directly.
The ideal-norm method follows Frei and Tinkov\'a~\cite{v5ft}; no theorem
from that work is assumed. The proof uses the published fixed-field prime
estimate of Mitsui, the one-parameter equidistribution and factorization
results of Green--Tao, and published ordinary-prime transference results.
The field and its ray modulus are fixed before every asymptotic limit.
The exact source versions checked are identified in the bibliography
and the accompanying theorem-application inventory.

Write $\mathrm N(z)=z\bar z$ for $z\in\mathbb Z[i]$, and put
$d_-(z)=\mathrm N(z-1)$ and $d_+(z)=\mathrm N(z+1)$.
An Artin or Frobenius element in this appendix is arithmetic Frobenius:
on residue fields it sends $x$ to $x^{\mathrm N\mathfrak p}$.

\begin{theorem}[The simultaneous prime input]\label{thm:v5-simultaneous-primes}
Let $c_1,\ldots,c_k$ be distinct integers, where $k\geq2$.
Let $q\geq1$, let $M$ be an odd squarefree positive integer with $7\nmid M$,
and let $H,Y>0$. There are positive integers $u,w$ such that
\[
 q\mid w,\qquad u/w^2>H,\qquad p_j=u+c_jw^2>Y
 \quad(1\leq j\leq k),
\]
where every $p_j$ is prime. Each $p_j$ has a Gaussian prime generator
$\pi_j$ of norm $p_j$ such that
\[
 \pi_j\equiv9\pmod {16},\qquad
 \pi_j\equiv1+3i\pmod7,\qquad
 \pi_j\equiv i\pmod M.
\]
No odd prime divides both $d_-(\pi_j)$ and $d_+(\pi_j)$.
For $j\ne h$, no odd prime divides both $d_\epsilon(\pi_j)$ and
$d_\eta(\pi_h)$, for any $\epsilon,\eta\in\{-,+\}$.
Moreover,
\[
 v_2(d_-(\pi_j))=6,\quad v_2(d_+(\pi_j))=2,\quad
 \gcd(d_-(\pi_j),d_+(\pi_h))=4\quad\hbox{for all }j,h,
\]
and all these norms are coprime to $7M$.
\end{theorem}

The quantifiers concern fixed data. No estimate uniform in $q,M,k$ or the
$c_j$ is asserted. We first prove the result for $q=1$. Replacing the
coefficients by $q^2c_j$ and then replacing the parameter $v$ by $w=qv$
proves the general case, since the ratio before this substitution can be
made greater than $q^2H$.

\subsection{Local counts and the medium primes}\label{subsec:v5-local-sieve}

Fix a nonzero integer $\kappa$ and write
\[
 Q_\kappa=a^2+b^2-c^2-d^2-\kappa w^2,\quad
 G_\epsilon=(a+\epsilon)^2+b^2,\quad
 H_\eta=(c+\eta)^2+d^2,
 \qquad \epsilon,\eta\in\{1,-1\}.
\]
We use $\chi(\ell)=(-1/\ell)$ for odd primes $\ell$.

\begin{lemma}[Finite-field counts]\label{lem:v5-local-counts}
If $\ell\nmid2\kappa$, then $Q_\kappa=0$ has $\ell^4-1$ nonzero
solutions over $\mathbb F_\ell$. At most $4\ell^2$ of its solutions satisfy
$G_\epsilon=H_\eta=0$. If $\chi(\ell)=-1$, the latter set has exactly one
element.

Let $\gamma_1,\ldots,\gamma_k$ be distinct integers, including $0$ and
$-\kappa$, and put $L_j=a^2+b^2+\gamma_jw^2$. Outside a fixed set of primes
containing the divisors of $2\kappa\prod_{j<h}(\gamma_j-\gamma_h)$, the
proportion $g(\ell)$ of nonzero zeros of $Q_\kappa$ on which some $L_j$
vanishes is
\begin{equation}\label{eq:v5-local-g}
 g(\ell)=
 \frac{(\ell-1)(\ell-\chi(\ell))
 [2\nu_0+(k-2)(\ell-\chi(\ell))]+\nu_0^2-1}{\ell^4-1},
 \qquad \nu_0=\ell+(\ell-1)\chi(\ell).
\end{equation}
Consequently $g(\ell)=(k+2\chi(\ell))/\ell+O_k(\ell^{-2})$.
\end{lemma}
\begin{proof}
The norm $x^2+y^2$ has $\nu_0$ representations of zero and
$\ell-\chi(\ell)$ representations of every nonzero element. This follows
either by factoring over $\mathbb F_\ell$ when $\chi=1$, or by the norm map
from $\mathbb F_{\ell^2}$ when $\chi=-1$. Summing these fibres for
$Q_\kappa$ gives $\ell^4$ solutions including the origin.
In the split case, each of the equations $G_\epsilon=0$ and $H_\eta=0$
is a union of two lines. On each of the four pairs of lines the remaining
equation is
$-2\epsilon a+2\eta c-\kappa w^2=0$, leaving at most $\ell^2$ choices.
In the inert case these equations force
$(a,b,c,d)=(-\epsilon,0,-\eta,0)$, and then $w=0$.

For $w\ne0$, the conditions $L_j=0$ are disjoint. For the two distinguished
coefficients one norm is zero and the other is nonzero, giving
$2\nu_0(\ell-\chi)$ choices. Each other coefficient gives
$(\ell-\chi)^2$ choices. There are $\ell-1$ possible nonzero $w$.
For $w=0$, both norms are zero, giving $\nu_0^2-1$ nonzero vectors.
This proves \eqref{eq:v5-local-g}. In particular, when $k=2$ and
$\chi=-1$, its value is $2/(\ell^2+1)$, not zero.
\end{proof}

\begin{lemma}[Primitive congruence count]\label{lem:v5-primitive-count}
Let $\psi$ be a smooth nonnegative compactly supported function away from
the origin, and let $m$ be squarefree and prime to $2\kappa$.
For each $\ell\mid m$ prescribe an arbitrary subset $\Omega_\ell$ of the
nonzero zeros of $Q_\kappa$ over $\mathbb F_\ell$. The weighted number of
primitive integral vectors on $Q_\kappa=0$ with these reductions equals
\begin{equation}\label{eq:v5-primitive-count}
 C_\kappa\sigma_\infty(\psi)B^3
 \prod_{\ell\mid m}\frac{|\Omega_\ell|}{\ell^4-1}
 +O_{\kappa,\psi,\varepsilon}
       (B^{5/2+\varepsilon}m^{5/2+\varepsilon}).
\end{equation}
The constant $C_\kappa$ is positive and independent of $m$ and the subsets.
\end{lemma}
\begin{proof}
Apply Browning--Loughran~\cite{v5bl}, Theorem~4.1, to $Q_\kappa$ in five
variables. It is isotropic, since $(1,0,1,0,0)$ is a zero, and nonsingular,
since $\kappa\ne0$. Its discriminant has no prime divisor outside
$2\kappa$. The theorem permits $n=5$ and arbitrary nonzero residue subsets;
its error is uniform in the modulus and in these subsets.
Impose primitivity by $\sum_{h\mid(a,b,c,d,w)}\mu(h)$.
If $(h,m)>1$ there are no permitted vectors. Otherwise the subsets become
$h^{-1}\Omega_\ell$, of the same cardinalities, and the scale becomes $B/h$.
The sum of the error terms is convergent. Extending the main-term sum to
infinity costs $O(B)$.
At a good prime the unrestricted local density is
\[
 \sigma_\ell=\frac{1-\ell^{-4}}{1-\ell^{-3}}.
\]
Indeed, nonsingular primitive zeros lift with four free coordinates, and
the vectors divisible by $\ell$ contribute the geometric factor
$\ell^{-3}$. The primitive density is therefore $1-\ell^{-4}$.
This turns the factor $|\Omega_\ell|/\ell^4$ of the unprimitive count into
$|\Omega_\ell|/(\ell^4-1)$ in \eqref{eq:v5-primitive-count}.
The remaining finitely many local factors are positive because the form
has a nonsingular rational zero. Empty subsets give the assertion directly.
\end{proof}

In using this count, the two norms lie between fixed positive multiples of
$B^2$, and $|w|\leq\tau B$, where $0<\tau\leq\tau_0$.
A smooth majorant can be chosen with
$\sigma_\infty(\psi)\ll\tau$ uniformly in $\tau$: integrate first in the
two radial norm variables, on which the derivative of their difference is
bounded away from zero. The constants in error terms may depend on
derivatives of this majorant, and hence on $\tau$.

\begin{lemma}[Medium-prime bound]\label{lem:v5-medium-sieve}
For fixed coefficients and sufficiently large fixed $L$, the contribution
of prime patterns with a common adjacent-norm divisor
$L<\ell\leq B^{1/100}$ is at most
\begin{equation}\label{eq:v5-medium-sieve}
 \frac{C\tau B^3}{(\log B)^k L\log L}+O_\tau(B^{27/10}).
\end{equation}
The count is unweighted by logarithms, and $C$ is independent of $L,\tau$.
The bound covers both distinct positions and the two signs at one position.
\end{lemma}
\begin{proof}
For two positions, use $Q_\kappa$ with their fixed nonzero coefficient
difference and impose $G_\epsilon=H_\eta=0\pmod\ell$.
For squarefree $d$ supported on good primes other than $\ell$, impose also
that some $L_j$ vanishes at every prime dividing $d$. The preceding lemmas
give a sieve sequence
\[
 A_{\ell,d}=X_\ell g(d)+r_{\ell,d},\qquad
 X_\ell\ll\tau B^3/\ell^2,\qquad
 |r_{\ell,d}|\ll_{\tau,\varepsilon}
 B^{5/2+\varepsilon}(\ell d)^{5/2+\varepsilon}.
\]
Enlarge the fixed exceptional set so that $0<g(p)<1$ outside it.
Set $h(p)=g(p)/(1-g(p))$ and extend $h$ multiplicatively to squarefree
integers. The elementary Selberg upper sieve has denominator
\[
 \mathcal G_\ell(z)=\sum_{\substack{d<z\ \text{squarefree}\\
                         (d,\ell\prod_{p\in S}p)=1}}h(d)
 \gg (\log z)^k.
\]
Here is the dimension calculation. Williams~\cite{v5williams}, Theorem~1
with fixed modulus $4$ and residues $1,3$, gives Mertens' estimates in the
two reduced classes. Taking logarithms gives
$\sum_{p\leq y}g(p)=k\log\log y+O(1)$ and
$\sum_{p\leq y}g(p)\log p\ll\log y$.
For the latter bound only the elementary Chebyshev estimate is needed:
$\prod_{n<p\leq2n}p\mid\binom{2n}{n}<4^n$ gives
$\vartheta(x)\ll x$ by dyadic summation, and partial summation gives
$\sum_{p\leq y}(\log p)/p\ll\log y$.
In the Euler product $\prod_{p\leq y}(1+h(p))$, normalize the weights
$h(d)$ to a probability distribution. Its expected value of $\log d$ is
$\sum_{p\leq y}g(p)\log p$. Taking $y$ to be a sufficiently small fixed
power of $z$ leaves at least half of the mass on $d<z$, by Markov's
inequality. Removing one prime $\ell$ changes the estimate by a bounded
factor.

For completeness the sieve weights may be taken as
\[
 \lambda_d=\frac{\mu(d)}{\mathcal G_\ell(z)}
 \prod_{p\mid d}(1-g(p))^{-1}
 \sum_{\substack{r<z/d\ (r,d)=1\ r\ \text{squarefree}\\
                      (r,\ell\prod_{p\in S}p)=1}}h(r).
\]
They have $\lambda_1=1$, $|\lambda_d|\ll_\varepsilon d^\varepsilon$,
and direct expansion gives
$\sum_{d,e}\lambda_d\lambda_e g([d,e])=1/\mathcal G_\ell(z)$.
Thus the number with no prime divisor below $z$ in any $L_j$ is at most
\[
 \frac{C\tau B^3}{\ell^2(\log z)^k}
 +O_{\tau,\varepsilon}
  (B^{5/2+\varepsilon}\ell^{5/2+\varepsilon}z^{7+4\varepsilon}).
\]
All the prime values under consideration exceed $z$. Put $z=B^{1/100}$
and sum over $L<\ell\leq B^{1/100}$. The main term uses
$\sum_{\ell>L}\ell^{-2}\ll1/(L\log L)$, another consequence of the
same Chebyshev bound; for example
$\varepsilon=1/1000$ makes the error smaller than $B^{27/10}$.

At one position, an odd common divisor of $(a-1)^2+b^2$ and
$(a+1)^2+b^2$ forces $a=0$ and $b^2=-1\pmod\ell$.
There are at most $2\ell$ such triples $(a,b,w)$ over $\mathbb F_\ell$.
For $w\ne0$ the zero sets of the $k$ prime forms are disjoint, and the
ordinary three-dimensional local sieve density is
$(k+\chi(\ell))/\ell+O_k(\ell^{-2})$.
Counting residue classes in a smooth box gives the corresponding
three-dimensional congruence estimate. The same Selberg weights and the
same, weaker error in \eqref{eq:v5-medium-sieve} apply.
\end{proof}

\subsection{The large-prime tail}\label{subsec:v5-large-sieve}

We use Lemmas~2.1 and 4.1 of Browning--Heath-Brown~\cite{v5bhb}.
The adjacent-norm polynomials are inhomogeneous, so we give the
fixed-coordinate argument rather than invoke their homogeneous
quadric-sieve theorem directly.
Their Lemma~2.1 says that coprime polynomials in $\mathbb Z[X_1,\ldots,X_m]$
of bounded degree and coefficient height at most $H$ have at most
\begin{equation}\label{eq:v5-box-sieve}
 O\left(\frac{V\log(VH)}{L\log L}
           +\frac{V\log(VH)}{B_{\min}}\right)
\end{equation}
integer points in a box of side parameters $B_i\geq2$ at which a prime
greater than $L\geq2$ divides every polynomial; $V=\prod B_i$.
``Coprime'' here excludes a common constant prime factor as well as a
common nonconstant factor. Their Lemma~4.1 covers the zeros modulo a
squarefree $q$ of a rank-$m$ quadratic form, with $m\geq3$, by at most
$(3m)^{\omega(q)}q^{m-2}$ lattices of determinant $q^{m-1}$.
The number with largest Euclidean successive minimum exceeding $T$ is
\begin{equation}\label{eq:v5-lattice-tail}
 O_m((3m)^{\omega(q)}q^{2m-3}T^{-m}).
\end{equation}
Every prime dividing $q$ must be odd and preserve the rank of the form.
The lattices in that lemma contain $q\mathbb Z^m$; in particular their
largest successive minima are at most $q$.

\begin{lemma}[Large common divisors on the quadric]\label{lem:v5-large-sieve}
Fix the implied constant in the box of size $O(B)$.
For fixed $\kappa\ne0$ and $\epsilon,\eta\in\{1,-1\}$, uniformly for
$B,L\geq2$, the number of integral vectors in this box satisfying
$Q_\kappa=0$ and
$\ell\mid G_\epsilon,H_\eta$ for some prime $\ell>L$ is
\begin{equation}\label{eq:v5-large-sieve}
 O_{\kappa,\alpha}\left(
    \frac{B^{3+\alpha}}{L\log L}+B^{11/4+\alpha}\right)
 \qquad(\alpha>0).
\end{equation}
\end{lemma}
\begin{proof}
Sign changes of $a,c$ reduce to $\epsilon=\eta=1$.
Put $x_0=a+c$, $x_1=a-c$, $x=b-d$, $z=b+d$, and
$R=\kappa w^2-xz$. Then $x_0x_1=R$.
Parity restrictions can be discarded in an upper bound.
For a fixed nonzero value of either $x_0$ or $x_1$, the triple $(x,z,w)$
determines the other coordinate uniquely whenever it gives an integral
solution. We shall use this injection when discarding divisibility
conditions.

First suppose that $\ell\mid x_0,x_1$. Then $\ell^2\mid R$.
For each fixed integer $v=O(B^2)$, the equation $R=v$ has
$O_{\kappa,\alpha}(B^{1+\alpha})$ triples: for each $w$, count the
factorizations $xz=\kappa w^2-v$ by the divisor bound; the zero right-hand
side occurs for at most two $w$ and contributes $O(B)$ each.
If $R\ne0$, it has $O_\alpha(B^\alpha)$ factorizations $x_0x_1=R$.
Since $0<|h|\ll B^2/\ell^2$ and $\ell\ll B$, summation gives
\[
 O_{\kappa,\alpha}\left(
 B^{3+2\alpha}\sum_{\ell>L}\ell^{-2}\right)
 =O_{\kappa,\alpha}(B^{3+2\alpha}/(L\log L)).
\]
The same bound for triples with $R=0$, followed by the $O(B)$ choices
with $x_0x_1=0$, gives $O_{\kappa,\alpha}(B^{2+\alpha})$.
We may discard all these $R=0$ solutions before proceeding.

In the remaining cases choose $t=x_1$ or $t=x_0$ such that $\ell\nmid2t$.
Take $t=x_1$ first. Here
$G_1-H_1=\kappa w^2+2t$ and $4t^2G_1=K$.
Thus the two norm congruences are equivalent to divisibility of
\[
 P=\kappa w^2+2t,\qquad
 K=(\kappa w^2-xz+t^2+2t)^2+t^2(x+z)^2.
\]
The identity
\[
 K-(t^2+x^2)(t^2+z^2)=P(P+2t^2-2xz)
\]
shows that it suffices to sieve the coprime pair
$P,(t^2+x^2)(t^2+z^2)$. They have no common nonconstant factor over
$\mathbb Q$, because the first depends nontrivially only on $w$, whereas
the second is independent of $w$. Their contents have no common prime,
since the second has coefficient $1$ on $x^2z^2$.
For $t=x_0$ we still have $4t^2G_1=K$, whereas
$t(G_1-H_1)=\kappa(t+2)w^2-2xz$. The pair is therefore
\[
 P'=\kappa(t+2)w^2-2xz,\qquad K.
\]
If $t\ne-2$, $P'$ is an irreducible ternary quadratic form of rank three
and $K$ has nonzero constant term $t^2(t+2)^2$; homogeneity rules out
divisibility. If $t=-2$, neither $x$ nor $z$ divides $K$.
The coefficient $1$ of $x^2z^2$ in $K$ excludes every common content
prime. In both charts the fixed value $t$ is nonzero.

In either chart, $t\mid R$. Put
\[
 S_R=\{p:p\mid2\kappa\},\qquad
 q=\prod_{\substack{p\mid t\ p\notin S_R}}p.
\]
The matrix of $R$ has determinant $-\kappa/4$, so $R$ has rank three
at every prime in this product, and those primes are odd. Also $q\ll B$.
For a fixed $q$, every eligible nonzero $t=O(B)$ is supported on
$S_R\cup\{p:p\mid q\}$. For every $\delta>0$, Rankin's bound gives
\[
 \#\{t:q(t)=q,\ 0<|t|\ll B\}
 \ll B^\delta
   \prod_{p\in S_R\ \text{or}\ p\mid q}(1-p^{-\delta})^{-1}
 \ll_{\kappa,\delta} B^\delta q^\delta
 \ll_{\kappa,\delta} B^{2\delta}.
\]
For the middle inequality, remove the finitely many primes with
$p^{-\delta}>1/2$ and bound the remaining factors by $2^{\omega(q)}
\ll_\delta q^\delta$. Renaming the small exponent gives
$O_{\kappa,\alpha}(B^\alpha)$ choices of $t$.

Apply \eqref{eq:v5-lattice-tail} with $m=3$ to $q\mid R$.
If $\lambda_1\leq\lambda_2\leq\lambda_3$ are the successive minima,
then $\lambda_3\leq q\ll B$ and
$\lambda_1\lambda_2\lambda_3\asymp q^2$.
A basis comparable to these minima therefore gives
$O(\prod_i(1+B/\lambda_i))=O(B^3/q^2)$ points in the box.
The lattices with $\lambda_3>T$ contribute at most
\[
 B^\alpha\sum_{q\ll B}
 O\left(9^{\omega(q)}q^3T^{-3}\frac{B^3}{q^2}\right)
 \ll_{\kappa,\alpha} B^{5+2\alpha}T^{-3}.
\]
For the other lattices choose a basis with lengths comparable to the
successive minima. The inverse basis matrix bounds the new coordinates by
$O(B/\lambda_i)$, so the resulting rectangular box has volume
$O(B^3/q^2)$ and minimum side parameter $\gg B/T$.
Enlarge its side parameters to at least $2$, as required in
\eqref{eq:v5-box-sieve}; since $\lambda_i\ll B$, this changes volume
only by a fixed factor.
The substitution into the polynomial pair is invertible over $\mathbb Q$.
For $\ell\nmid2t$ its determinant, of absolute value $q^2$, is invertible
modulo $\ell$ too. Thus no new common content prime outside $2t$ appears.
Remove common content supported on $2t$, which cannot affect the primes
being counted. Coefficient heights are powers of $B$ with fixed exponents.
Formula \eqref{eq:v5-box-sieve} bounds this lattice's contribution by
\[
 O\left(\frac{B^3\log B}{q^2L\log L}
                 +\frac{B^2T\log B}{q^2}\right).
\]
There are at most $9^{\omega(q)}q$ lattices. Since
$9^{\omega(q)}\ll_\alpha q^\alpha$, summing over $q\ll B$ and the
$O_{\kappa,\alpha}(B^\alpha)$ values of $t$ gives
$O_{\kappa,\alpha}(B^{3+3\alpha}/(L\log L)+B^{2+3\alpha}T)$.
Choose $T=B^{3/4}$ and rename the small exponent to obtain
\eqref{eq:v5-large-sieve}.
\end{proof}

\begin{proposition}[Weighted obstruction tail]\label{prop:v5-obstruction-tail}
For fixed distinct $c_j$ there are $L_*\geq2$ and $\tau_0,C>0$
such that, for every fixed $0<\tau\leq\tau_0$, uniformly for $L\geq L_*$,
the logarithmically weighted number of pairs
\[
 N/2\leq u\leq N,\qquad 1\leq w\leq\tau\sqrt N,
 \qquad p_j=u+c_jw^2\text{ all prime},
\]
having any of the odd common divisors forbidden in
Theorem~\ref{thm:v5-simultaneous-primes} greater than $L$, is at most
\begin{equation}\label{eq:v5-obstruction-tail}
 \frac{C\tau N^{3/2}}{L\log L}
       +O_\tau(N^{3/2-1/800}).
\end{equation}
The weight is $\prod_j\log p_j$, and the bound allows every Gaussian
generator, so it also holds for any prescribed choice of generators.
\end{proposition}
\begin{proof}
Take $B=\sqrt N$ and $L_0=B^{1/100}$, and choose $\tau_0$ so that
every prime value in the displayed region is between fixed positive
multiples of $B^2$. Thus $\prod_j\log p_j\ll (\log B)^k$.
Choose $L_*$ above the fixed exceptional primes in the medium-prime bound.
When $L\leq L_0$, that bound gives
\[
 \frac{C\tau B^3}{L\log L}+O_\tau(B^{27/10}(\log B)^k)
\]
for $L<\ell\leq L_0$. Its proof is uniform in this interval of $L$:
the remainder sum is bounded by the sum over all $\ell\leq L_0$.
When $L>L_0$ this part is empty.

For distinct positions, apply Lemma~\ref{lem:v5-large-sieve} at $L_0$
with $\alpha=1/1000$. After multiplication by the weight, the bound is
\[
 O\left(\left(
   \frac{B^{3+1/1000}}{B^{1/100}\log L_0}
       +B^{11/4+1/1000}\right)(\log B)^k\right)
 =O(B^{3-1/200}).
\]
This bounds all the remaining primes, also when $L>L_0$.
At one position, an odd divisor of both adjacent norms divides
$a$ and $b^2+1$. The coprime pair $a,b^2+1$ in variables $(a,b,w)$,
inserted in \eqref{eq:v5-box-sieve} at $L_0$, gives
\[
 O\left(\left(
 \frac{B^3\log B}{L_0\log L_0}+B^2\log B\right)(\log B)^k\right)
 =O(B^{3-1/200}).
\]
No factor $\tau$ is needed in these power-saving errors.

For a fixed pair of positions, the vector $(a,b,c,d,w)$ determines
$u=a^2+b^2-c_jw^2$; at one position $(a,b,w)$ does so.
Thus the vector counts bound the pattern counts. There are finitely many
pairs of positions and signs. Each split rational prime has eight
Gaussian elements of that norm, so counting all generator choices
changes the bound by at most a fixed factor. Every vector used in the
medium-prime count is primitive, since a common integer divisor greater
than one would have its square divide a prime norm.
All logarithmic factors in the remainders are absorbed for sufficiently
large $B$, with thresholds depending only on the fixed data and $\tau$.
Weakening the error to $B^{3-1/400}=N^{3/2-1/800}$ proves the claim.
\end{proof}

\subsection{Gaussian ray weights and their progression main terms}
\label{subsec:gaussian-progression-mainterms}

Let $E=\Q(i)$ and let $H$ be a fixed positive integer divisible by $16$.
Let $K/E$ be the ray class field for $H\Z[i]$.  This field is Galois over
$\Q$, since complex conjugation preserves both the modulus and its
principal ray subgroup.  Put
\[
 A=\operatorname{Gal}(K/E)
   \simeq (\Z[i]/H\Z[i])^\times/\{1,-1,i,-i\},
 \qquad G=\operatorname{Gal}(K/\Q).
\]
We use arithmetic Frobenius throughout.  Let $C\subseteq A$ be invariant
under complex conjugation, so that $C$ is a union of conjugacy classes in
$G$.  Define $\Lambda'_C(p)=\log p$ for unramified rational primes whose
Frobenius lies in $C$, and put $\Lambda'_C(n)=0$ otherwise.  In particular
this function is supported on rational primes which split in $E$.

For ideals coprime to $H$ let $[\mathfrak a]\in A$ denote their ray class,
and set
\[
 T_C(n)=\sum_{\substack{\mathrm N\mathfrak a=n\\(\mathfrak a,H)=1}}
       \Lambda_E(\mathfrak a)1_C([\mathfrak a]).
\]
Here $\Lambda_E(\mathfrak p^j)=\log\mathrm N\mathfrak p$ and it is zero
on all other ideals.  Every unramified split rational prime has two
conjugate prime ideals of norm $p$, so
\begin{equation}\label{eq:gaussian-ray-prime-factor}
 T_C(p)=2\Lambda'_C(p)\qquad(p\nmid2H).
\end{equation}
The remaining difference between $T_C$ and $2\Lambda'_C$ is supported on
rational prime powers of exponent at least two and the finitely many
excluded primes.  Its total absolute value up to $X$ is
$O_H(\sqrt X\log^2(2X))$.  This bound follows by summing the at most two
prime ideals above each rational prime and their powers; no density
theorem is used in this comparison.

For $\chi\in\widehat A$ put
\[
 c_C(\chi)=|A|^{-1}\sum_{a\in C}\overline{\chi(a)},\qquad
 T_\chi(n)=\sum_{\mathrm N\mathfrak a=n}
                  \Lambda_E(\mathfrak a)\chi(\mathfrak a),
\]
where each character is extended by zero on ideals not coprime to $H$.
Finite Fourier inversion gives the exact identity
\begin{equation}\label{eq:gaussian-ray-fourier}
                 T_C(n)=\sum_{\chi\in\widehat A}c_C(\chi)T_\chi(n).
\end{equation}
The field and all these characters are fixed before $X$ tends to infinity.

\begin{lemma}[Fixed-field Hecke Siegel--Walfisz estimate]
\label{lem:gaussian-hecke-sw}
Let $E=\Q(i)$, let $\chi$ be a fixed finite-order ray character of $E$
with finite modulus $\mathfrak m$, and let $B>0$.
There is $c=c(\chi,B)>0$ such that, uniformly for $X\geq3$,
$1\leq Y\leq X$, positive integers $q\leq(\log X)^B$, and Dirichlet
characters $\rho$ modulo $q$,
\begin{equation}\label{eq:gaussian-hecke-sw-estimate}
 \sum_{\mathrm N\mathfrak a\leq Y}
   \Lambda_E(\mathfrak a)\chi(\mathfrak a)\rho(\mathrm N\mathfrak a)
 =r_{\chi,\rho}Y+
 O_{\chi,B}\bigl(X\exp(-c\sqrt{\log X})\bigr).
\end{equation}
Here $\Lambda_E(\mathfrak p^j)=\log\mathrm N\mathfrak p$ for a prime
ideal $\mathfrak p$ and $j\geq1$, and it is zero otherwise; characters
are extended by zero on ideals not coprime to their moduli.
The number $r_{\chi,\rho}$ is one if the primitive Hecke character
inducing $\chi(\rho\circ\mathrm N)$ is trivial, and zero otherwise.
There is also $c_B>0$ such that, with the same ranges of $X,Y,q,\rho$,
\begin{equation}\label{eq:rational-sw-estimate}
 \sum_{n\leq Y}\Lambda(n)\rho(n)
 =1_{\rho\ \mathrm{principal}}Y+
 O_B\bigl(X\exp(-c_B\sqrt{\log X})\bigr).
\end{equation}
In particular, either error is $O(X(\log X)^{-D})$ for every fixed
$D>0$, with the corresponding dependence on $D$ allowed.
\end{lemma}

\begin{proof}
We use the published fixed-field estimate of Mitsui
\cite[\S1, Lemma~15, p.~12]{Mitsui1956}.
In his notation, for a fixed number field $K$, a fixed $A>0$, and any
ray character $\vartheta$ whose finite modulus has norm at most
$(\log x)^A$, it states
\begin{equation}\label{eq:mitsui-weighted-primes}
 \sum_{\mathrm N\mathfrak p\leq x}
       \vartheta(\mathfrak p)\log\mathrm N\mathfrak p
 =E(\vartheta)x+
 O_{K,A}\bigl(x\exp(-c_{K,A}\sqrt{\log x})\bigr),
\end{equation}
where $E(\vartheta)=1$ for the principal character and is zero otherwise.
The sum is over prime ideals of every residue degree.
Mitsui defines these characters and $E(\vartheta)$ on p.~2, includes
imprimitive characters, and states the uniformity in the character
and its modulus immediately before Lemma~14 on p.~11.
Infinite primes may occur in the ray modulus; only its finite part
enters the norm bound.

For completeness, the elementary passage from Mitsui's proved
smoothed estimate (Lemma~14, pp.~11--12) to
\eqref{eq:mitsui-weighted-primes} also applies to complex characters.
Put
\[
 P(x,\vartheta)=\sum_{\mathrm N\mathfrak p\leq x}
  \vartheta(\mathfrak p)\log\mathrm N\mathfrak p
                         \log\frac{x}{\mathrm N\mathfrak p}.
\]
Lemma~14 gives $P(x,\vartheta)=E(\vartheta)x+
O_{K,A}(xe^{-c_1\sqrt{\log x}})$.
Subtract its values at $xe^h$ and $x$, where
$h=e^{-c_1\sqrt{\log x}/2}$, and divide by $h$.
The resulting left side is the prime sum in
\eqref{eq:mitsui-weighted-primes} plus boundary terms of absolute
value at most
\[
 \sum_{x<\mathrm N\mathfrak p\leq xe^h}\log\mathrm N\mathfrak p
       \ll_K(xh+1)\log x.
\]
Indeed, for each integer norm there are at most $[K:\Q]$ prime ideals.
The main term is $E(\vartheta)x(e^h-1)/h=E(\vartheta)x+O(xh)$,
and the error divided by $h$ is
$O_{K,A}(xe^{-c_1\sqrt{\log x}/2})$.
Reducing the positive exponent constant proves
\eqref{eq:mitsui-weighted-primes}.

Prime powers of exponent at least two contribute
$O_K(\sqrt x\log^2(2x))$: sum over exponents $j\geq2$ and use the
elementary bound $O_K(x^{1/j})$ for the number of prime ideals of norm
at most $x^{1/j}$.
Consequently \eqref{eq:mitsui-weighted-primes} holds with the prime
sum replaced by
$\sum_{\mathrm N\mathfrak a\leq x}
\Lambda_K(\mathfrak a)\vartheta(\mathfrak a)$.

The character $\rho\circ\mathrm N$ on $E$ has ray modulus $(q)$:
if $\alpha\equiv1\pmod{q\Z[i]}$, then
$\mathrm N\alpha\equiv1\pmod q$.
Thus $\chi(\rho\circ\mathrm N)$ has a modulus dividing
$\mathfrak m(q)$, whose norm is $\mathrm N\mathfrak m\,q^2$.
For $\sqrt X\leq Y\leq X$, the bounds
$q\leq(\log X)^B$ and $\log Y\geq\frac12\log X$ imply
\[
 \mathrm N\mathfrak m\,q^2\leq(\log Y)^{2B+1}
\]
for all sufficiently large $X$, with the threshold depending only
on $\mathfrak m$ and $B$.
Apply the preceding estimate with $K=E$, $x=Y$, and $A=2B+1$.
Its error is at most the error in
\eqref{eq:gaussian-hecke-sw-estimate}, after reducing $c$.
For $1\leq Y<\sqrt X$, the elementary bound
$\sum_{\mathrm N\mathfrak a\leq Y}\Lambda_E(\mathfrak a)
\ll_E Y\log^2(2Y)$ suffices, and the remaining bounded values of
$X$ are absorbed in the implied constant.

The principal character at the chosen modulus is exactly the
character induced by the trivial primitive character, giving the
stated $r_{\chi,\rho}$.
Equivalently one can work at the primitive conductor and remove
the extra Euler factors: the corresponding prime-power towers
contribute only
$O_\chi(\log(q+2)\log(Y+2))$, which is absorbed in the error.

For \eqref{eq:rational-sw-estimate}, take $K=\Q$ and use the ray
character attached to $\rho$, allowing the real infinite place
for odd characters.  Its finite modulus norm is $q$.
The same proof uses $A=B+1$ for $Y\geq\sqrt X$ and the elementary
bound for smaller $Y$.  Finally,
$e^{-c\sqrt{\log X}}\ll_{c,D}(\log X)^{-D}$ for every fixed $D>0$.
\end{proof}


The preceding lemma gives the following consequence, including $Y=0$ by the empty-sum convention.
  For every $B,D>0$, every fixed finite-order
ray character $\chi$ of $E$, and every Dirichlet character $\rho$ of
modulus $q\leq(\log X)^B$,
\begin{equation}\label{eq:gaussian-hecke-sw-input}
 \sum_{\mathrm N\mathfrak a\leq Y}
       \Lambda_E(\mathfrak a)\chi(\mathfrak a)
                         \rho(\mathrm N\mathfrak a)
  = r_{\chi,\rho}Y+O_{H,B,D}(X(\log X)^{-D})
 \quad(0\leq Y\leq X),
\end{equation}
where $r_{\chi,\rho}=1$ if the induced Hecke character
$\chi(\rho\circ\mathrm N)$ is trivial and is zero otherwise.
The omitted Euler factors only change the sum by a polylogarithmic
quantity.  The constant may be ineffective.  We require uniformity in
$\rho,q,Y$, but not in $H$.

\begin{lemma}[Progression density from the fixed-field input]
\label{lem:gaussian-ray-progression}
For $(a,q)=1$ define
$J_q=K\cap\Q(\zeta_q)$ and let $\sigma_a(\zeta_q)=\zeta_q^a$.
For fixed $B,D>0$, uniformly for $q\leq(\log X)^B$ and $0\leq Y\leq X$,
\[
 \sum_{\substack{n\leq Y\\ n\equiv a\pmod q}}\Lambda'_C(n)
 =\frac{\#\{g\in C:g|_{J_q}=\sigma_a|_{J_q}\}}
             {[K\Q(\zeta_q):\Q]}Y
       +O_{H,B,D}(X(\log X)^{-D}).
\]
\end{lemma}
\begin{proof}
Orthogonality of Dirichlet characters and
\eqref{eq:gaussian-ray-fourier} show that the coefficient of $Y$ for
$T_C$ is
\[
 \frac1{\varphi(q)}\sum_{\chi\in\widehat A}c_C(\chi)
       \sum_{\rho\bmod q}\overline{\rho(a)}r_{\chi,\rho}.
\]
The finite sum over $\chi$ has size depending only on $H$, while the
average over $\rho$ cancels its own factor $\varphi(q)$ in the error.

To evaluate this coefficient without a convention about inverse Artin
symbols, put $L=K\Q(\zeta_q)$ and $B_q=\operatorname{Gal}(L/E)$.
Restriction identifies $B_q$ with
\[
 \{(g,s)\in A\times\operatorname{Gal}(\Q(\zeta_q)/\Q):
                       g|_{J_q}=s|_{J_q}\}.
\]
The characters $\chi(g)\rho(s)$ restrict to characters of $B_q$.
Their average on $B_q$ is one precisely when their pullback by Artin
reciprocity is the trivial Hecke character; otherwise it is zero.
This is exactly $r_{\chi,\rho}$.  Expanding $1_C(g)$ and
$1_{s=\sigma_a}$ by Fourier inversion therefore identifies the displayed
coefficient with the proportion of elements of $B_q$ satisfying
$g\in C$ and $s=\sigma_a$.  Its numerator is the numerator in the
lemma and its denominator is $|B_q|=[L:E]$.
Dividing by two using \eqref{eq:gaussian-ray-prime-factor} changes the
denominator to $[L:\Q]$.  The prime-power error is smaller than the stated
error.  This proves the formula.
\end{proof}

Choose a fixed integer $W$ divisible by the conductor of $K^{\mathrm{ab}}$.
It is in particular divisible by $4$.
For a unit $b\pmod W$ put
\[
\begin{aligned}
 \mu_C(b)&=
 \frac{\#\{g\in C:g|_{K^{\mathrm{ab}}}=\sigma_b|_{K^{\mathrm{ab}}}\}}
      {[K:K^{\mathrm{ab}}]},\\
 f_{C,b,W}(n)&=\frac{\varphi(W)}W\Lambda'_C(Wn+b),\\
 \lambda_{b,W}(n)&=\frac{\varphi(W)}W\Lambda'(Wn+b).
\end{aligned}
\]
The prime-supported ordinary von Mangoldt function is denoted by
$\Lambda'$.  Replacing it by the usual $\Lambda$ costs only the same
negligible prime-power error.
The centered weight is
\begin{equation}\label{eq:gaussian-centered-weight}
       h_{C,b,W}=f_{C,b,W}-\mu_C(b)\lambda_{b,W}.
\end{equation}
Here $0\leq\mu_C(b)\leq1$, so
$|h_{C,b,W}(n)|\leq\lambda_{b,W}(n)=O_W(\log(2N))$ on $[N]$.

\begin{corollary}[Centered sums on progressions]
\label{cor:gaussian-centered-progressions}
For every $D>0$,
\[
       \sum_{n\in P}h_{C,b,W}(n)=O_{H,W,D}(N(\log N)^{-D})
\]
uniformly over arithmetic progressions $P$ contained in $[N]$.
\end{corollary}
\begin{proof}
Write $P$ as a residue class modulo $q$ intersected with an interval.
If $q>(\log N)^{D+3}$, its cardinality is at most
$N/(\log N)^{D+3}+1$, and the pointwise logarithmic bound proves the
claim directly.  Otherwise the condition $n\in P$ becomes one residue
class $a\pmod{Wq}$ and an interval for $Wn+b$.

If $(a,Wq)>1$, every prime in this residue class divides $Wq$.  There
are at most $O(\log(Wq))$ such primes, with total logarithmic weight
polylogarithmic in $N$, which is negligible here.  Suppose instead that
$(a,Wq)=1$.
Since $\Q(\zeta_{Wq})$ is abelian and already contains
$K^{\mathrm{ab}}$, one has
$K\cap\Q(\zeta_{Wq})=K^{\mathrm{ab}}$.
Also $a\equiv b\pmod W$, so the lemma's main coefficient is
\[
   \frac{\mu_C(b)}{\varphi(Wq)}.
\]
For the ordinary prime weight, the Siegel--Walfisz main coefficient
is $1/\varphi(Wq)$.  The two terms of
\eqref{eq:gaussian-centered-weight} therefore have equal main terms.
Apply the progression lemma at both interval endpoints, with an error
exponent larger than $D$ and with $X=O_W(N)$, and subtract.  This proves
the claimed bound.  All moduli used are at most a fixed power of
$\log N$, and $W$ is fixed throughout this argument.
\end{proof}


% Standalone section fragment; needs amsmath, amssymb, amsthm.
% Bibliography keys: GTOrbits2012 and GTMobius2012.
\subsection{Transfer from Gaussian ideal norms to polynomial nilsequences}
\label{sec:gn-transfer}

Throughout this subsection $E=\mathbb Q(i)$, $\mathcal O_E=\mathbb Z[i]$,
and $H_0$ is a fixed positive integer. A ray character $\chi$ modulo
$H_0\mathcal O_E$ is extended by zero to ideals not coprime to
$H_0\mathcal O_E$. The trivial character on all ideals is included by
taking $H_0=1$. We use the conventions of \cite{GTOrbits2012}: $G$ is
connected, simply connected and nilpotent, $\Gamma$ is a discrete cocompact
subgroup, and the filtration consists of closed connected rational
subgroups. Lipschitz norms include the supremum norm.

\begin{lemma}[Gaussian norm transfer]\label{lem:gn-norm-transfer}
Fix positive integers $m,d$. There are $\kappa,\sigma>0$, depending only
on $m,d$, and constants $C,X_0,\delta_0>0$, depending on $H_0,m,d$, with
the following property. Suppose
\[
 X\geq X_0,\qquad X^{-\sigma}\leq\delta\leq\delta_0,
 \qquad 1\leq Q\leq\delta^{-\kappa},\qquad
 1\leq q\leq\delta^{-\kappa},\quad q\in\mathbb Z.
\]
Let $G/\Gamma$ have dimension $m$, a filtration $G_\bullet$ of degree at
most $d$, and a $Q$-rational adapted Mal'cev basis. Suppose
$g\in\operatorname{poly}(\mathbb Z,G_\bullet)$ is totally
$\delta$-equidistributed on $[X]$. If $F:G/\Gamma\to\mathbb C$ is
Lipschitz and $\int F=0$, then for every $0\leq A<B\leq X$ and every
$b\in\mathbb Z$,
\begin{equation}\label{eq:gn-transfer}
 \left|\sum_{\substack{A<\mathrm N\mathfrak a\leq B\\
                 \mathrm N\mathfrak a\equiv b\pmod q}}
       \chi(\mathfrak a)F(g(\mathrm N\mathfrak a)\Gamma)\right|
 \leq C X\delta^\kappa\|F\|_{\mathrm{Lip}}.
\end{equation}
There is no lower bound on $B-A$ or on the number of progression terms.
Inserting $\log\mathrm N\mathfrak a$ on the left costs at most an
additional factor $O(\log(2X))$.
\end{lemma}

We prove the lemma using only the one-parameter quantitative Leibman
theorem \cite[Theorem 2.9]{GTOrbits2012}, together with elementary
polynomial and lattice arguments. No multiparameter equidistribution
theorem is used.

\begin{lemma}[Quadratic coefficient extrapolation]\label{lem:gn-quadratic-coefficients}
Let $X\geq1$, $P(t)=a_0+a_1t+a_2t^2\in\mathbb Z[t]$, $a_2\ne0$, and suppose
\[
 |a_i|\leq B_0X^{1-i/2},\qquad
 L\geq X^{1/2}/B_0,\qquad B_0\geq2.
\]
For every scalar polynomial $p:\mathbb Z\to\mathbb R/\mathbb Z$ of
degree at most $d$, there is an integer $1\leq k\leq C_dB_0^{C_d}$ such
that
\[
 \|kp\|_{C^\infty[X]}
   \leq C_dB_0^{C_d}\|p\circ P\|_{C^\infty[L]}.
\]
Here, for $p(n)=\sum_j\alpha_j\binom nj$, we use
$\|p\|_{C^\infty[X]}=\max_{1\leq j\leq d}X^j
\|\alpha_j\|_{\mathbb R/\mathbb Z}$.
\end{lemma}
\begin{proof}
Write $\binom{P(t)}j=\sum_\ell s_{j,\ell}\binom t\ell$. These
coefficients are integers since the left side is integer-valued. Its
degree is $2j$, and comparison of leading coefficients gives
\[
 s_{j,2j}=\frac{(2j)!}{j!}a_2^j\ne0,
 \qquad |s_{j,\ell}|\ll_d B_0^jX^{j-\ell/2}.
\]
For the bound, expand powers of $P$ and change between the monomial and
binomial bases. The change-of-basis matrices have fixed rational
coefficients. Descending induction in $\ell$ shows that every term
contributing to $s_{j,\ell}$ has weighted input degree at least $\ell$,
which gives the displayed estimate.

If $p=\sum_j\alpha_j\binom nj$ and
$p\circ P=\sum_\ell\beta_\ell\binom t\ell$, then
\[
 \beta_{2j}=\sum_{h=j}^d s_{h,2j}\alpha_h.
\]
Put $k_j=\prod_{h=j}^d s_{h,2h}$. Descending induction in $j$ gives
\[
 X^j\|k_j\alpha_j\|_{\mathbb R/\mathbb Z}
   \ll_d B_0^{O_d(1)}\|p\circ P\|_{C^\infty[L]}.
\]
Indeed the $\beta_{2j}$ term is bounded using
$X^jL^{-2j}\leq B_0^{2j}$, and a higher coefficient is multiplied by
$s_{h,2j}=O_d(B_0^hX^{h-j})$, canceling its $X^{-h}$ error.
All multipliers used to clear higher rows are integers of size
$B_0^{O_d(1)}$. Multiplying to the common nonzero integer $|k_1|$ proves
the assertion. The constant coefficient $\alpha_0$ is not used.
\end{proof}

\begin{lemma}[Small horizontal character obstruction]\label{lem:gn-horizontal-obstruction}
There are $c_0,\sigma_0>0$, depending only on $m,d$, such that the
following holds for sufficiently large $X$ and sufficiently small
$\delta\geq X^{-\sigma_0}$. Suppose the basis is $T$-rational and
a nontrivial horizontal character $\eta$ satisfies
\[
 |\eta|\leq T,\qquad \|\eta\circ g\|_{C^\infty[X]}\leq T,
 \qquad T\leq\delta^{-c_0}.
\]
Then $g$ is not totally $\delta$-equidistributed on $[X]$.
\end{lemma}
\begin{proof}
Choose representatives of the nonconstant binomial coefficients of
$\eta\circ g$ satisfying $|\alpha_j|\leq TX^{-j}$. On the initial
interval of length $\lfloor X/(100dT)\rfloor$ its values lie in an arc
of length less than $1/10$, translated by its constant coefficient.
Choose a Lipschitz function on $\mathbb R/\mathbb Z$ which equals one on
that arc, has mean zero and bounded Lipschitz norm. Its composition with
$\eta$ has Lipschitz norm at most $T^{C(m,d)}$, by the Mal'cev metric
estimates in \cite[Appendix A]{GTOrbits2012}. Its Haar mean is zero:
a nontrivial horizontal character maps Haar measure onto Haar measure
on the circle. Choose $c_0$ with $C(m,d)c_0<1/4$ and $c_0<1/4$.
The initial interval is longer than $\delta X$, while the discrepancy
one exceeds $\delta T^{C(m,d)}$. Taking $\sigma_0$ sufficiently small
absorbs the floors. This contradicts total equidistribution.
\end{proof}

\begin{proof}[Proof of Lemma~\ref{lem:gn-norm-transfer}]
Fix a constant $C_*$ which dominates the exponents in
Lemma~\ref{lem:gn-quadratic-coefficients}, the degree-$2d$ case of
\cite[Theorem 2.9]{GTOrbits2012}, and the preceding obstruction.
Choose
\[
 0<\kappa<\min\{(1000C_*^2)^{-1},c_0/(12C_*)\},\qquad
 0<\sigma<\min\{(1000C_*^2)^{-1},\sigma_0\},
 \qquad 2\kappa\sigma<1/4.
\]
These choices depend only on $m,d$. Decrease $\delta_0$, depending on
$H_0,m,d$, whenever a fixed multiplicative constant must be absorbed.

We first record a pullback consequence. Suppose $P$ and $L$ satisfy
Lemma~\ref{lem:gn-quadratic-coefficients} with
$B_0\leq\delta^{-5\kappa}$. Then
$(g(P(t))\Gamma)_{t\in[L]}$ is $\delta^\kappa$-equidistributed.
Otherwise the quantitative Leibman theorem gives a nontrivial $\eta$
with
\[
 |\eta|,\ \|\eta\circ g\circ P\|_{C^\infty[L]}
     \ll\delta^{-C_*\kappa}.
\]
The basis is $\delta^{-\kappa}$-rational as required. The composed
sequence has degree at most $2d$ for the stretched filtration
$\widetilde G_j=G_{\lceil j/2\rceil}$, with the same adapted basis.
This follows from the elementary one-parameter coordinate
characterization in \cite[Lemma 6.7]{GTOrbits2012}.
Extrapolation gives a nonzero integer multiple $k\eta$ with
\[
 |k\eta|,\ \|k\eta\circ g\|_{C^\infty[X]}
       \ll\delta^{-6C_*\kappa},
\]
with $C_*$ chosen sufficiently large initially. It remains nontrivial. The choice
$6C_*\kappa<c_0/2$ and a reduction of $\delta_0$ make this contradict
Lemma~\ref{lem:gn-horizontal-obstruction}. Only ordinary
equidistribution of this pullback is needed.

Every nonzero Gaussian ideal has exactly four generators. The function
$z\mapsto\chi((z))$, extended by zero when $z$ is not coprime to $H_0$,
is periodic modulo $H_0\mathbb Z[i]$ and invariant under units. Thus the
sum in \eqref{eq:gn-transfer} is exactly
\[
 \frac14\sum_{\substack{z=x+iy\ne0\\A<x^2+y^2\leq B\\
                         x^2+y^2\equiv b\pmod q}}
       \chi((z))F(g(x^2+y^2)\Gamma).
\]
The factor $1/4$ also applies to ideals with generators on an axis;
we have introduced no angular fundamental domain.

Set $\rho=\delta^\kappa$ and
$M=\lfloor\rho\sqrt X/H_0\rfloor$. On each horizontal row $y\in
\mathbb Z$, and in each residue class of $x$ modulo $H_0$, partition
the points into consecutive blocks of $M$ points with step $H_0$.
The character is constant on a block. The intersection of the annulus
with a horizontal row consists of at most two intervals, so at most
four blocks per residue class meet an endpoint. There are $O(\sqrt X)$
relevant rows. Consequently the boundary blocks contain
\[
 O(H_0\sqrt X\,M)=O(\rho X)
\]
points. Estimate their contribution trivially. Excluding $z=0$, if
necessary, costs one more bounded summand, absorbed by $\rho X$.

Every remaining contributing block lies wholly in the annulus. Split
its index modulo $q$. The norm congruence is constant on each resulting
class. Such a class has the form
\[
 x=x_0+H_0qn,\quad 1\leq n\leq L,\qquad
 P(n)=(x_0+H_0qn)^2+y^2,
 \qquad L\geq\frac{\rho\sqrt X}{4H_0q}.
\]
The hypotheses on $\delta$ and $\sigma$ ensure $M\geq2q$ and absorb
the floors for large $X$. A harmless endpoint adjustment gives
$|x_0|\leq2\sqrt X$. The coefficients of $P$ satisfy
\[
 |a_0|\ll X,\qquad |a_1|\ll H_0q\sqrt X,\qquad a_2=H_0^2q^2.
\]
Since $H_0$ is fixed and $q\leq\delta^{-\kappa}$, these bounds and the
lower bound for $L$ fit the pullback consequence with
$B_0\leq\delta^{-5\kappa}$ after decreasing $\delta_0$.
Hence the sum on the class is
$O(L\delta^\kappa\|F\|_{\mathrm{Lip}})$. The classes partition the
interior points, of which there are $O(X)$; this elementary bound follows
from the enclosing square. Summing and adding the boundary contribution
proves \eqref{eq:gn-transfer}.

The number of row endpoints is independent of $B-A$, explaining why no
lower interval length is required. Finally apply Abel summation to the
uniform bound with variable upper endpoint to insert $\log\mathrm N
\mathfrak a$ at cost $O(\log(2X))$.
\end{proof}

\begin{remark}\label{rem:gn-product-and-progression}
For Type II, apply the lemma in dimension $2m$ to
$(g_1(n),g_2(n))$ and $F_1\otimes\overline{F_2}$, when this entire
product orbit is totally equidistributed. Mean-zero factors ensure that
the product test function has mean zero; they do not imply
equidistribution of the product orbit. Use the smaller exponents for
dimensions $m$ and $2m$. In dimension zero mean-zero test functions
vanish, so the assertion is immediate.

The constraints $rn\equiv b\pmod q$ and $r'n\equiv b\pmod q$ either
are inconsistent or specify one residue class modulo a divisor of $q$.
Intersecting their interval restrictions only changes $A,B$. Thus all
Type I and Type II truncations are covered by the lemma. If
$\delta=(\log X)^{-B_1}$ and $Q,q\leq(\log X)^{C_1}$, any desired
saving $X(\log X)^{-A_1}$ follows once
$B_1\kappa>\max\{C_1,A_1\}$, with $H_0$ fixed.
\end{remark}


% Depends on the separately proved Lemma lem:gn-norm-transfer.
% No Frei--Tinkova theorem is assumed in this fragment.
% Their ideal-norm Type I/II strategy is retained as methodological attribution.
\subsection{Gaussian ideal Vaughan decomposition and Type I/II estimates}
\label{gi-section}

We give the ideal decomposition and the subsequent integer-parameter
recombination in full.  The architecture is that of
\cite[Section~3]{moebius}, with the Gaussian ideal-norm adaptation
attributed to Frei--Tinkova.  The analytic inputs used here are the
separately proved Lemma~\ref{lem:gn-norm-transfer} and the published
one-parameter results \cite[Theorem~2.9 and Lemma~4.5]{GTOrbits2012}.
In particular, no theorem of Frei--Tinkova is an assumed input.

Fix the Gaussian ray modulus $H_0$.  All ideals in this subsection are
nonzero integral ideals of $\mathbb Z[i]$.  Write $\mathrm N\mathfrak a$
for the absolute ideal norm, $\mu_E$ for ideal Mobius, and
$\Lambda_E(\mathfrak p^j)=\log\mathrm N\mathfrak p$ for $j\geq1$, with
$\Lambda_E$ zero on other ideals.  A ray character $\chi$ is extended by
zero on ideals not coprime to $H_0$.  It is completely multiplicative
and satisfies $|\chi|\leq1$.  We use a Lipschitz norm that includes the
supremum norm.  All nilmanifolds have connected, simply connected
nilpotent covering group, a cocompact lattice, and an adapted rational
Mal'cev basis, with the filtration and polynomial-sequence conventions
of \cite{GTOrbits2012}.

\begin{lemma}[Ideal multiplicities and moments]\label{gi-moments}
Let $a_E(n)=\#\{\mathfrak a:\mathrm N\mathfrak a=n\}$.  For $X\geq1$,
\begin{align}
 a_E(n)&\leq\tau(n),\label{gi-multiplicity}\\
 \sum_{n\leq X}a_E(n)^2&\leq X(1+\log X)^3,&
 \sum_{n\leq X}\frac{a_E(n)^2}{n}&\leq(1+\log X)^4.\label{gi-two-moments}
\end{align}
Moreover
\begin{equation}\label{gi-ideal-count}
 A_E(X):=\#\{\mathfrak a:\mathrm N\mathfrak a\leq X\}
       =\frac{\pi}{4}X+O(\sqrt X+1).
\end{equation}
In particular $A_E(X)\ll X$, and the number of ideals of norm in
$(B/2,B]$ is comparable to $B$ for all sufficiently large $B$.
\end{lemma}
\begin{proof}
At a split rational prime $p$ one has $a_E(p^k)=k+1$; at an inert prime
it is one for even $k$ and zero for odd $k$; and $a_E(2^k)=1$.
Multiplicativity proves \eqref{gi-multiplicity}.  Prime by prime,
$(k+1)^2\leq\binom{k+3}{3}$, so $\tau(n)^2\leq d_4(n)$.
Summing over ordered quadruples with product at most $X$, and summing
the last coordinate first, bounds $\sum_{n\leq X}d_4(n)$ by
$X(1+\log X)^3$.  Bounding each of the four harmonic sums separately
gives $\sum_{n\leq X}d_4(n)/n\leq(1+\log X)^4$.
Every ideal has exactly four Gaussian generators.  Counting lattice
points in the disk of radius $\sqrt X$, and dividing by four, proves
\eqref{gi-ideal-count}: the unit squares meeting its boundary lie
between circles of radii $\sqrt X-\sqrt2$ and $\sqrt X+\sqrt2$ and
number $O(\sqrt X+1)$.
\end{proof}

\begin{lemma}[Exact ideal Vaughan identity]\label{gi-vaughan}
Let $N\geq8$, $U=V=N^{1/3}$, $L=\log(2N)$, and let
$f:\mathbb N\longrightarrow\mathbb C$ satisfy $|f|\leq1$ and vanish
outside $(N/2,N]$.  Then
\begin{equation}\label{gi-vaughan-sum}
 \sum_{\mathfrak a}\Lambda_E(\mathfrak a)\chi(\mathfrak a)
                       f(\mathrm N\mathfrak a)=T_0-T_1+T_2,
\end{equation}
where
\begin{align*}
 T_0&=\sum_{\mathrm N\mathfrak d\leq U}\mu_E(\mathfrak d)\chi(\mathfrak d)
       \sum_{\mathfrak b}\log(\mathrm N\mathfrak b)\chi(\mathfrak b)
                  f(\mathrm N\mathfrak d\,\mathrm N\mathfrak b),\\
 T_1&=\sum_{\mathrm N\mathfrak d\leq UV}\alpha_{\mathfrak d}\chi(\mathfrak d)
       \sum_{\mathfrak b}\chi(\mathfrak b)
                  f(\mathrm N\mathfrak d\,\mathrm N\mathfrak b),\\
 T_2&=\sum_{V<\mathrm N\mathfrak d\leq N/U}\beta_{\mathfrak d}\chi(\mathfrak d)
       \sum_{\mathrm N\mathfrak b>U}\mu_E(\mathfrak b)\chi(\mathfrak b)
                  f(\mathrm N\mathfrak d\,\mathrm N\mathfrak b),
\end{align*}
with
\begin{align*}
 \alpha_{\mathfrak d}
   &=\sum_{\substack{\mathfrak b\mathfrak c=\mathfrak d\\
              \mathrm N\mathfrak b\leq U,\ \mathrm N\mathfrak c\leq V}}
                        \mu_E(\mathfrak b)\Lambda_E(\mathfrak c),&
 \beta_{\mathfrak d}
   &=\sum_{\substack{\mathfrak c\mid\mathfrak d\\\mathrm N\mathfrak c>V}}
                        \Lambda_E(\mathfrak c).
\end{align*}
Both coefficients have absolute value at most $L$.
\end{lemma}
\begin{proof}
Convolution on ideals satisfies
$\mu_E*1=\mathbf1_{\mathfrak a=(1)}$, $1*\Lambda_E=\log\mathrm N$,
and $\mu_E*\log\mathrm N=\Lambda_E$, as one checks on prime powers.
Splitting $\mu_E$ at $U$ and $\Lambda_E$ at $V$ gives the exact identity
\begin{equation}\label{gi-convolution-identity}
 \Lambda_E=\mu_{\leq U}*\log\mathrm N+\Lambda_{\leq V}
       -\mu_{\leq U}*1*\Lambda_{\leq V}
       +\mu_{>U}*1*\Lambda_{>V}.
\end{equation}
The second term vanishes on the support of $f$.  Complete multiplicativity
of $\chi$, including its zero extension, permits separation of each
convolution factor and yields \eqref{gi-vaughan-sum} with exactly the
stated cutoffs.  Finally
$\sum_{\mathfrak c\mid\mathfrak d}\Lambda_E(\mathfrak c)
=\log\mathrm N\mathfrak d$ and every summand is nonnegative.  This
proves the coefficient bounds, with no discarded endpoint or error term.
\end{proof}

\begin{proposition}[Quantitative extraction]\label{gi-extraction}
There are absolute constants $c_0>0$ and $N_0$ such that the following
holds.  Under Lemma~\ref{gi-vaughan}, suppose $N\geq N_0$,
$0<\epsilon\leq1$, and
\begin{equation}\label{gi-large-correlation}
 \left|\sum_{\mathfrak a}\Lambda_E(\mathfrak a)\chi(\mathfrak a)
                         f(\mathrm N\mathfrak a)\right|\geq\epsilon NL.
\end{equation}
Put $\theta=c_0\epsilon^4/L^{24}$.  At least one of the following holds.
\begin{enumerate}
\item Some dyadic $R\leq2N^{1/3}$ has at least $\theta R$ integers
$r\in(R/2,R]$ satisfying
\[
 \left|\sum_{\mathfrak b}\log(\mathrm N\mathfrak b)\chi(\mathfrak b)
                   f(r\mathrm N\mathfrak b)\right|\geq\theta N/r.
\]
\item Some dyadic $R\leq2N^{2/3}$ has at least $\theta R$ integers
$r\in(R/2,R]$ satisfying the same inequality without the logarithm.
\item There are dyadic $R,B$ with
\begin{equation}\label{gi-dyadic-ranges}
 N^{1/3}<R\leq2N^{2/3},\qquad N/2<RB<4N,
\end{equation}
and at least $\theta R^2$ pairs $r,r'\in(R/2,R]$ satisfying
\begin{equation}\label{gi-type-two-correlation}
 \left|\sum_{B/2<\mathrm N\mathfrak b\leq B}
       f(r\mathrm N\mathfrak b)\overline{f(r'\mathrm N\mathfrak b)}
       \right|\geq\theta B.
\end{equation}
\end{enumerate}
\end{proposition}
\begin{proof}
One of $|T_i|$ is at least $\epsilon NL/3$.  For $T_0$ put
\[
 G(r)=\frac{r}{NL}\left|\sum_{\mathfrak b}\log(\mathrm N\mathfrak b)
                   \chi(\mathfrak b)f(r\mathrm N\mathfrak b)\right|;
\]
for $T_1$ use $r/N$ times the absolute unweighted inner sum.  In either
case Lemma~\ref{gi-moments} gives $G\ll1$ and
\[
 \epsilon\ll\sum_{r\leq N^e}\frac{a_E(r)G(r)}r,
 \qquad e=1/3\text{ or }2/3.
\]
Cauchy--Schwarz and \eqref{gi-two-moments} imply
$\sum_{r\leq N^e}G(r)^2/r\gg\epsilon^2/L^4$.
There are $O(L)$ dyadic cells, so one satisfies
\[
 \sum_{R/2<r\leq R}G(r)^2\gg R\epsilon^2/L^5.
\]
Since $G\ll1$, at least $cR\epsilon^2/L^5$ terms have
$G(r)\geq c\epsilon/L^{5/2}$.  Both the count and the claimed amplitudes
follow after decreasing $c_0$.

For $T_2$, split both ideal norms into dyadic cells.  There are $O(L^2)$
rectangles, and a nonzero cell satisfies \eqref{gi-dyadic-ranges} and
$B\geq N^{1/3}/4$.  A rectangle contributes at least $c\epsilon N/L$;
we use the weaker bound $c\epsilon N/L^2$.  Grouping $\mathfrak d$ by
$r=\mathrm N\mathfrak d$ writes this contribution as
\[
 T=\sum_{R/2<r\leq R}A_r\sum_{\mathfrak b\in\mathcal B}c_{\mathfrak b}
                   f(r\mathrm N\mathfrak b),\qquad
 \mathcal B=\{\mathfrak b:B/2<\mathrm N\mathfrak b\leq B\}.
\]
Here $|A_r|\leq L a_E(r)$ and $|c_{\mathfrak b}|\leq1$; the restriction
$\mathrm N\mathfrak b>U$ belongs to $c_{\mathfrak b}$.  All product
cutoffs remain in $f$, so the domain is rectangular.  By
\eqref{gi-two-moments}, $\sum_r|A_r|^2\ll RL^5$.  Thus, with
$F_{r,\mathfrak b}=f(r\mathrm N\mathfrak b)$,
\begin{equation}\label{gi-first-cauchy}
 T_*:=\sum_r\left|\sum_{\mathfrak b\in\mathcal B}
                          c_{\mathfrak b}F_{r,\mathfrak b}\right|^2
          \gg \epsilon^2 RB^2/L^9.
\end{equation}
Put $M=|\mathcal B|\asymp B$.  Expansion and Cauchy--Schwarz in the
ordered pair $(\mathfrak b,\mathfrak b')$ give
\begin{align}
 T_*^2
 &=\left|\sum_{\mathfrak b,\mathfrak b'}
 c_{\mathfrak b}\overline{c_{\mathfrak b'}}
 \sum_r F_{r,\mathfrak b}\overline{F_{r,\mathfrak b'}}\right|^2\notag\\
 &\leq M^2\sum_{\mathfrak b,\mathfrak b'}
       \left|\sum_rF_{r,\mathfrak b}\overline{F_{r,\mathfrak b'}}\right|^2
 =M^2\sum_{r,r'}\left|\sum_{\mathfrak b}
                  F_{r,\mathfrak b}\overline{F_{r',\mathfrak b}}\right|^2.
 \label{gi-second-cauchy}
\end{align}
Consequently
\[
 \mathbb E_{r,r'}\left|\mathbb E_{\mathfrak b\in\mathcal B}
              F_{r,\mathfrak b}\overline{F_{r',\mathfrak b}}\right|^2
       \gg\epsilon^4/L^{18}.
\]
Every summand is at most one.  Thresholding gives
\eqref{gi-type-two-correlation} and its stated count, with the weaker
common parameter $\theta=c_0\epsilon^4/L^{24}$.  The character and Mobius
coefficients have disappeared.  All conjugates in
\eqref{gi-second-cauchy} are necessary.
\end{proof}

\paragraph{Progressions and interval restrictions.}\label{gi-progressions}
In the application,
\[
 f(n)=1_{N/2<n\leq N}\,1_{u<n\leq v}\,1_{n\equiv b\pmod q}
                              F(g(n)\Gamma).
\]
For Type I the congruence $rn\equiv b\pmod q$ is either insoluble or
one class modulo $q/\gcd(r,q)$.  In Type II its intersection with
$r'n\equiv b\pmod q$ is either empty or one class modulo
\[
 \operatorname{lcm}\bigl(q/\gcd(r,q),q/\gcd(r',q)\bigr),
\]
which divides $q$.  The remaining restrictions are interval intersections.
Lemma~\ref{lem:gn-norm-transfer} covers them without a lower bound on
their lengths.  On the product nilmanifold the function is
$F(x)\overline{F(y)}$, whose Haar integral is
$|\int F|^2=0$.  After normalizing $\|F\|_{\mathrm{Lip}}=1$, its
Lipschitz norm for the sum metric is at most two.  Thus all mean-zero,
modulus and interval hypotheses for product norm transfer are satisfied.

\begin{lemma}[Integer-parameter recombination]\label{gi-recombination}
For each $m,d$ there are $C_{\rm rec},K_{\rm rec}\geq2$ with the following
property.  Let $0<\eta<\eta_0$, $N\geq\eta^{-K_{\rm rec}}$, and let $g$
have degree at most $d$ on an $m$-dimensional nilmanifold with an
$\eta^{-1}$-rational adapted basis.  Each of the following implies that
$(g(n)\Gamma)_{n\leq N}$ is not totally $\eta^{C_{\rm rec}}$-equidistributed.
\begin{enumerate}
\item For at least $\eta R$ values $r\in(R/2,R]$, where
$R\leq2N^{2/3}$, the orbit $(g(rn)\Gamma)_{n\leq\lceil N/r\rceil}$
is not totally $\eta$-equidistributed.
\item For at least $\eta R^2$ pairs $r,r'\in(R/2,R]$, where
$N^{1/3}/4\leq R\leq2N^{2/3}$ and $N/2<RB<4N$, the product orbit
$(g(rn),g(r'n))_{n\leq B}$ is not totally $\eta$-equidistributed.
\end{enumerate}
\end{lemma}
\begin{proof}
All constants in this proof depend only on $m,d$.  Failure of total
equidistribution gives a progression $a+q[\ell]\subseteq[M]$ with
$\ell\geq\eta M$, $q\leq2/\eta$, and $|a|\leq2M$ on which ordinary
$\eta$-equidistribution fails.  Apply
\cite[Theorem~2.9]{GTOrbits2012}, in dimension $m$ or $2m$.  A nonzero
horizontal character has magnitude and restricted smoothness norm at
most $\eta^{-C_1}$, after shrinking $\eta_0$ to absorb fixed factors.

Here is the required removal of a possibly large shift.  Write the scalar
phase $p(a+qt)$ in monomials.  Multiplication by $d!$ clears conversion
from its binomial coefficients.  Inverting $t=(n-a)/q$ and multiplying
by $q^d$ clears every remaining denominator.  The error in the coefficient
of $n^i$ is bounded by
\begin{equation}\label{gi-shift-removal}
 C_d\eta^{-C_1}\sum_{j=i}^d q^{d-j}|a|^{j-i}\ell^{-j}
       \leq\eta^{-C_2}M^{-i},\qquad i\geq1.
\end{equation}
The integer parts disappear modulo one.  Thus a nonzero integer multiple
of the character has magnitude and $C^\infty[M]$ norm at most
$\eta^{-C_2}$.  This argument does not treat $a=O(M)$ as a bounded-height
rational parameter.  There are at most $\eta^{-C_3}$ possible resulting
characters, since their lattice has rank at most $2m$.  Pigeonholing
therefore fixes one character on a set of density at least $\eta^{C_4}$.

In Type I write
$\psi(g(n))=\sum_{j=0}^d\beta_jn^j$ modulo one.  For the selected set
$S\subseteq(R/2,R]$, coefficient conversion gives
\begin{equation}\label{gi-one-coefficients}
 \|d!\beta_jr^j\|_{\mathbb R/\mathbb Z}
       \leq\eta^{-C_5}(R/N)^j,\qquad 1\leq j\leq d.
\end{equation}
If $R\leq\eta^{-C_6}$ for a sufficiently large fixed $C_6$, choose one
$r\in S$ and multiply the character by $d!r^d$.  Multiplication of the
$j$th bound by $r^{d-j}$ gives magnitude and $C^\infty[N]$ norm bounded
by a fixed power of $\eta^{-1}$.  This treats every small dyadic scale,
including $R=1$.

Otherwise apply \cite[Lemma~4.5]{GTOrbits2012} to
$r\mapsto d!\beta_jr^j$ on $[\lceil R\rceil]$.  The selected set has
density at least a fixed power of $\eta$, and its values lie in an arc
of length at most $2\eta^{-C_5}(R/N)^j$.  Taking $K_{\rm rec}$ large,
the bound $R\leq2N^{2/3}$ makes this length at most half the density.
The choice of $C_6$ also places the recurrence length beyond the fixed
inverse density powers.  The published recurrence lemma supplies
nonzero integers $t_j$, $|t_j|\leq\eta^{-C_7}$, such that
\begin{equation}\label{gi-recurrent-coefficients}
 \|t_j(d!)^2\beta_j\|_{\mathbb R/\mathbb Z}
                    \leq\eta^{-C_8}N^{-j}.
\end{equation}
The extra factorial clears the leading-coefficient conversion.  Taking
the product of the $t_j$ and the fixed factorials gives a single nonzero
horizontal character of magnitude and $C^\infty[N]$ norm at most
$\eta^{-C_9}$.

For Type II the pigeonholed character is $\psi_1\oplus\psi_2$, not both
zero.  Write the two scalar phases with coefficients $\beta_j,\beta'_j$.
The information replacing \eqref{gi-one-coefficients} is
\begin{equation}\label{gi-two-coefficients}
 \|d!(\beta_jr^j+\beta'_jr'^j)\|_{\mathbb R/\mathbb Z}
                        \leq\eta^{-C_5}(R/N)^j.
\end{equation}
Interchange the two coordinates if necessary so that $\psi_1\ne0$.
Fix an $r'$ with at least $\eta^{C_4}R$ partners.  Equation
\eqref{gi-two-coefficients} puts $d!\beta_jr^j$ in an arc centered at
$-d!\beta'_jr'^j$.  The same published recurrence lemma permits this
arbitrary center and gives \eqref{gi-recurrent-coefficients} for
$\psi_1$.  Its large-length requirement follows from
$R\geq N^{1/3}/4$ and a sufficiently large $K_{\rm rec}$.  There is no
substitution of $(r-r_0)^j$ for $r^j-r_0^j$.

Finally let $\Psi$ denote the nonzero character just obtained.  Its
binomial expansion shows that for
$1\leq n\leq c\eta^{2C_9+4}N$ the phase $\Psi(g(n))$ stays within
$1/10$ of $\Psi(g(0))$.  Choose a mean-zero Lipschitz circle function
equal to one on that arc and compose it with $\Psi$.  Its mean is zero
because a nonzero horizontal character pushes Haar measure to circle
Haar measure.  Its Lipschitz norm is bounded by a fixed power of
$\eta^{-1}$ by \cite[Appendix~A]{GTOrbits2012}.  Choose $C_{\rm rec}$
larger than the length and Lipschitz exponents, and then enlarge
$K_{\rm rec}$ so the initial interval is nonempty.  Its average one
contradicts total $\eta^{C_{\rm rec}}$-equidistribution.  This proves
both assertions.  All exponent constants are finite combinations of
the published one-parameter constants and degree-$d$ polynomial algebra.
\end{proof}

\begin{proposition}[Power-form Gaussian ideal minor-arc bound]
\label{gi-minor-arc}
For each $m,d$ there is $c_{\rm G}>0$ with the following property.
Fix $H_0$ and a finite-order Gaussian ray character $\chi$ of modulus
$H_0$, extended by zero.  Let $G/\Gamma$ have dimension $m$, filtration
degree $d$, and a $Q$-rational adapted basis, with $Q\geq2$.  Suppose
$(g(n)\Gamma)_{n\leq N}$ is totally $\delta$-equidistributed, where
$0<\delta<1$.  If $F$ is Lipschitz and $\int F=0$, then for every
interval progression $P\subseteq[N]$ of common difference at most $Q$,
\begin{equation}\label{gi-power-bound}
 \left|\sum_{\mathrm N\mathfrak a\leq N}\Lambda_E(\mathfrak a)
       \chi(\mathfrak a)1_P(\mathrm N\mathfrak a)
       F(g(\mathrm N\mathfrak a)\Gamma)\right|
 \ll_{H_0,m,d} QN\log(2N)^7\delta^{c_{\rm G}}\|F\|_{\mathrm{Lip}}.
\end{equation}
No lower bound on the length of $P$ or on $\delta$ is required.
\end{proposition}
\begin{proof}
Normalize $\|F\|_{\mathrm{Lip}}=1$.  Take the smaller positive exponents
$\kappa,\sigma$ in Lemma~\ref{lem:gn-norm-transfer} for dimensions $m$
and $2m$, and decrease $\kappa$ to at most one.  Partial summation gives
the same transfer bound with an extra factor $\log(2M)$ for a
$\log\mathrm N\mathfrak b$ weight: the unweighted lemma holds for
every partial interval, and $\int_1^Mdt/t\leq\log(2M)$.
Let $C_{\rm rec},K_{\rm rec}$ be as in Lemma~\ref{gi-recombination},
and choose
\begin{equation}\label{gi-exponent-choice}
 a_*=\frac1{4C_{\rm rec}},\qquad
 c_* =\frac{a_*\kappa}{32},\qquad
 s_* =\min\left\{\frac1{100},\frac1{4a_*K_{\rm rec}},
                             \frac{\sigma}{12a_*}\right\}.
\end{equation}
First assume $N^{-s_*}\leq\delta\leq\delta_0$, with the fixed cutoff
$\delta_0$ sufficiently small, and restrict the sum to $(N/2,N]$.
We claim that its absolute value is
\begin{equation}\label{gi-shell-bound}
                  \ll QNL^7\delta^{c_*},\qquad L=\log(2N).
\end{equation}
The elementary bound is $O(NL)$.  Thus the claim is immediate when
$QL^6\delta^{c_*}\geq1$, and in the other case
$Q\leq\delta^{-c_*}$ and $L\leq\delta^{-c_*/6}$.
If \eqref{gi-shell-bound} fails with a sufficiently large implied
constant, Proposition~\ref{gi-extraction} applies with
$\epsilon$ a fixed large multiple of $QL^6\delta^{c_*}$; if this
exceeds the permitted range, the elementary bound already proves the
claim.  The extracted density and normalized amplitude are at least
a fixed positive constant times $\delta^{4c_*}$, hence at least
$\delta^{5c_*}$ after decreasing $\delta_0$.

Set $\eta=\delta^{a_*}$.  If an extracted orbit were totally
$\eta$-equidistributed, Lemma~\ref{lem:gn-norm-transfer}, with the
progression analysis above, would bound the extracted sum by
\[
 O(M\delta^{a_*\kappa}L^e),\qquad e=0\text{ or }1.
\]
Here $M=\lceil N/r\rceil\geq N^{1/3}/2$ in Type I and
$M=B\geq N^{1/3}/4$ in Type II.  The rationality and modulus hypotheses
hold because $c_*<a_*\kappa$.  The lower-parameter hypothesis holds
because $a_*s_*\leq\sigma/12$; increasing the fixed threshold for $N$
absorbs the constants two and four.  Since $a_*\kappa=32c_*$, the
displayed bound is smaller than $M\delta^{5c_*}$, even for $e=1$.
This is a contradiction.  Thus every extracted orbit fails total
$\eta$-equidistribution.  Its density is at least $\eta$, since
$a_*>5c_*$.  Also $N\geq\eta^{-K_{\rm rec}}$ by
$a_*s_*K_{\rm rec}\leq1/4$, and $Q\leq\eta^{-1}$.
Lemma~\ref{gi-recombination} now contradicts total
$\delta$-equidistribution, because
$\eta^{C_{\rm rec}}=\delta^{1/4}>\delta$.  This proves
\eqref{gi-shell-bound}.  Fixed values $\delta\geq\delta_0$ are covered
by the elementary bound with a larger constant.

We spell out the passage through all scales.  Use exact real endpoints
$M_j=2^{-j}N$ and retain shells with $M_j\geq2\sqrt\delta N$.
The discarded norms are at most $2\sqrt\delta N$ and their total
absolute weight is $O(\sqrt\delta NL)$ by Lemma~\ref{gi-moments}.
There is no rounding of ideal-norm cutoffs.  On $[\lfloor M_j\rfloor]$
the orbit is totally $2\sqrt\delta$-equidistributed: every progression
of the requisite length has at least $\delta N$ terms for sufficiently
large $N$, and the original discrepancy bound is stronger.
For $\delta\geq N^{-s_*/4}$ the lower-parameter condition for the shell
bound holds on every retained scale, since
$M_j\geq N^{1-s_*/8}$ up to fixed constants.  The progression is merely
intersected with the shell; its common difference remains at most $Q$.
Summing the shell estimates and using $\sum_jM_j\leq2N$ gives
\begin{equation}\label{gi-full-intermediate}
                  O(QNL^7\delta^{c_*/2}).
\end{equation}
For $1/N<\delta<N^{-s_*/4}$ replace $\delta$ by
$\delta^{s_*/8}\geq N^{-s_*/8}$.  Total equidistribution at the
original parameter implies it at this larger one.  The exponent in
\eqref{gi-full-intermediate} becomes $c_*s_*/16$.
For $\delta\leq1/N$, singleton progressions imply
$|F(g(n)\Gamma)|\leq\delta$ at each integer $n$, so the elementary
$O(NL)$ weight bound suffices.  All remaining bounded $N$ are absorbed
by the fixed constant.  Thus one may take
\begin{equation}\label{gi-final-exponent}
                         c_{\rm G}=c_*s_*/16>0.
\end{equation}
Its dependence is only on $m,d$ and the corresponding norm-transfer
exponents; $H_0$ enters constants and thresholds, not this exponent.
\end{proof}

\begin{corollary}[Ordinary prime term]\label{gi-ordinary-primes}
Under the nilmanifold hypotheses of Proposition~\ref{gi-minor-arc},
there is $c_{\mathbb Q}>0$, depending only on $m,d$, such that
\begin{equation}\label{gi-ordinary-bound}
 \left|\sum_{n\leq N}\Lambda'(n)1_P(n)F(g(n)\Gamma)\right|
 \ll_{m,d} QN\log(2N)^7\delta^{c_{\mathbb Q}}\|F\|_{\mathrm{Lip}},
\end{equation}
where $\Lambda'(p)=\log p$ on rational primes and is zero otherwise.
The same bound holds for ordinary von Mangoldt $\Lambda$.
\end{corollary}
\begin{proof}
For positive integers the convolution identities
$\mu*1=1_{n=1}$, $1*\Lambda=\log$, and $\mu*\log=\Lambda$ give
\eqref{gi-convolution-identity} with exactly the displayed norm cutoffs
replaced by integer cutoffs.  Its coefficients are bounded by
$\sum_{c\mid d}\Lambda(c)=\log d$.  Integer multiplicity is one,
so both moment bounds in Lemma~\ref{gi-moments} remain valid.  The
two Cauchy--Schwarz calculations \eqref{gi-first-cauchy} and
\eqref{gi-second-cauchy} apply to the resulting finite integer sums
and give Proposition~\ref{gi-extraction} with the same conservative
parameter $c_0\epsilon^4/L^{24}$.  Lemma~\ref{gi-recombination} concerns
only integer dilation parameters and is unchanged.

The unweighted replacement for norm transfer follows directly from
total equidistribution.  If an interval progression $J\subseteq[X]$
has $|J|\geq\zeta X$, then
\[
 \left|\sum_{n\in J}F(g(n)\Gamma)\right|
       \leq |J|\zeta\|F\|_{\mathrm{Lip}}
       \leq X\zeta\|F\|_{\mathrm{Lip}}.
\]
If $|J|<\zeta X$, its cardinality and the supremum norm give the same
bound.  This treats every modulus and every truncated interval.
Partial summation supplies the logarithmic weight; the identical
long-progression/short-progression dichotomy treats the mean-zero
product function $F\otimes\overline F$ on the product nilmanifold.
Thus the exponent choices and all shell steps in the proof of
Proposition~\ref{gi-minor-arc} apply, with any sufficiently small
fixed positive transfer exponents, to give \eqref{gi-ordinary-bound}
for $\Lambda$.  No extension of a Mobius theorem to $\Lambda$ is assumed.

Finally proper prime powers contribute at most
$O(\sqrt N\log(2N)^2\|F\|_{\mathrm{Lip}})$, by counting at most
$N^{1/k}$ possible bases for each $2\leq k\leq\log_2N$.
Shrink $c_{\mathbb Q}$ below $1/4$.  If $\delta>1/N$, this error is
absorbed because $N\delta^{c_{\mathbb Q}}\geq N^{1-c_{\mathbb Q}}$.
If $\delta\leq1/N$, singleton equidistribution and
$\sum_{n\leq N}\Lambda'(n)\leq N\log(2N)$ prove the bound directly.
This gives the assertion for $\Lambda'$ as well.
\end{proof}



% Section fragment. Requires the gn-* lemmas from norm_transfer.tex.
\subsection{Fixed-modulus affine interpolation}
\label{sec:ga-affine}

Fix integers $W\geq1$ and $0\leq b<W$. The modulus $W$ remains fixed
as the size parameter tends to infinity. We first justify the affine
interpolation needed to apply ideal-norm estimates to a $W$-tricked
sequence.

\begin{lemma}[Linear coefficient extrapolation]\label{lem:ga-linear-coefficients}
Let $X\geq1$ and let $p:\mathbb Z\to\mathbb R/\mathbb Z$ have degree at most $d$.
Suppose $P(t)=u+vt$ with $u,v\in\mathbb Z$, $v\ne0$,
$|u|\leq C X$, $|v|\leq E$, and $L\geq X/E$, where $E\geq2$.
There is an integer $1\leq k\leq C_dE^{C_d}$ such that
\[
 \|kp\|_{C^\infty[X]}
 \leq C_d(C+E)^{C_d}\|p\circ P\|_{C^\infty[L]}.
\]
\end{lemma}
\begin{proof}
Expand $\binom{u+vt}j$ in the binomial basis in $t$. All coefficients
are integers, the diagonal coefficient is $v^j$, and the coefficient in
row $\ell$ is $O_d((C+E)^jX^{j-\ell})$. Starting in degree $d$, clear
the nonzero diagonal entries successively. Their product is a nonzero
integer bounded by $E^{O_d(1)}$. At degree $j$ use
$X^j/L^j\leq E^j$; a higher-row error is multiplied by
$O_d((C+E)^hX^{h-j})$, canceling its $X^{-h}$ factor. Multiplication
to the common product of diagonal entries proves the claim, exactly
as in Lemma~\ref{lem:gn-quadratic-coefficients}.
\end{proof}

\begin{lemma}[Affine interpolation preserves quantitative distribution]
\label{lem:ga-affine-interpolation}
Fix dimension and degree bounds $m,d$. There are $a,\sigma>0$, depending
only on $m,d$, and thresholds depending on $W,m,d$, with the following
property. Suppose a nilmanifold and its rational filtration satisfy the
Green--Tao conventions, its adapted basis is $Q$-rational, and
$g\in\operatorname{poly}(\mathbb Z,G_\bullet)$. Suppose also that
\[
 N^{-\sigma}\leq\delta\leq\delta_0(W,m,d),\qquad
 Q\leq\delta^{-a},
\]
and that $g$ is totally $\delta$-equidistributed on $[N]$. There is
$\widetilde g\in\operatorname{poly}(\mathbb Z,G_\bullet)$ satisfying
\[
 \widetilde g(Wn+b)=g(n)\quad(n\in\mathbb Z)
\]
which is totally $\delta^a$-equidistributed on $[WN+b]$.
\end{lemma}
\begin{proof}
Use the one-parameter coordinate characterization in
\cite[Lemma 6.7]{GTOrbits2012}. Each second-kind Mal'cev coordinate of
$g(n)$ is a polynomial, with the prescribed degree bound from the
filtration. Evaluate those same polynomials at real arguments and apply
the inverse coordinate map, obtaining $g_{\mathbb R}(t)$. Set
\[
 \widetilde g(t)=g_{\mathbb R}((t-b)/W).
\]
Affine composition preserves the coordinate degree bounds, so the same
characterization shows that this is a polynomial sequence for the
original filtration. It gives the stated identity on integer $n$.
Equivalently, one may use the Taylor product with real powers; those
powers exist in a connected simply connected nilpotent group.

Put $X=WN+b$ and $\epsilon=\delta^a$. Suppose the asserted total
$\epsilon$-equidistribution fails. Then on a progression
\[
 u+v[L]\subseteq[X],\qquad L\geq\epsilon X,
 \quad u,v\in\mathbb Z,\quad
 1\leq v\leq2/\epsilon,\quad |u|\leq2X,
\]
the sequence is not $\epsilon$-equidistributed. By the one-parameter
quantitative Leibman theorem \cite[Theorem 2.9]{GTOrbits2012}, there is
a nontrivial horizontal character $\eta$ satisfying
\[
 |\eta|\ll\epsilon^{-C},\qquad
 \|\eta\circ\widetilde g\circ(u+vt)\|_{C^\infty[L]}
       \ll\epsilon^{-C}.
\]
The rationality hypothesis holds since $Q\leq\epsilon^{-1}$.
Lemma~\ref{lem:ga-linear-coefficients} gives a positive integer
$k\ll\epsilon^{-C}$ such that
\[
 \|k\eta\circ\widetilde g\|_{C^\infty[X]}
      \ll\epsilon^{-C}.
\]
Forward composition with the fixed integer polynomial $Wn+b$ gives
\[
 \|k\eta\circ g\|_{C^\infty[N]}
      \ll_{W,d}\epsilon^{-C}.
\]
For this last bound expand $\binom{Wn+b}j$ in the binomial basis; its
matrix is integral, and $X$ is comparable to $WN$. The character
$k\eta$ remains nontrivial.

Let $C_{\mathrm{aff}}$ bound all exponents in these estimates, including
the character size. Fix explicitly
\[
 0<a<\min\{1,c_0/(2C_{\mathrm{aff}})\},\qquad
 0<\sigma<\min\{\sigma_0,1/(2a)\},
\]
where $c_0,\sigma_0$ are from
Lemma~\ref{lem:gn-horizontal-obstruction}. Decrease $\delta_0$ to
absorb fixed $W$-dependent constants. The resulting character bounds
contradict that lemma and total $\delta$-equidistribution of $g$.
The inequality $a\sigma<1/2$ ensures $\epsilon X\gg N^{1/2}$, so all
progression parametrizations and floor estimates used above hold for
large $N$. This proves the assertion on the entire interval $[WN+b]$,
including progressions outside the residue class $b$ modulo $W$.
\end{proof}

The Gaussian input in the following lemma is furnished by
Proposition~\ref{gi-minor-arc}; its rational counterpart is
Corollary~\ref{gi-ordinary-primes}.

\begin{lemma}[$W$-tricked minor-arc interface]\label{lem:ga-wtrick-ma}
For $E=\mathbb Q(i)$ and a fixed ray character $\chi$, put
\[
 T_\chi(t)=\sum_{\mathrm N\mathfrak a=t}
             \Lambda_E(\mathfrak a)\chi(\mathfrak a).
\]
Assume the mean-zero ideal estimate
\begin{equation}\label{eq:ga-ideal-input}
 \left|\sum_{\substack{\mathrm N\mathfrak a\leq X\\
                          \mathrm N\mathfrak a\in R}}
  \Lambda_E(\mathfrak a)\chi(\mathfrak a)
             F_0(h(\mathrm N\mathfrak a)\Gamma)\right|
 \ll Q_1X\log(2X)^7\eta^{c_E}\|F_0\|_{\mathrm{Lip}}
\end{equation}
for a positive $c_E$, whenever $h$ is totally $\eta$-equidistributed,
the basis is $Q_1$-rational, and $R$ is any interval progression of
modulus at most $Q_1$. Then, under the hypotheses of
Lemma~\ref{lem:ga-affine-interpolation}, for every Lipschitz $F_0$
with Haar mean zero and every interval progression
$P\subseteq[N]$ of modulus $q\leq Q$,
\begin{equation}\label{eq:ga-wtrick-ideal}
 \left|\sum_{n\in P}T_\chi(Wn+b)F_0(g(n)\Gamma)\right|
 \ll_{H_0,W,m,d}QN\log(2N)^7\delta^{ac_E}
                     \|F_0\|_{\mathrm{Lip}}.
\end{equation}
The same implication holds for $E=\mathbb Q$ and the corresponding
ordinary von Mangoldt estimate.
\end{lemma}
\begin{proof}
Take $X=WN+b$ and the sequence $\widetilde g$ constructed above. The
conditions $t=Wn+b$ and $n\in P$ specify an interval progression $R$
of modulus $Wq$. Apply \eqref{eq:ga-ideal-input} with
$Q_1=\max\{Q,Wq,2\}\ll_WQ$ and $\eta=\delta^a$. The ideal sum
then equals the left side of \eqref{eq:ga-wtrick-ideal} exactly, while
$X\ll_WN$ and $\log(2X)\ll_W\log(2N)$. No assertion about
fractional-argument distribution has been assumed: it was proved in
Lemma~\ref{lem:ga-affine-interpolation}.
\end{proof}

\begin{corollary}[Centered ray-prime weights]\label{cor:ga-centered-primes}
Suppose a conjugation-invariant fixed Gaussian ray condition has the
checked finite-character expansion
\[
 \Lambda'_C(t)=\frac12\sum_\chi c_\chi T_\chi(t)+e(t),
 \qquad \sum_{t\leq X}|e(t)|\ll_{H_0}\sqrt X\log(2X)^2.
\]
Assume both the Gaussian estimate \eqref{eq:ga-ideal-input} and its
ordinary rational counterpart, with positive exponents $c_E,c_{\mathbb Q}$.
Let $G/\Gamma$ and $g$ satisfy the hypotheses of
Lemma~\ref{lem:ga-affine-interpolation}, and let $F_0$ be Lipschitz with
Haar mean zero.
Then any fixed linear combination of the $W$-tricked ray-prime weight
and the $W$-tricked ordinary-prime weight obeys, for every interval
progression $P\subseteq[N]$ of modulus at most $Q$,
\[
 \left|\sum_{n\in P}w_N(n)F_0(g(n)\Gamma)\right|
 \ll QN\log(2N)^7\delta^{c_W}\|F_0\|_{\mathrm{Lip}},
 \qquad c_W=a\min\{c_E,c_{\mathbb Q}\}>0,
\]
for mean-zero $F_0$, in a regime
$\delta\geq N^{-\sigma_1}$ and $Q\leq\delta^{-a}$ with fixed
$\sigma_1>0$. Constants may depend on $H_0,W$ and the fixed linear
combination. In particular this applies to the proposed centering
\[
 w_N(n)=\frac{\varphi(W)}W\Lambda'_C(Wn+b)
             -\mu_C(b)\lambda_{b,W}(n),
\]
once the normalization of $\lambda_{b,W}$, the scalar $\mu_C(b)$ and
the displayed character expansion have been fixed and checked.
\end{corollary}
\begin{proof}
Apply Lemma~\ref{lem:ga-wtrick-ma} to the finitely many characters and
to the rational von Mangoldt function. The Gaussian factor $1/2$
accounts for the two conjugate prime ideals above an eligible split
rational prime. Replacing ordinary von Mangoldt by its prime-only
version also has a proper-prime-power error
$O_W(\sqrt N\log(2N)^2)$. Thus the stated estimate holds with the
additional error
$O_{H_0,W}(\sqrt N\log(2N)^2\|F_0\|_\infty)$.
Choose $\sigma_1\leq\sigma$ so small that $\sigma_1c_W<1/4$.
Then $N\delta^{c_W}\geq N^{3/4}$ throughout the claimed regime, and
the square-root errors are absorbed. All normalization factors are
fixed constants because $W$ is fixed. This proves the minor-arc
estimate but does not assert the centered progression estimate
\eqref{eq:gf-sw}, which must be proved separately from the arithmetic
main terms.
\end{proof}

\begin{remark}\label{rem:ga-choice-order}
Fix the arithmetic moduli first, then the positive exponents in this
subsection. They do not depend on the later factorization exponent
$B$. Proposition~\ref{prop:gf-factorization} chooses $B$ large enough
to satisfy the resulting upper complexity condition. Its output
$M$ is polylogarithmic in $N$, so its equidistribution parameter also
obeys $\delta\geq N^{-\sigma_1}$ for sufficiently large $N$.
The affine interpolation argument is motivated by the corresponding
step of Frei--Tinkova, but uses no theorem from that work as an input.
\end{remark}


% Section fragment. Bibliography: GTOrbits2012, GTMobius2012.
\subsection{Quantitative factorization for centered weights}
\label{sec:gf-factorization}

We isolate the analytic implication needed after the arithmetic
progression main terms have been established. Throughout,
$L=\log(2N)$, all nilpotent groups are connected and simply connected,
and all filtrations are by closed connected rational subgroups. Constants
may depend on fixed arithmetic moduli and on the fixed dimension and
degree bounds.

\begin{proposition}[Factorization implication]\label{prop:gf-factorization}
Fix dimension and degree bounds $m,d$. Let $w_N:[N]\to\mathbb C$ satisfy
the following three hypotheses.
\begin{enumerate}
\item For a fixed $K_w\geq0$, $|w_N(n)|\ll L^{K_w}$.
\item For every $A_0,C_0>0$, uniformly over all interval progressions
$P\subseteq[N]$ of common difference at most $L^{C_0}$,
\begin{equation}\label{eq:gf-sw}
 \left|\sum_{n\in P}w_N(n)\right|
       \ll_{A_0,C_0}N L^{-A_0}.
\end{equation}
No lower bound on $|P|$ is imposed here.
\item There are fixed $c>0$, $C_1,K_1\geq0$ such that if a nilmanifold
of dimension at most $m$ and filtration degree at most $d$ has a
$Q$-rational adapted Mal'cev basis, $h$ is a polynomial sequence totally
$\delta$-equidistributed on $[N]$, and $F_0$ is Lipschitz with Haar mean
zero, then
\begin{equation}\label{eq:gf-ma}
 \left|\sum_{n\in P}w_N(n)F_0(h(n)\Gamma)\right|
 \ll N\delta^c Q^{C_1}L^{K_1}\|F_0\|_{\mathrm{Lip}}
\end{equation}
for every interval progression $P\subseteq[N]$ of modulus at most $Q$
and length at least $N/Q$. It is sufficient for this hypothesis to hold
only when
\begin{equation}\label{eq:gf-ma-regime}
 \delta\geq N^{-\sigma},\qquad Q\leq\delta^{-\kappa}
\end{equation}
for fixed positive $\sigma,\kappa$. An estimate allowing arbitrary
progression lengths is stronger than required.
\end{enumerate}
Then, for every $A,C_{\mathrm{init}}>0$, uniformly over initial
nilmanifolds of dimension at most $m$ and degree at most $d$ with an
$L^{C_{\mathrm{init}}}$-rational adapted basis, all polynomial sequences
$g$, and all Lipschitz $F$,
\begin{equation}\label{eq:gf-conclusion}
 \left|\sum_{n\leq N}w_N(n)F(g(n)\Gamma)\right|
     \ll_{A,C_{\mathrm{init}}}N L^{-A}\|F\|_{\mathrm{Lip}}.
\end{equation}
\end{proposition}

\begin{proof}
Normalize $\|F\|_{\mathrm{Lip}}=1$ and put
\[
 D=A+K_w+3,\qquad M_0=L^{\max\{1,C_{\mathrm{init}}\}}.
\]
Leave a factorization exponent $B$ unspecified for the moment. By
\cite[Theorem 1.19]{GTOrbits2012}, there is an integer $M$ with
\[
 M_0\leq M\ll_{B,m,d}M_0^{E(B,m,d)}
\]
and a factorization $g=\varepsilon g'\gamma$ into polynomial sequences.
Here $\varepsilon$ is $(M,N)$-smooth, $\gamma(n)\Gamma$ has period
$q\leq M$, and $g'$ is totally $M^{-B}$-equidistributed in a rational
subgroup quotient. The new adapted basis consists of $M$-rational
combinations of the initial one.

For each residue $j$ modulo $q$, choose a reduced rational representative
$\gamma_j$ of $\gamma(j)\Gamma$. Split that residue class into
$\lceil ML^D\rceil$ consecutive blocks of nearly equal cardinality.
For sufficiently large $N$ after $B$ is fixed, the resulting progressions
$P_{j,k}$ satisfy
\begin{equation}\label{eq:gf-cells}
 \#\{P_{j,k}\}\ll M^2L^D,\qquad
 |P_{j,k}|\gg\frac{N}{M^2L^D},\qquad
 \operatorname{diam}(P_{j,k})\ll\frac{N}{ML^D}.
\end{equation}
Their common difference is $q$. Indeed each original residue class has
cardinality comparable to $N/q$, and rounding the balanced block sizes
changes these estimates only by an absolute factor, since all powers of
$M$ are polylogarithmic in $N$ once $B$ is fixed.

Choose $n_{j,k}\in P_{j,k}$ and put $a_{j,k}=\varepsilon(n_{j,k})$.
The right-invariant ambient metric and smoothness imply
\begin{align*}
 &|F(\varepsilon(n)g'(n)\gamma_j\Gamma)
       -F(a_{j,k}g'(n)\gamma_j\Gamma)|\\
 &\hspace{1cm}\leq d(\varepsilon(n),\varepsilon(n_{j,k}))
       \ll L^{-D}\qquad(n\in P_{j,k}).
\end{align*}
There is no factor depending on the coordinates of $g'(n)$ in this
estimate: right invariance removes it. The total weighted freezing
error is therefore $O(NL^{K_w-D})=O(NL^{-A-3})$.

We use the quantitative conjugation claim proved in
\cite[p.~547]{moebius}, whose inputs are the rational-subgroup,
rational-element and metric comparisons of
\cite[Appendix A]{GTOrbits2012}. There are fixed $a>0$, $C_2\geq1$,
depending only on $m,d$, such that, with
\[
 H_j=\gamma_j^{-1}G'\gamma_j,\quad
 \Lambda_j=H_j\cap\Gamma,\quad
 h_j(n)=\gamma_j^{-1}g'(n)\gamma_j,
\]
the quotient $H_j/\Lambda_j$ has an adapted basis of rationality at most
$M^{C_2}$, $h_j$ is totally $M^{-aB+C_2}$-equidistributed, and
\[
 F_{j,k}(x\Lambda_j)=F(a_{j,k}\gamma_jx\Gamma)
 \quad\hbox{satisfies}\quad
 \|F_{j,k}\|_{\mathrm{Lip}}\leq M^{C_2},\qquad
 \|F_{j,k}\|_\infty\leq1.
\]
We have enlarged $C_2$ to absorb fixed multiplicative constants. These
exponents do not depend on $B$. In applying this published claim, one
must account for the commensurable lattice resulting from rational
conjugation; conjugation is not asserted to be an isometry of the
original quotients. Its finite-index and metric losses are included in
$a,C_2$.

Write $F_{j,k}=m_{j,k}+F_{j,k}^0$ by subtracting its Haar mean on
$H_j/\Lambda_j$. Then $|m_{j,k}|\leq1$,
$\int F_{j,k}^0=0$, and
$\|F_{j,k}^0\|_{\mathrm{Lip}}\leq2M^{C_2}$. This subtraction is needed
even if the initial function $F$ happened to have Haar mean zero.
For a dimension-zero factor the centered function simply vanishes.

Put
\[
 U=C_2+D+5,\qquad Q_*=M^U,\qquad b_*=aB-C_2.
\]
Since $M\geq L$, $Q_*$ dominates the rationality, cell modulus and
$N/|P_{j,k}|$. Applying \eqref{eq:gf-ma} to all centered cell functions
and using \eqref{eq:gf-cells} gives a total of at most
\begin{equation}\label{eq:gf-total-ma}
 C N L^{K_1}M^{-b_*c+C_1U+C_2+2+D}.
\end{equation}
Every cell is counted here. In particular, the bound $N$ in the
minor-arc hypothesis has not been replaced by an unsupported cell
cardinality.

We now fix $B$ sufficiently large that
\begin{equation}\label{eq:gf-choice-B}
 b_*>1,\qquad
 b_*c\geq C_1U+C_2+2+D+A+K_1+4,
 \qquad b_*\kappa\geq U+1
\end{equation}
when the restricted regime \eqref{eq:gf-ma-regime} is used. The last
condition is omitted for an unrestricted minor-arc estimate. These
inequalities involve only fixed exponents and not $E(B,m,d)$.
They make \eqref{eq:gf-total-ma} $O(NL^{-A-4})$. They also ensure
$Q_*\leq(M^{-b_*})^{-\kappa}$. The other regime condition,
$M^{-b_*}\geq N^{-\sigma}$, holds for all sufficiently large $N$:
after fixing $B$ the upper bound on $M$ is a fixed power of $\log N$.

Only at this point use that upper bound, writing $M\leq L^{E_*}$ after
enlarging the fixed exponent $E_*$. There are at most
$O(L^{2E_*+D})$ cells of modulus at most $L^{E_*}$. Apply
\eqref{eq:gf-sw} with
\[
 C_0=E_*+1,\qquad A_0=A+2E_*+D+4.
\]
Summing the constant terms $m_{j,k}\sum_{n\in P_{j,k}}w_N(n)$ gives
$O(NL^{-A-4})$. Combined with the centered estimates and freezing error,
this proves \eqref{eq:gf-conclusion}.
\end{proof}

\begin{remark}[Order of parameter choices]\label{rem:gf-order}
The order in the proof is
\[
 (A,K_w,C_{\mathrm{init}},c,C_1,K_1,a,C_2,\kappa,\sigma)
 \longrightarrow (D,U)\longrightarrow B\longrightarrow E_*
 \longrightarrow(A_0,C_0)\longrightarrow N_0.
\]
Thus there is no requirement that $B$ dominate the factorization upper
exponent $E(B,m,d)$. The centered progression estimate allows every
logarithmic saving, so its exponent can be chosen after $E_*$ is known.
The minor-arc input must have a fixed power form such as
\eqref{eq:gf-ma}; unspecified polylogarithmic specializations alone
would not justify this argument.
\end{remark}


\begin{corollary}[Fixed-modulus centered nilsequence estimate]
\label{cor:ga-fixed-centered-nilsequence}
Fix the Gaussian ray field and a conjugation-invariant union $C$ of ray
classes as above. Let $W$ be a fixed multiple of the conductor of
$K^{\mathrm{ab}}/\Q$, and let $b$ be a unit residue modulo $W$.
For every $A,C_{\mathrm{init}}>0$ and fixed dimension and degree bounds,
uniformly over the nilmanifolds and polynomial sequences in
Proposition~\ref{prop:gf-factorization} with an
$(\log(2N))^{C_{\mathrm{init}}}$-rational adapted basis,
\[
 \left|\sum_{n\leq N}h_{C,b,W}(n)F(g(n)\Gamma)\right|
 \ll_{H,W,A,C_{\mathrm{init}},m,d}
 N(\log(2N))^{-A}\|F\|_{\mathrm{Lip}}.
\]
The constants are uniform in $b$ and in the polynomial orbit. No
uniformity as $W$ varies is asserted.
\end{corollary}
\begin{proof}
We check the three hypotheses of
Proposition~\ref{prop:gf-factorization}. The pointwise inequality
$|h_{C,b,W}|\leq\lambda_{b,W}$ gives its first hypothesis with $K_w=1$.
Corollary~\ref{cor:gaussian-centered-progressions} gives its second
hypothesis for every $A_0,C_0$, in fact for every interval progression.

For the third hypothesis, take $H_0=H$ in
Corollary~\ref{cor:ga-centered-primes}. Its character expansion is
\eqref{eq:gaussian-ray-fourier}, with the factor $1/2$ and the
prime-power error supplied by \eqref{eq:gaussian-ray-prime-factor}
and the bound following that equation. Its two assumed minor-arc
estimates are Proposition~\ref{gi-minor-arc} and
Corollary~\ref{gi-ordinary-primes}. Apply that corollary to the fixed
linear combination defining $h_{C,b,W}$ in
\eqref{eq:gaussian-centered-weight}. It gives \eqref{eq:gf-ma} with
$c=c_W$, $C_1=1$ and $K_1=7$ in the regime
$\delta\geq N^{-\sigma_1}$, $Q\leq\delta^{-a}$.
The affine-interpolation result initially restricts $\delta$ to a
fixed sufficiently small interval. For larger $\delta$, the bound
$|h_{C,b,W}|\ll_W\log(2N)$ gives \eqref{eq:gf-ma} directly after
enlarging its fixed constant. For a dimension-zero quotient the
mean-zero test function vanishes. For all other dimensions and degrees
within the fixed bounds, take the minimum of the finitely many positive
exponents $c_W,a,\sigma_1$ and the maximum of the fixed constants.
Thus the restricted regime \eqref{eq:gf-ma-regime} holds with fixed
positive $\kappa,\sigma$, uniformly throughout the required class.

Proposition~\ref{prop:gf-factorization} now proves the displayed
estimate. There are only finitely many unit residues modulo the fixed
$W$, so their constants can be bounded by a single constant. In
particular, division by $N$ and the relabeling $W=V$, $N=L$ give the
fixed-complexity estimate \eqref{eq:gaussian-centered-nil-input}.
\end{proof}


% The analytic bridge below uses the fixed-modulus centered nilsequence
% estimate just proved for the Gaussian ray fields under consideration.
\subsection{From Gaussian nilsequence estimates to the polynomial count}
\label{subsec:gaussian-gowers-bridge}

Fix the Gaussian ray field $K/\Q(i)$ and the conjugation-invariant ray
classes used in the preceding argument. Let $D$ be divisible by the
conductor of $K^{\rm ab}/\Q$. For each relevant union $C$ of classes,
let $\Lambda'_C$ be its rational prime weight and let $\mu_C(b)$ be the
conditional density in \eqref{eq:v5-mu}. In particular
$0\leq\mu_C(b)\leq1$. We use the following conclusion of the preceding
Gaussian norm, progression and factorization arguments, as assembled in
Corollary~\ref{cor:ga-fixed-centered-nilsequence}: for every fixed
multiple $V$ of $D$, every unit $b\pmod V$, every fixed bounded-complexity
nilsequence $F(g(n)\Gamma)$, and every $B>0$,
\begin{equation}\label{eq:gaussian-centered-nil-input}
\begin{aligned}
 \left|\mathbb E_{n\leq L}h_{C,V,b}(n)F(g(n)\Gamma)\right|
 &\ll_{V,B,\mathrm{complexity}}(\log L)^{-B},\\
 h_{C,V,b}(n)&=\frac{\varphi(V)}V
 \bigl(\Lambda'_C(Vn+b)-\mu_C(b)\Lambda'(Vn+b)\bigr).
\end{aligned}
\end{equation}
The bound is uniform in the polynomial orbit and functions with the stated
complexity bound. The field, its ray modulus and $V$ are fixed in this
statement. We do not assume any uniform rate as $V$ varies.

\paragraph{Published transference inputs.}
We use Green--Tao, \emph{Linear equations in primes}~\cite{greentao},
Proposition~6.4 and Appendix~D for the ordinary prime majorant,
Proposition~7.1 and Appendix~C for its generalized von Neumann bound,
and Lemma~4.4 with its subsequent normal-form reduction.
These are statements about the ordinary prime majorant and arbitrary
functions dominated by it; they do not assert an Artin-prime theorem.
We also use Dodos--Kanellopoulos~\cite{v5dk}, Theorem~5.1, in the
following form, checked in arXiv:1610.00487v4, pp.~12--13. The final
typeset statement has not been compared. For target interval norm $U^s[L]$, $s\geq2$, its hypothesis on a
nonnegative majorant on a prime cyclic group of size
$L'\in[CL,2CL]$, $C\geq20$, is
\begin{equation}\label{eq:gaussian-dk-correct-exponent}
 \|\nu-1\|_{U^{2s}(\Z/L'\Z)}\leq\eta.
\end{equation}
For a real function $|f|\leq\nu$ with $\|f\|_{U^s[L]}\geq\delta$,
the conclusion is a correlation bounded below with a bounded Lipschitz
linear nilsequence on an $(s-1)$-step nilmanifold. Its nilmanifold,
Lipschitz bound and positive correlation threshold depend only on
$s,C,\delta$. The interval norm is the normalized zero-extension norm.
In particular a target $U^{u+1}$ requires majorant norm $U^{2u+2}$.

\paragraph{An auxiliary cutoff, distinct from the polynomial cutoff.}
Put
\[
 V_z=D\prod_{\substack{p\leq z\\p\nmid D}}p.
\]
Here $z$ is an auxiliary cutoff used only to prove the next linear
replacement statement. It is not the cutoff $v$ defining $W,A$ in
\eqref{eq:v5-global-gowers}.
For every fixed integer $s\geq2$, one can choose $z=z(L)\to\infty$
sufficiently slowly that
\begin{equation}\label{eq:gaussian-refined-gowers}
 \max_{C,\ (b,V_z)=1}\|h_{C,V_z,b}\|_{U^s[L]}=o(1).
\end{equation}
The maximum is over the fixed finite list of ray-class unions and the
finitely many unit residues for the current $V_z$.

Here are all normalization and limit details. First choose the complexity
of the ordinary prime majorant larger than both the requirements of the
intended generalized von Neumann bound and
$\max(4^s,4s)$. Choose its ambient constant $C\geq20$ large enough for
both published majorant and generalized von Neumann statements. If
$\lambda_{V,b}(n)=(\varphi(V)/V)\Lambda'(Vn+b)$, then
$|h_{C,V,b}|\leq\lambda_{V,b}$. The majorant simultaneously dominates
a fixed number of such ordinary weights on $[L^{3/5},L]$, with a
multiplicative constant depending only on the chosen complexities.
Consequently there is $\alpha>0$, depending only on the fixed
complexities and the fixed factor in the modulus, such that
$\alpha 1_{[L^{3/5},L]}h_{C,V_z,b}$ is dominated in absolute value by
$\nu$ on $[L]$. The same $\alpha$ works for all cutoffs and unit residues.
In applying the inverse theorem to a lower bound $\delta$ for the
unscaled norm, use the lower bound $\alpha\delta/2$ for this truncated,
scaled function, once the truncation error is at most $\delta/2$.

For sufficiently large $z$, $V_z=E\prod_{p\leq z}p$ with $E$ fixed.
The ordinary divisor-sum construction works with this $V_z$: at primes
dividing $E$, the unit residue excludes every divisor term containing
that prime, and at all other primes its local factors are unchanged.
The same observation applies to the correlation upper bounds in
Appendix~D. Thus the stated pseudorandomness conditions remain valid,
uniformly in unit residues, with constants allowed to depend on $E$.
Expand the centered $(2s)$-cube of $\nu-1$. It contains at most $4^s$
distinct forms, in $2s+1$ variables, with coefficients zero or one.
Every subproduct has mean $1+o(1)$ by the linear-forms condition;
inclusion--exclusion gives \eqref{eq:gaussian-dk-correct-exponent} with
a left side tending to zero. This estimate is on the full cyclic group.

The removed initial interval contributes $o(1)$ to the interval $U^s$
norm. Indeed the weights are $O(\log L)$ if $V_z=L^{o(1)}$, and a cube
with its initial vertex in this interval has relative count
$O(L^{-2/5})$. The $2^s$th power of the norm of the removed function is
at most $O((\log L)^{2^s}L^{-2/5})$. Its correlation with a bounded
nilsequence is also $o(1)$.

If \eqref{eq:gaussian-refined-gowers} fails by some fixed $\delta>0$,
Theorem~5.1 gives a fixed-complexity nilsequence correlation bounded
away from zero, contradicting \eqref{eq:gaussian-centered-nil-input}.
There is no illicit varying-modulus use of that input. For each integer
$z$ and each of the finitely many residues modulo $V_z$, first take a
size threshold for the required nilsequence estimates. Diagonalize
these thresholds also over $\delta=1,1/2,1/3,\ldots$ and over the finite
list of norm orders needed. Impose in addition
$z(L)\leq\log\log\log L$ and all size thresholds of the majorant.
At stage $j$, choose the cutoff $z_j$ large enough for the limiting
majorant errors to be below half the inverse-theorem tolerances needed
for the first $j$ tests; only then choose their size threshold $L_j$.
The resulting increasing step function tends to infinity and proves
the assertion for all sufficiently large $L$. The fixed smooth metric
in the inverse theorem is comparable to a fixed Mal'cev metric, and a
linear orbit with arbitrary base point is a polynomial orbit, so the
nilsequence input applies to exactly the supplied correlations.
When this construction is used at the refined scale $L\asymp N/V_{z_j}$
below, choose the physical threshold $N_j$ so large that every constant
multiple of $N/V_{z_j}$ used by the normal-form reduction is at least
$L_j$ for $N\geq N_j$. Impose the same growth constraints there.
Thus the cutoff is constant on each physical-scale stage; no implicit
equation $z=z(N/V_z)$ is being assumed.

\begin{lemma}[Fixed-system linear replacement]
\label{lem:gaussian-linear-replacement}
Let $\Psi=(\psi_1,\ldots,\psi_t):\Z^d\to\Z^t$ have fixed nonzero integer
linear parts of finite complexity. Allow its constant terms to vary
with $N$, subject to $\|\Psi\|_N\leq L_0$ for a fixed $L_0$.
Uniformly over convex regions $\mathcal K\subseteq[-N,N]^d$ with
$\Psi(\mathcal K)\subseteq[1,L_0N]^t$,
\begin{equation}\label{eq:gaussian-linear-replacement}
 \sum_{\mathbf n\in\mathcal K\cap\Z^d}
 \left(\prod_{i=1}^t\Lambda'_{C_i}(\psi_i(\mathbf n))
 -\prod_{i=1}^t\mu_{C_i}(\psi_i(\mathbf n)\bmod D)
                         \Lambda'(\psi_i(\mathbf n))\right)=o(N^d).
\end{equation}
The same conclusion holds for a mixed product in which any selected
factors have already been replaced by their ordinary-prime models.
\end{lemma}
\begin{proof}
We can discard points where some $\psi_i\leq N^{9/10}$.
For a nonconstant fixed linear form their number is
$O_{\Psi,L_0}(N^{d-1}(N^{9/10}+1))$; the product of weights is
$O((\log N)^t)$. This is $o(N^d)$, uniformly in the constant terms.
Apply the auxiliary $V_z$-trick to the variables:
$\mathbf n=V_z\mathbf m+\mathbf r$.
If one output residue is not a unit modulo $V_z$, its prime weight
vanishes on the remaining range, since the output is larger than
every prime dividing $V_z$. For an admissible residue vector write
\[
 \psi_i(V_z\mathbf m+\mathbf r)
       =V_z\widetilde\psi_{i,\mathbf r}(\mathbf m)+b_{i,\mathbf r},
       \qquad1\leq b_{i,\mathbf r}\leq V_z.
\]
The new linear coefficients are the original ones. Its normalized
constant terms remain bounded at scale $N/V_z$. Because $D\mid V_z$,
the factor $\mu_{C_i}$ is constant on this residue vector.

Telescope the difference of the two products. After normalization by
$(\varphi(V_z)/V_z)^t$, each term has one factor $h_{C_i,V_z,b_i}$;
all other factors are nonnegative and at most
$\lambda_{V_z,b_j}$. Introduce a common ordinary-prime majorant for
these at most $t$ weights and truncate as above. Multiply every
factor by a common fixed $\alpha_0>0$ so that each has absolute value
at most that majorant, and apply the normal-form extension of
Green--Tao Lemma~4.4 and Proposition~7.1. Here $\alpha_0$ depends
only on the fixed system, the chosen complexities and the fixed
conductor, not on $z$ or the residue vector. The small factor has
the required norm by \eqref{eq:gaussian-refined-gowers}. Dividing
the resulting estimate by $\alpha_0^t$ changes only a fixed constant.
The normal-form extension adds only a fixed number of variables with
fixed coefficients; its convex-region reduction is the one following
Lemma~4.4, and divides out their box volume. The error on each residue
vector is therefore $o((N/V_z)^d)$, uniformly in the vector and convex
region. The discarded initial segment stays outside the retained
output range since $V_z=N^{o(1)}$.

One can make this reduction explicit. For each pair $i\ne j$ choose
an integer vector $v_{ij}$ with $\dot\psi_j(v_{ij})=0$ and
$\dot\psi_i(v_{ij})\ne0$; fixed nonproportional linear parts ensure
such vectors exist with fixed size. Replace $\mathbf m$ by
$\mathbf x+\sum_{i\ne j}y_{ij}v_{ij}$ and allow each new $y_{ij}$
to range over an integer interval of length comparable to $N/V_z$.
The block $y_{ij}$ for fixed $i$ certifies normal form. The region
defined by the original convex constraint on this sum and the new box
is convex. The integer change $(\mathbf x,\mathbf y)\mapsto
(\mathbf x+\sum y_{ij}v_{ij},\mathbf y)$ is bijective, so its sum
is exactly the original sum times the number of points in the new box.
Division by that number gives the asserted error. If there is only one
form, choose an integer vector $v$ with $\dot\psi_1(v)\ne0$ and replace
$\mathbf m$ by $\mathbf x+yv$. The same box and bijection argument
gives the reduction in this case.

No growing factor is lost when summing residue vectors. If
$\mathcal R_z$ is the set of admissible vectors, then
\begin{equation}\label{eq:gaussian-admissibility-mass}
 \frac{|\mathcal R_z|}{V_z^d}
       \left(\frac{V_z}{\varphi(V_z)}\right)^t
   =\prod_{p\mid V_z}\beta_p(\Psi)=O_{\Psi,L_0}(1).
\end{equation}
Exclude the finitely many primes dividing a chosen nonzero
coefficient of each linear part, and a chosen nonzero two-by-two
minor for each pair of linear parts. At every remaining prime each
output hyperplane has probability $p^{-1}$, and any two meet with
probability $p^{-2}$. Hence $\beta_p=1+O_t(p^{-2})$, uniformly in
the constant terms. This also covers $t=1$, when there are no pair
minors and the good-prime factor is exactly one.
The finitely many remaining factors are bounded by $(p/(p-1))^t$.
The Chinese remainder theorem proves the equality, also for the fixed
prime powers dividing $D$. Multiplying the per-residue error by this
bounded mass gives $o(N^d)$. The proof works unchanged for the mixed
products. All auxiliary thresholds can be chosen uniformly over bounded
constant terms and convex regions, because the published majorant and
generalized von Neumann conclusions have this uniformity.
\end{proof}

\paragraph{The exact local cube needed in the appendix.}
Now fix the appendix cutoff $v$, and put $W,A,n,R$ as in its proof.
This is a new limiting argument; the auxiliary cutoff in the proof of
Lemma~\ref{lem:gaussian-linear-replacement} has already been eliminated
from that lemma's statement. Let
\[
 f_{C,b}(x)=\frac{\varphi(W)}W\Lambda'_C(Wx+b)1_{R\leq x\leq n},
 \qquad g_{C,b}=f_{C,b}-\mu_C(b).
\]
For fixed $a\geq2$ and $C_1<C_2<C_3$, uniformly over
$q\leq A^{C_1}$ and $nA^{-C_3}\leq J\leq nA^{-C_2}$,
\begin{equation}\label{eq:gaussian-global-gowers}
 \|g_{C,b}\|_{U^a_{q[J]}}^{2^a}
       \ll_{a,C_1,C_2,C_3,\rho,D}\frac1{\log v}+o_v(1).
\end{equation}
The norm uses the difference multiset $q(h-h')$, $h,h'\in[J]$, on
$\Z/n\Z$.

To prove this, remove the cutoff $x<R$ at cost
$O((\log n)^{2^a}R/n)=o_v(1)$. Expand the centered cube and partition
its domain according to every wrap-around integer $m_\omega$.
The relevant affine forms are precisely
\[
 W\left(x+q\sum_{r=1}^a\omega_r(h_r-h'_r)-m_\omega n\right)+b,
 \qquad\omega\in\{0,1\}^a.
\]
They have pairwise nonproportional linear parts. Each region is convex;
there are $O_a(1)$ regions because $qJ/n<1$. For fixed $v$, the number
of possible $q$ is finite and $J/n$ stays in a fixed positive compact
interval. Apply Lemma~\ref{lem:gaussian-linear-replacement} to every
subproduct and region. Its error is $o_v(n^{2a+1})$ uniformly in $J$,
and division by $nJ^{2a}$ gives $o_v(1)$. Since every output is $b$
modulo $D$, all its conditional densities are the same $\mu_C(b)$.
Thus the centered Artin cube equals $\mu_C(b)^{2^a}$ times the
centered ordinary-prime cube, plus $o_v(1)$.

The ordinary cube is evaluated exactly as in Tao--Ziegler~\cite{v5tz},
Proposition~6, published pages 39--40, using the ordinary-prime linear
correlation theorem, for example the published formulation in
Bienvenu~\cite{v6bienvenu}, Theorem~1.1, pages 79--80.
The factors at $p\mid W$ cancel the normalization.
For $p\nmid Wq$, pairwise independent gradients give local factor
$1+O_a(p^{-2})$, uniformly in the wrap constants. For $p\mid q$,
$p>v$, the factor is $1+O_a(p^{-1})$. Since
$\log q=O_{C_1,\rho,D}(v)$, there are $O(v/\log v)$ such prime divisors.
Their product is $1+O(1/\log v)$, and the other large-prime factors
give $1+O(1/v)$. Centering cancels the common main term. Prime powers,
if the ordinary theorem is stated for $\Lambda$, contribute $o_v(1)$
by their $O(\sqrt n\log^2 n)$ total weight and the fixed linear-fiber
bound. This proves \eqref{eq:gaussian-global-gowers}, and hence supplies
\eqref{eq:v5-global-gowers} without an Artin linear-correlations theorem.

\paragraph{Completion by the published polynomial transfer.}
The remainder of the proof of Proposition~\ref{prop:v5-artin-patterns},
starting with its arbitrary-function dense-model step, now applies
unchanged to these fixed Gaussian ray fields. Specifically use the
published Tao--Ziegler~\cite{v5tz} Theorem~11 and Corollary~3, then Theorem~9 and
Theorem~13. The model for an arbitrary $f\leq\lambda_{W,b}$ is used;
one does not assert that an Artin weight has model one. Equation
\eqref{eq:gaussian-global-gowers} shows that its bounded model, minus
$\mu_C(b)$, has small global norms. The model is bounded by a fixed
constant depending on the chosen sieve exponent. Divide its difference
from $\mu_C(b)$ by a fixed bound for that difference before applying
Theorem~9, whose hypothesis is a bound by one. After restoring this
fixed factor, concatenation and model approximation give the required
local norm bounds. None of these scaling constants depends on $v$.

Keep the original parameter order: local box order and coefficient
bound first, then $B_0>C+2$, then the global order $a$, and finally
small $\rho$. Only then let $v$ increase. The quantities
$M_0=\lfloor T/W\rfloor$, $M_1=\lfloor M_0/A\rfloor$ and
$J=\lfloor A^{-2B_0a}M_1^2\rfloor$ satisfy the size hypotheses recorded
in the appendix when $T/\sqrt N\geq4\sqrt W/A$.
For each fixed $v$ the estimates are uniform for $T$ in that allowed
range. The final assertion $T=o(\sqrt N)$ is imposed only along the
subsequent diagonal $v=v(N)\to\infty$, not while $v$ is frozen.

After $u=Wx+b$, $w=Wm+c$, the polynomials are
$c_jWm^2+2c_jcm+(b+c_jc^2-b_j)/W$, with distinct leading coefficients.
The bounded-model transfer kills every term containing a centered
factor. Output-boundary strips have width $O(T^2)=o(N)$.
For completeness, dominate them first by strips of width $\epsilon N$.
Here is the needed localization of the ordinary majorant estimate.
This argument also covers negative coefficients $c_j$. For the true
integer count use an enlarged period $n_* = \lceil B N/W\rceil$, where
$B>2+\max_j|c_j|$ is fixed. We may assume $T\leq\sqrt N$.
Nonpositive outputs have prime weight zero, and every positive output
is at most $(1+\max_j|c_j|)N$. Discard also tuples for which some
positive output is at most $W\lceil n_*^\rho\rceil+W$. For each $w$
there are $O_k(W n_*^\rho+W)$ such $u$, so their total weight is
$O_k((W n_*^\rho+W)T(\log N)^k)=o_v(NT)$, uniformly in the
allowed $T$ range. All remaining true outputs have their ordinary
prime majorants on $[n_*]$ without wrap-around. Use the same
localization calculation on the original period $n$ to bound the
artificial wrapped terms of the cyclic count. The parameters $n_*$
and $n$ differ by a fixed factor, so both applications satisfy the
same polynomial-forms hypotheses. In the following calculation $n$
denotes either chosen ambient period.
Write $\nu_0(x)$ for the majorant factor with constant shift if there
is one, and set $\nu_0=1$ otherwise. There is at most one such factor.
Put $F(x)=\mathbb E_m\prod_i\nu_i(x+P_i(m))$, omitting that factor.
Tao--Ziegler, Theorem~15 in Appendix~A (published page~57), applied to
$\nu_0F$, $\nu_0F^2$ and $\nu_0$, gives
\[
 \mathbb E_x\nu_0(x)(F(x)-1)^2=o_{v\to\infty}(1)+o_v(1).
\]
In the square use independent variables $m,m'$; the differences
between all its nonconstant shift polynomials are nonconstant,
including those from distinct copies. Weighted Cauchy--Schwarz now
shows, for every interval $I$ in the $x$ variable,
\[
 \mathbb E_x1_I(x)\nu_0(x)F(x)
 =\mathbb E_x1_I(x)\nu_0(x)+o_{v\to\infty}(1)+o_v(1).
\]
For fixed relative interval length the one-dimensional divisor-sum
mean estimate gives
$\mathbb E_x1_I\nu_0=|I|/n+o_{v\to\infty}(1)+o_v(1)$;
it follows either directly from the ordinary majorant construction
or from its linear-forms estimate on that fixed interval.
The normalized prime weights are bounded by fixed multiples of
their sieve majorants, with constants depending on the fixed sieve
exponent. Thus a strip of width $\epsilon N$ has total weight
$[O(\epsilon)+o_{v\to\infty}(1)+o_{v,\epsilon}(1)]NT$, with that
fixed domination factor included in the implied constants.
The errors here have the same iterated order as the rest of the
polynomial transfer. Then take $\epsilon\downarrow0$ as described
below. This avoids a spurious factor $(\log N)^k$ from a pointwise
bound. The same positive majorant treats interval endpoints.

Choose increasing cutoffs $v_j$ for which the global/local errors and
$4\sqrt{W(v_j)}/A(v_j)$ tend to zero, and then choose size thresholds
valid for all the finitely many residues, step values and transfer
parameters at that cutoff.
At stage $j$ impose the boundary estimates for every
$\epsilon\in\{1,1/2,\ldots,1/j\}$, preserving all earlier tests.
Let $v(N)=v_j$ between successive thresholds and take
$\omega(N)\geq4\sqrt{W(v(N))}/A(v(N))$, tending to zero sufficiently
slowly. This gives the uniform assertion for
$\omega(N)\sqrt N\leq T=o(\sqrt N)$.
Indeed, for an arbitrary such $T$, first fix $k$; eventually
$T^2/N<1/k$ and the estimate for the strip of width $N/k$ applies.
Take the upper limit as $N\to\infty$ and then $k\to\infty$.
There is no requirement that $T^2/N<1/j(N)$ along the chosen stages.
The main term is precisely \eqref{eq:v5-artin-series}; the admissible
residue mass is bounded, so summation does not amplify errors.
The $u\in[N/2,N]$ version follows by using that convex interval directly,
or by subtraction with thresholds enlarged for $N/2$.

This bridge uses the separately established Gaussian nilsequence input
and published ordinary-prime transference. It uses no statement or
estimate from Frei--Tinkov\'a as a theorem.


\subsection{Fixed Artin conditions in quadratic patterns}
\label{subsec:v5-artin-transfer}



Here the field is fixed throughout the limiting argument. Let $K$ be the
ray class field of $\Q(i)$ for $H\Z[i]$, where $16\mid H$. Put
$G=\operatorname{Gal}(K/\Q)$, and let $C_1,\ldots,C_k$ be
conjugation-invariant subsets of $\operatorname{Gal}(K/\Q(i))$,
regarded as conjugacy-invariant subsets of $G$.
Let $D$ be divisible by the conductor of $K^{\rm ab}$.
For a unit $b\pmod D$, define
\begin{equation}\label{eq:v5-mu}
 \mu_j(b)=\frac{\#\{g\in C_j:g|_{K^{\rm ab}}=\sigma_b|_{K^{\rm ab}}\}}
                  {[K:K^{\rm ab}]},
 \qquad \sigma_b(\zeta_D)=\zeta_D^b,
\end{equation}
and set $\mu_j(b)=0$ for nonunits. Define $\Lambda'_j(p)=\log p$ for
unramified primes with arithmetic Frobenius in $C_j$, and zero otherwise.

\begin{proposition}[Fixed-field polynomial count]\label{prop:v5-artin-patterns}
There is a function $\omega(N)\to0$, depending on the fixed data, such that
whenever $\omega(N)\sqrt N\leq T=o(\sqrt N)$,
\begin{equation}\label{eq:v5-artin-patterns}
 \sum_{u\leq N}\sum_{1\leq w\leq T}
       \prod_{j=1}^k\Lambda'_j(u+c_jw^2)
       =\mathfrak S\,NT+o(NT),
\end{equation}
where
\begin{equation}\label{eq:v5-artin-series}
 \mathfrak S=
 \left(\frac D{\varphi(D)}\right)^k
 \mathbb E_{u,w\bmod D}\prod_j\mu_j(u+c_jw^2)
 \prod_{p\nmid D}\beta_p,
 \quad
 \beta_p=\mathbb E_{u,w\bmod p}
       \prod_j\frac p{p-1}1_{p\nmid u+c_jw^2}.
\end{equation}
The same statement on $N/2\leq u\leq N$ has main term
$\mathfrak SNT/2$.
\end{proposition}
\begin{proof}
The Gaussian linear replacement of
Lemma~\ref{lem:gaussian-linear-replacement} and its local-cube consequence
\eqref{eq:gaussian-global-gowers} supply the analytic input. We record
the normalization and polynomial-transfer parameters here.

Choose a small positive $\rho$ and first keep $v$ fixed. Put
\[
 W=D\prod_{\substack{p\leq v\\p\nmid D}}p,\quad
 A=W^{1/\rho},\quad n=\lfloor N/W\rfloor,\quad R=n^\rho,
\]
and, for unit $b\pmod W$, put
\[
 f_{j,b}(x)=\frac{\varphi(W)}W\Lambda'_j(Wx+b)1_{R\leq x\leq n},
 \qquad g_{j,b}=f_{j,b}-\mu_j(b).
\]
View these functions on $\mathbb Z/n\mathbb Z$, with representatives
$1,\ldots,n$. If $a\geq2$ and $C_1<C_2<C_3$ are fixed,
\eqref{eq:gaussian-global-gowers} gives
\begin{equation}\label{eq:v5-global-gowers}
 \|g_{j,b}\|_{U^a_{q[J]}}^{2^a}
       \ll_a(\log v)^{-1}+o_v(1),
 \quad q\leq A^{C_1},\quad
 nA^{-C_3}\leq J\leq nA^{-C_2}.
\end{equation}
The implicit constant may also depend on the other fixed complexities
and the fixed ray field. The norm, wrap-around regions, conditional
densities, and order of the auxiliary cutoff are all treated explicitly
in Section~\ref{subsec:gaussian-gowers-bridge}.

We now use the arbitrary-function dense-model theorem of
Tao--Ziegler~\cite{v5tz}. Their published Theorem~11 and Corollary~3
give a single bounded model accurate in both global Gowers norms and the
local polynomial box norms. Their Theorem~9 is the bounded-function
concatenation statement, and their Theorem~13 bounds the polynomial
average by one of those local norms. All majorants here are the ordinary
prime majorants, since $0\leq\Lambda'_j\leq\Lambda'$.
The fixed factor $D$ in $W$ is allowed in their divisor-sum construction:
for sufficiently large $v$, $W=E\prod_{p\leq v}p$ with a fixed $E$ whose
prime factors are at most $v$. A unit residue makes every small-prime
divisibility condition identically false, including at the extra prime
powers. At the remaining primes the divisor-sum local factors are
unchanged. Polynomial coefficient bounds absorb the fixed factor $E$.
Thus the polynomial-forms estimates for the majorant have the same
$1+o(1)$ conclusion.

The following choice records the size hypotheses of these results.
First choose their local box order $d'$ and a coefficient exponent $C$;
choose an integer $B_0>C+2$, then a sufficiently large global order $a$,
and finally a sufficiently small $\rho$. Set $\rho_* =\rho/B_0$ and
$\mathcal A=A^{B_0}$. Put
$M_0=\lfloor T/W\rfloor$, $M_1=\lfloor M_0/A\rfloor$.
For $T/\sqrt N\geq4\sqrt W/A$ one has
\[
 A^{-1}\sqrt n\leq M_0\leq\sqrt n,
 \qquad A^{-3}\sqrt n\leq M_1\leq\sqrt n.
\]
The model at cutoff $n^{\rho_*}$ dominates the function already truncated
at $n^\rho$. Corollary~3 supplies approximation for the degree-one local
step polynomials, whose coefficients are $O(A^C)$, and simultaneously for
$q\leq A^{B_0a}$ with
$J=\lfloor A^{-2B_0a}M_1^2\rfloor$.
In particular
\[
 nA^{-2B_0a-7}\leq J\leq nA^{-2B_0a}.
\]
This is within \eqref{eq:v5-global-gowers}. Let $f'_{j,b}$ be the
model supplied by Corollary~3 at cutoff $n^{\rho_*}$. Its bound is
$0\leq f'_{j,b}\leq B_*-1$, where $B_*\geq2$ depends on the fixed
ray field, complexities and on $\rho_*$, but not on $v,n,b$ or the polynomials
within the stated coefficient bounds. This is the bounded model of
Tao--Ziegler's Proposition~4: its construction first rescales the
prime weight by a fixed positive constant and then divides the model
by that constant. In particular, it need not be bounded by one.
Since $0\leq\mu_j(b)\leq1$, the function
\[
 F_{j,b}=\frac{f'_{j,b}-\mu_j(b)}{B_*}
\]
is bounded in magnitude by one. The triangle inequality, the global
part of Corollary~3, and \eqref{eq:v5-global-gowers} show that its
required global norms tend to zero under the diagonal limit below.
Apply Theorem~9 to $F_{j,b}$
with parameters $d=d_0=2$, $M=M_1$, $A=\mathcal A$, and global
order $a$. Its smallness threshold is independent of $\mathcal A$;
$B_*$ has already been fixed before $v$ is chosen. Multiplication
by $B_*$ and the local part of Corollary~3 therefore show that the
local norm of $g_{j,b}$ tends to zero. Theorem~13, with the original
parameter $A$ and coarse scale $M_0$, makes every polynomial average
containing a centered factor $o(1)$. These are the reasons for choosing
the complexities before $v$.

For the final $W$-trick write $u=Wx+b$ and $w=Wm+c$.
If $b_j$ is the residue of $b+c_jc^2$ modulo $W$, the resulting integer
polynomial is
\[
 c_jWm^2+2c_jcm+(b+c_jc^2-b_j)/W.
\]
Its leading coefficient differences are nonzero, and its coefficients
are $O(W)$. Expand the product into its means and centered factors and
apply the preceding conclusion. End intervals of length $O(T^2)=o(N)$
are controlled by the weighted second-moment localization and the
fixed-strip argument in Section~\ref{subsec:gaussian-gowers-bridge}. Summing residue classes does not
amplify the error: their normalized admissible mass is
$\prod_{p\mid W}\beta_p=O(1)$.
Indeed, outside the fixed coefficient-difference primes,
\[
 \beta_p=\left(\frac p{p-1}\right)^k
                \frac{(p-1)(p-k+1)}{p^2}=1+O_k(p^{-2}).
\]
This also proves convergence of the product in
\eqref{eq:v5-artin-series}.
Let $N$ tend to infinity for each fixed $v$ and then let $v$ tend to
infinity sufficiently slowly. A diagonal choice supplies $\omega(N)$
and yields \eqref{eq:v5-artin-patterns}. Subtract the corresponding
formula with upper limit $N/2$, enlarging $\omega$ if necessary, to
obtain the last assertion.
\end{proof}

\begin{remark}[Published numbering]\label{rem:v5-tz-numbering}
The labels in the published 2018 paper differ from those in arXiv
1603.07817v1. The relevant dictionary is
\[
\begin{array}{c|ccccc}
 \text{arXiv} &\text{Theorem 1.4}&\text{Theorem 2.5}&\text{Theorem 3.1}
               &\text{Corollary 3.9}&\text{Theorem 4.1}\\
 \text{published}&\text{Theorem 4}&\text{Theorem 9}&\text{Theorem 11}
               &\text{Corollary 3}&\text{Theorem 13}
\end{array}
\]
The published locations are respectively pages $4,21,28,38,41$.
The preprint's equations (3.5) and (4.1) are published equations (36)
and (57). The dense model used above is the theorem's assertion for an
arbitrary function bounded by the prime weight. The special model $1$
for the ordinary prime weight is not an assertion about Artin weights.
\end{remark}

\subsection{The ray-class lower bound}\label{subsec:v5-ray-density}

For ray class fields we use the published existence and reciprocity
statements in Kopp--Lagarias~\cite{v5kl}, Theorems~1.1, 1.3 and 7.2,
and the modulus exact sequence of Theorem~6.5, specialized to the maximal
order $\mathbb Z[i]$. There are no real places. This replaces an appeal
to unpublished course notes. In this principal ideal domain the ray group
for the finite modulus $H\mathbb Z[i]$ is
\begin{equation}\label{eq:v5-ray-group}
 (\mathbb Z[i]/H\mathbb Z[i])^\times/\{1,-1,i,-i\}.
\end{equation}
It also follows directly by sending a residue $r$ to the ideal class of
$(r)$: surjectivity is principality, and the kernel consists exactly of
unit residues. The Artin map sends this ideal class to arithmetic
Frobenius. Conjugation acts by $r\mapsto\bar r$.

Choose a fixed $L_1$ larger than all coefficient-difference primes and
larger than $11$. Let $H_0$ be divisible by $16\cdot7M$ and by all odd
primes at most $L_1$. Prescribe a unit residue $r\pmod{H_0}$ which is
$9$ modulo $16$, $1+3i$ modulo $7$, and $i$ modulo every other odd
prime dividing $H_0$. Compatible higher powers can be chosen arbitrarily.
The associated rational-prime Frobenius condition is the two-element
conjugacy class $C_0=\{h_r,h_{\bar r}\}$ in the Galois group $G_0$ of
the ray field $K_0$ over $\mathbb Q$.
The two elements are distinct: an equality of their ray classes would
require $r=\epsilon\bar r\pmod{H_0}$ for a Gaussian unit $\epsilon$;
modulo $16$ this forces $\epsilon=1$, contradicted modulo $7$.
The field is Galois over $\mathbb Q$, because conjugation preserves the
modulus and the defining ray subgroup.

Let $D_0$ be the conductor of $K_0^{\rm ab}$ and choose a unit
$b_0\pmod{D_0}$ inducing the restriction of $h_r$.
The value in \eqref{eq:v5-mu} is then
$\mu_0=|C_0|/|G_0'|>0$.
For $L\geq L_1$ let
$T_L=\prod_{L_1<\ell\leq L,\,\ell\nmid H_0}\ell$ and let $K_L$ be the
ray field for $H_0T_L\mathbb Z[i]$. For each new prime write
\[
 B_\ell=(\mathbb Z[i]/\ell\mathbb Z[i])^\times,
 \qquad N_\ell=\ker(B_\ell\xrightarrow{\mathrm N}\mathbb F_\ell^\times).
\]
The norm map is onto and $|N_\ell|=\ell-\chi(\ell)$.

\begin{lemma}[Abelian fibres and their surviving proportion]
\label{lem:v5-ray-fibres}
Put $N_L=\prod_{\ell\mid T_L}N_\ell$ and
$G_L=\operatorname{Gal}(K_L/\mathbb Q)$.
Then
\[
 N_L\subseteq G_L',\qquad
 G_L'/N_L\simeq G_0',\qquad
 K_L^{\rm ab}=K_0^{\rm ab}\mathbb Q(\zeta_{T_L}).
\]
Let $C_L$ consist of those elements above $C_0$ whose normalized Gaussian
residue $z$ satisfies $\mathrm N(z-1)\mathrm N(z+1)\not\equiv0\pmod\ell$
at every $\ell\mid T_L$. For the density $\mu_L$ of this set,
\begin{equation}\label{eq:v5-ray-density}
 \mu_L(b)\geq\mu_0
       \prod_{\ell\mid T_L}\left(1-\frac4{\ell-1}\right)
 \qquad (b\equiv b_0\pmod{D_0},\ (b,D_0T_L)=1).
\end{equation}
\end{lemma}
\begin{proof}
In the CRT presentation of \eqref{eq:v5-ray-group}, the unit group
embeds in the $H_0$ factor, because it already embeds modulo $16$.
Once the representative above $r$ or $\bar r$ is normalized, the kernel
of reduction to the base ray group is exactly $\prod B_\ell$.
There is no further quotient by the four units.

The map $z\mapsto z/\bar z$ from $B_\ell$ onto $N_\ell$ is surjective.
For split primes this is $(x,y)\mapsto(x/y,y/x)$.
For inert primes it is $z\mapsto z^{1-\ell}$ in the cyclic group
$\mathbb F_{\ell^2}^\times$.
These maps are commutators with complex conjugation, so $N_L\subseteq G_L'$.
The norm on each new residue factor gives the cyclotomic action
$\zeta_\ell\mapsto\zeta_\ell^{\mathrm N(z)}$.
This is also seen by applying arithmetic Frobenius to $\zeta_\ell$;
it is compatible with the ray map by reciprocity.
The reduction map gives $K_0\subseteq K_L^{N_L}$. The cyclotomic
action just computed is trivial on $N_L$, so also
$\mathbb Q(\zeta_{T_L})\subseteq K_L^{N_L}$.
The fields $K_0$ and $\mathbb Q(\zeta_{T_L})$ are disjoint over
$\mathbb Q$. Indeed, $K_0$ is unramified at every prime dividing
$T_L$, whereas the inertia groups at these primes generate the
Galois group of $\mathbb Q(\zeta_{T_L})/\mathbb Q$. Thus their
intersection is trivial. Moreover, the kernel computation above gives
\[
 [K_L^{N_L}:\mathbb Q]
 =[K_0:\mathbb Q]\prod_{\ell\mid T_L}\frac{|B_\ell|}{|N_\ell|}
 =[K_0:\mathbb Q]\varphi(T_L).
\]
This is the degree of $K_0\mathbb Q(\zeta_{T_L})$, proving
\[
 K_L^{N_L}=K_0\mathbb Q(\zeta_{T_L}),\qquad
 G_L/N_L\simeq G_0\times\prod_{\ell\mid T_L}\mathbb F_\ell^\times.
\]
Since $N_L\subseteq G_L'$, taking commutators in this quotient gives
$G_L'/N_L\simeq G_0'$ and
$K_L^{\rm ab}=K_0^{\rm ab}\mathbb Q(\zeta_{T_L})$.

Fix a nonzero norm $s\in\mathbb F_\ell^\times$.
For $z=a+bi$ in that norm fibre,
$\mathrm N(z\pm1)=s+1\pm2a$.
For each sign, vanishing fixes $a$ and leaves at most two values of $b$.
Thus at most four of the $\ell-\chi(\ell)$ elements are excluded.
The conditions are invariant under conjugation, so $C_L$ is a union of
rational conjugacy classes.
In every compatible abelian fibre its numerator in \eqref{eq:v5-mu} is
at least $|C_0|\prod(\ell-\chi(\ell)-4)$, and its denominator is
$|G_0'|\prod(\ell-\chi(\ell))$.
Their quotient proves \eqref{eq:v5-ray-density}.
\end{proof}

\begin{proposition}[Finite-cutoff lower bound]\label{prop:v5-cutoff-lower}
There is a constant $c>0$ independent of $L$ such that the singular series
in Proposition~\ref{prop:v5-artin-patterns} for the fixed field $K_L$ and
the set $C_L$ at every position satisfies
\[
 \mathfrak S_L\geq c(\log L)^{-4k}.
\]
\end{proposition}
\begin{proof}
At the base modulus retain just $u=b_0$, $w=0\pmod{D_0}$.
At each new prime retain all pairs on which the $k$ norm values are units.
Apply \eqref{eq:v5-ray-density} separately to each position.
The norm-fibre denominators cancel as shown in the previous lemma; a
factor inverse to the degree of $K_L$ is not left in the bound.
The normalized proportion of the retained pairs at each new prime is
$\beta_\ell=1+O_k(\ell^{-2})$. Each is positive, since $w=0$, $u\ne0$
is allowed. Their product, together with the tail Euler product, is
bounded below by a fixed positive number. Williams' theorem with modulus
$1$, or its product for the two reduced classes modulo $4$, gives
\[
 \prod_{\ell\mid T_L}(1-4/(\ell-1))^k
       \gg(\log L)^{-4k}.
\]
The base modulus contributes another fixed positive factor.
\end{proof}

\subsection{The final order of limits}\label{subsec:v5-diagonal}

\begin{proof}[Proof of Theorem~\ref{thm:v5-simultaneous-primes}]
Fix the coefficients, $M$, and the base modulus first. Choose one fixed
$L\geq\max(L_1,L_*)$ so large that
\[
 \frac C{L\log L}<\frac c{32(\log L)^{4k}},
\]
where $C,c$ are from the obstruction and lower-bound propositions.
Now fix $K_L,C_L$. Only after this choice do we use the asymptotic
polynomial theorem or let its small-prime cutoff tend to infinity.

Choose numbers $\tau_j\downarrow0$, all at most $\tau_0$.
Choose successively $N_j$ so large that for $N\geq N_j$
$\omega(N)\leq\tau_j/2$,
that $\tau_j\sqrt N\geq2$, and that the error term in
\eqref{eq:v5-obstruction-tail} is less than
$c\tau_jN^{3/2}/(32(\log L)^{4k})$.
This is possible for each fixed $\tau_j$.
For $N_j\leq N<N_{j+1}$ use $\tau(N)=\tau_j$ and define
$T=\lfloor\tau(N)\sqrt N\rfloor$.
Then $T=o(\sqrt N)$ and $T\geq\tau\sqrt N/2$.
Apply Proposition~\ref{prop:v5-artin-patterns} to this piecewise choice
of $T$. For all sufficiently large $N$ its lower count is at least
$cNT/(4(\log L)^{4k})$ and exceeds the full obstruction tail.
Hence a pattern remains with no forbidden odd
common divisor greater than $L$.
The ray conditions exclude every odd divisor at most $L$ from each
individual adjacent norm. At $7$ the adjacent norms of $1+3i$ are $2$
and $6$; at every other odd base prime the adjacent norms of $i$ are $2$.
Thus no forbidden odd common divisor remains at all.

The residue modulo $16$ gives, on writing $a=9+16A$ and $b=16B$,
\[
 (a+1)^2+b^2\equiv4\pmod8,\qquad
 (a-1)^2+b^2\equiv64\pmod{128}.
\]
These yield the exact two-adic valuations and hence the cross-sign gcd
$4$. The residue modulo $M$ gives coprimality with $M$.
Finally $u\geq N/2$ and $w\leq\tau\sqrt N$ imply
$u/w^2\geq1/(2\tau^2)\to\infty$.
For sufficiently small $\tau_0$, all $u+c_jw^2$ are between fixed positive
multiples of $N$, so they exceed $Y$ when $N$ is large.
The dilation at the start imposes $q\mid w$.
\end{proof}

% BEGIN VERSION 9 ADDITION
\section{The odd-product quartics in 523 and 540 variables}\label{sec:v9}

The following construction incorporates the November 7 contribution of
Gabriel Conant, developed with ChatGPT. We give its arithmetic argument,
including the connection to the results proved above. Variable counts
include the input. The older quartics remain as source and proof baselines.

\begin{theorem}[The Version 9 quartics]\label{thm:v9-quartics}
There are specified integer polynomials $P_{523}$ and $P_{540}$, of total
degree four in respectively 523 and 540 variables, such that
\[
 q\in\Z\ \Longleftrightarrow\ \exists v\in\Q^{522}:P_{523}(q,v)=0
 \ \Longleftrightarrow\ \exists w\in\Q^{539}:P_{540}(q,w)=0.
\]
The first has 35,163 nonzero monomials and largest coefficient magnitude
$64\cdot3^{2^{27}}$. Its coefficient table uses the fixed integer
$S_0=3^{2^{26}}$. The second has 35,213 nonzero monomials and largest
coefficient magnitude $3^{2048}+1005072314631333<2^{3247}$.
Its complete decimal table occupies 708,848 bytes. Both have content one.
\end{theorem}

The proof is the composition of the source equivalences, local and global
arguments, and exact polynomial certificates below. The intermediate
554-variable construction is included to separate the odd-product
argument from the additional formal-depth refinement. It is not a further
claimed improvement over the two polynomials in the theorem.

\subsection{Inherited inputs and the exact changes of convention}\label{subsec:v9-inputs}

Use the normalized curve and generator of
Section~\ref{subsec:v5-point-source}, with $P=P_*$, not a new generator.
The coefficient basis is $(1,\sqrt2)$ and conjugation sends
$a+b\sqrt2$ to $a-b\sqrt2$. At unramified finite places $v(p)=1$;
at ramified places the integral valuation has value group $\Z$.
Heights are absolute logarithmic Weil heights, with the local degree
weights divided by $[K:\Q]=2$. Throughout this section set
$\mathcal H=H_{19}$ of Theorem~\ref{thm:quartic-height19}.

The source is the system in Section~\ref{subsec:v5-complete-source},
with the following explicit changes. Replace the comparison targets
$[4]P_*,[4]Q_i$ by $P_*,Q_i$; replace the orders 16 and 52 by 12 and 40
before the further refinement below. Retain the sixteen-test product,
the same five contacts, the scalar norm tree with orders $(1,1,1,1,32,32)$,
the Boolean correction and the preparation-zero pole equations.
Change $\beta$ to $2X^2+3$; use the point $S-P$ only in the scalar and
physical-anchor blocks; and change the fixed coefficient $\mathsf S$
to $S_0=3^{2^{26}}$. The point-chart replacement, the three new scales,
and the final order-one/order-forty test list are given below. All
unchanged quotient guards and permitted zero values of pole witnesses
are retained. This description, the displayed equations below, and the
pinned rational straight-line program specify the source without a
choice of arithmetic witnesses.

\begin{center}\small
\begin{tabular}{p{.29\textwidth}p{.63\textwidth}}
\toprule
Input & Exact result used and translation\\\midrule
Elliptic labels and domains & Section~\ref{subsec:v5-point-source},
Lemma~\ref{lem:v5-point-domains}, and \eqref{eq:o8-indexheight}:
$2E(K)=\Z P_*$, nonzero coordinate domains, and the same lower bound
$h(x([n]P_*))\ge n^2-6$.\\
Height of rational contacts & Theorem~\ref{thm:quartic-height19},
with $\mathcal H<2^{2093966}$, applied through the fractional linear
map and contact calculation of Lemma~\ref{lem:v5-height-transfer}.
Only that calculation is reused; its exponent 24 is replaced by 19.\\
Pole obstruction and realization & Lemmas~\ref{hyp:o8-pole},
\ref{lem:original-seven-simultaneous-realization}, and
\ref{lem:v5-one-place}. These include the restricted-converse dependency
proved in the earlier arithmetic sections; it has not been discarded.\\
Scalar and anchor & Lemma~\ref{v5:scalar32-bridge}, with its proof
applied to $2X^2+3$ and $S-P$, and
Lemma~\ref{lem:v5-three-witness-anchor} with the smaller exponent
calculated just below.\\
Local comparisons & Lemma~\ref{lem:v5-unramified-log} and the
recurrence of Lemma~\ref{lem:quartic-reciprocal-log}; orders and
target changes are justified below, not inferred by renaming indices.\\
Global choices & Theorem~\ref{thm:v5-simultaneous-primes},
Lemma~\ref{lem:v5-two-reductions}, Proposition~\ref{prop:v5-artin-patterns},
Section~\ref{subsec:v5-diagonal}, and
Lemma~\ref{v5:real-progression-bridge}. All fixed data precede the
asymptotic prime choice.\\
Final contradictions & Lemma~\ref{lem:o8-lucas} and the explicit
new height inequalities below.\\\bottomrule
\end{tabular}
\end{center}

The published results behind these inputs retain the versions and
convention qualifications in the theorem-application inventory. The new
contribution does not supply an assumed replacement for the restricted
converse, the height theorem, or the Gaussian prime argument.

\begin{lemma}[Direct targets and smaller anchor]\label{lem:v9-bridge}
For the changed source above, the local comparison and completeness
arguments remain valid with direct targets. The scalar anchor on
$R=[a-1]P_*$ implies
\[
 0\ne a-1\in2\,3^{2^{24}-1}\Z,
 \qquad \log_2|a-1|>25165823.5,
 \qquad h(\beta)\le\mathcal H|a-1|^{19/32}.
\]
It permits arbitrarily divisible positive scalar indices $a-1$.
\end{lemma}
\begin{proof}
For direct targets set $V_0=H-P_*$ and $V_i=U_i-Q_i$, with
$H=[h]P_*$, $U_i=[u_i]P_*$ and $Q_i=[n_i]P_*$. The quotient guards
give finite nonidentity points. At a consumed prime the full logarithm
ratios are $u_i/h$ and $(u_i-n_i)/(h-1)$. One of $h,h-1$ is a unit.
As in Lemma~\ref{lem:v5-two-anchors}, the cross congruence first makes
both approximate ratios integral, and then makes their relative
errors absolute. Subtraction of the two integral congruences gives
$L_i\equiv n_i$, without division by four. The logarithm recurrence
and denominator bound in Lemma~\ref{lem:quartic-reciprocal-log} give
the same estimates at orders 12 and 40. The test list below strengthens
the input order to 40 in the final refinement.

For completeness, the reduced elliptic groups have all four rational
two-torsion points, since $p>76176$ and the three roots $0,\pm d$
are distinct. Thus their doubling images have order $d_\pm(p)/4$.
The point $P_*$ is itself a double, so these integers annihilate its
reductions. The prime theorem and Lemma~\ref{lem:v5-two-reductions}
therefore give, on setting $M_\pm=\operatorname{lcm}_p(d_\pm(p)/4)$,
\[
 \gcd(M_-,M_+)=1,\quad\gcd(M_-,L)=16,\quad\gcd(M_+,L)=1.
\]
Here $64\mid L$, the exact two-parts are $2^6$ and $2^2$ before
division by four, and all odd prime divisors of $L$ were avoided.
Consequently $h\equiv0\pmod{M_-}$, $h\equiv1\pmod{M_+}$ and
$h\equiv h_0\pmod L$ are compatible when $64\mid h_0$.
The real-progression argument supplies large positive choices and
avoids the finite forbidden point and quotient loci. The later
prime-separation lemma justifies adding the factor $\Pi^{19}$.

For the anchor, repeat the proof of
Lemma~\ref{lem:v5-three-witness-anchor} with $C=3^{2^{25}}$ and
$\mathsf S=C^2=S_0$. Its quaternion criterion and anisotropic
binary form are unchanged. Thus the two rational coefficients of
$f(R)/C$ belong to $\Z_{(2)}\cap\Z_{(3)}$. The exact local calculation
of Section~\ref{subsec:o8-anchor-subgroup} now gives
$a-1\in2\,3^{2^{24}-1}\Z\setminus\{0\}$, and
\[
 \log_2|a-1|\ge1+(2^{24}-1)\log_2 3
 >1+\tfrac32(2^{24}-1)=25165823.5.
\]
For the scalar proof the new label is at least three, its five-contact
product never vanishes, and all its negative valuations at the
consumed odd primes are even. These are exactly the properties of
$X^2+2$ used in Lemma~\ref{v5:scalar32-bridge}; they hold for
$2X^2+3$ by the calculation below. That proof therefore gives
$M_5(\beta)^{32}\mid a-1$. Applying
Theorem~\ref{thm:quartic-height19} gives the asserted exponent $19/32$.
Its finite correction and continuity arguments apply to the same
norm tree with the scalar point $R$; choosing $a-1$ sufficiently
divisible meets the anchor and any further prescribed divisibility.
\end{proof}


\subsection{An exact replacement of the point chart}
Put $K=\Q(\sqrt2)$, $\delta=3-\sqrt2$ and $D=\delta^2=11-6\sqrt2$. The old chart, with letters in $K$, is
\begin{equation}\label{v9chart:eq:old}
 u=t^2,\qquad x=h^2,\qquad u(x-u/4)=D,\qquad y=h(u/2-x).
\end{equation}
The new chart is
\begin{equation}\label{v9chart:eq:new}
 x=h^2,\qquad j^2=x^2-D,\qquad y=hj.
\end{equation}
An equation over $K$ means its two rational component equations. The last expression in \eqref{v9chart:eq:new} need not be assigned its own coordinates.

\begin{lemma}\label{v9chart:lem:squares}
If $U,V\in K$, $U\ne0$, and $V^2=U^3+4DU$, then $U$ is a square in $K$.
\end{lemma}
\begin{proof}
The ring $\mathcal O=\Z[\sqrt2]$ is a PID: rounding both coefficients of an element of $K$ to nearest integers leaves a remainder of absolute norm at most $1/2$, giving a Euclidean algorithm. Its units are $\pm(1+\sqrt2)^n$, and every totally positive unit is a square. For completeness, multiply a positive unit by a power of $1+\sqrt2$ to place its first real embedding in $[1,1+\sqrt2)$. Its other embedding is the positive or negative reciprocal. Writing the unit as $a+b\sqrt2$ then forces $a,b\in\{0,1\}$, and the norm condition and excluded upper endpoint leave only the unit~1. The two signs of the fundamental unit under the real embeddings give the assertion about totally positive units.

Let $m=v_{\mathfrak p}(U)$, with integral normalization. Away from $(\sqrt2)$ and $(\delta)$ the two terms $U^3$ and $4DU$ have valuations $3m$ and $m$. If $m$ is odd, their unequal minimum is odd, a contradiction. At $(\sqrt2)$ these valuations are $3m$ and $m+4$, with equality only at the even value $m=2$. Again $m$ cannot be odd. At $(\delta)$ they are $3m$ and $m+2$, with equality only at $m=1$. In that exceptional case write $U=\delta u$ with $u$ a local unit. Then
\[
 V^2=\delta^3u(u^2+4).
\]
The residue field is $\mathbf F_7$, in which $-4$ is a nonsquare. Thus the right side has valuation three, also impossible. Every finite valuation of $U$ is even.

Under each real embedding, $D>0$ and $U(U^2+4D)=V^2$, so $U>0$. By the PID property, $U=\gamma^2\eta$ for $\gamma\in K^*$ and a unit $\eta$. This unit is totally positive, hence a square.
\end{proof}

\begin{proposition}\label{v9chart:prop:chart}
The charts \eqref{v9chart:eq:old} and \eqref{v9chart:eq:new} have exactly the same projected pairs $(x,y)\in K^2$. Every solution of the new chart has $h,j,x,y\ne0$, and $x$ is positive under both real embeddings.
\end{proposition}
\begin{proof}
Given \eqref{v9chart:eq:old}, put $j=t^2/2-h^2$. Then
\[
 j^2=h^4-t^2h^2+t^4/4=h^4-D,
\]
and $y=hj$. Conversely, given \eqref{v9chart:eq:new}, set $U=2(h^2+j)$ and $V=2hU$. The identity $j^2=h^4-D$ gives
\[
 U^2-4h^2U+4D=0,
\]
hence $V^2=U^3+4DU$. Also $U\ne0$, since $U=0$ would imply $D=0$. By Lemma~\ref{v9chart:lem:squares}, choose $t\in K$ with $t^2=U$. It follows that
\[
 t^2(h^2-t^2/4)=D,\qquad h(t^2/2-h^2)=hj,
\]
as required.

If $h=0$, then $j^2=-D$, contradicting either real embedding. If $j=0$, then $h^4=D$, so $h^2=\pm\delta$. Taking norms yields $\Norm(h)^2=7$, impossible for $\Norm(h)\in\Q$. The remaining assertions follow immediately.
\end{proof}

The proof is specific to this field and this value of $D$. It does not use the generally false assertion that square abscissa characterizes doubles on every elliptic curve. It also avoids adding a Mordell--Weil rank or generator hypothesis to the chart replacement.

\subsection{Using the chart without ordinate variables}
All nine point modules use the new chart. It has six rational coordinates: the two components of each of $h,j,x$, subject to four quadratic component equations. The old module had ten coordinates, including the components of $t,h,u,x,y$. Thus the change saves four coordinates per point, provided each downstream ordinate is consumed as the expression $hj$.

For the difference of chart points $Q,R$, retain the existing inverse equations for $x_Q-x_R$ and define the slope by
\[
 (x_Q-x_R)m=h_Qj_Q+h_Rj_R.
\]
Its numerator and defining equation are quadratic in the available coordinates. Define the difference abscissa by the existing square of $m$, and its ordinate by
\[
 y_{Q-R}=m(x_Q-x_{Q-R})-h_Qj_Q.
\]
This again is a quadratic assignment. The retained inverse ensures that the slope has exactly its old value.

For a chart point,
\[
 z=-x/y=-h/j.
\]
Thus the local parameter is defined by $jz=-h$. When a logarithm value is combined with its parameter, replace $y\ell=-xG$ by $j\ell=-hG$, where $G$ is the unchanged truncated logarithm polynomial. At the scalar ratio replace
\[
 x_H\rho=-y_Hz_S\qquad\text{by}\qquad h_H\rho=-j_Hz_S.
\]
All cancellations are reversible by Proposition~\ref{v9chart:prop:chart}. Difference-point denominators and all other guards remain unchanged. The fixed target has
\[
 h=-1+\tfrac74\sqrt2,\qquad j=\tfrac{39}{8}-\tfrac92\sqrt2;
\]
these give exactly its previous abscissa and ordinate.

Together with Proposition~\ref{v9chart:prop:chart}, these identities give rational witness maps in both directions. To reconstruct the old chart, choose the square root supplied by Lemma~\ref{v9chart:lem:squares}; then evaluate all old assignments. The choice of sign does not alter $(x,y)$ or any downstream value.

\subsection{Shorter weight circuits}\label{subsec:v9-laurent}
For each of nine detectors, write $e$ for its nonzero weight and $q$ for its order-one weight. The required outputs are $q_{12}=qe^{22}$ and $q_{40}=qe^{78}$. The old construction uses seven pure powers of $e$ and two final quotients. Instead construct
\[
 a_1=e^{-1},\quad a_3=a_1^2/e,\quad a_7=a_3^2/e,\quad
 a_{14}=a_7^2,\quad a_{29}=a_{14}^2/e,\quad a_{55}=a_{29}^2/a_3,
\]
and impose
\[
 a_{29}q_{12}=qa_7,\qquad a_{55}q_{40}=q_{12}e.
\]
These are six power coordinates and two output coordinates. They give the same outputs and save one coordinate per detector.

The nonzero condition is inherited, rather than presumed from a quotient notation. In the source, $w_0$ has a retained inverse and
\[
 w_0n=(w_0^2-2w_1^2)/2,\qquad (e\Delta)b=n^2.
\]
Since $w_0\ne0$ and 2 is a nonsquare in $\Q$, the first equation gives $n\ne0$. The second gives $e\Delta\ne0$. Therefore every Laurent denominator above is nonzero. Conversely all Laurent coordinates are uniquely defined for each source solution. The changes introduce no new zero branches and exclude no source solution.

\subsection{Scalar preparation with one fewer coordinate}\label{subsec:v9-preparation}
The scalar pole uses a nonzero input $b$ and a preparation witness $t$, with
\[
 U=\frac{b(1-bt^2)}{1+bt^2}.
\]
The old circuit stores $b^{-1},t_2=t^2,bt_2,U$. Replace these four coordinates by $t_2,T,V$, with
\[
 t_2=t^2,\qquad b(T-t_2)=1,\qquad TV=2,
\]
and use the affine expression $U=V-b$. The equation $b(T-t_2)=1$ forces $b\ne0$, and $TV=2$ forces $T\ne0$. Consequently $T=t^2+b^{-1}$, $1+bt^2=bT\ne0$, and
\[
 U=2/T-b=\frac{b(1-bt^2)}{1+bt^2}.
\]
Conversely every old solution extends uniquely by these values of $T,V$. Thus the modification saves one coordinate and preserves all downstream quantities. It also supplies the nonzero condition formerly imposed by the separate coordinate $b^{-1}$.


\subsection{An exact product-height identity}\label{v9:sec:height}
Put
\[
 K=\Q(\sqrt2),\quad d=3-\sqrt2,\quad D=d^2=11-6\sqrt2,
 \qquad E:y^2=x^3-Dx,
\]
and put
\[
 R_0=(1+\sqrt2,2\sqrt2),\qquad
 P=2R_0=\left(\frac{57-28\sqrt2}{8},\frac{-660+417\sqrt2}{32}\right).
\]
The model calculation in Section~\ref{sec:degree8-original} proves that $P=P_*$ has infinite order and $2E(K)=\Z P_*$; this is the same normalized curve and generator, not the earlier model of Proposition~\ref{prop:generator}. Write $\operatorname{h}$ for the absolute logarithmic Weil height. The following lemma uses the displayed coordinates and doubling identity, together with the infinite order of $P$.

\begin{lemma}\label{v9:lem:product}
For odd integers $n_1,\ldots,n_r$,
\[
 \operatorname{h}\!\left(\prod_{i=1}^r x([n_i]P)\right)
 =\sum_{i=1}^r\operatorname{h}(x([n_i]P)).
\]
\end{lemma}
\begin{proof}
Let $T_0=(0,0)$. At a finite place $v$ with $v(D)=0$, use the finite quotient $E(K_v)/E_1(K_v)$, where $E_1$ denotes the formal neighborhood of infinity for this integral equation.
Here we use the integral formal group construction of
\cite[Chapter IV, Sections 1--3]{silverman}. Its parameter is $z=-x/y$,
and $x=z^{-2}(1+O(z))$, $y=-z^{-3}(1+O(z))$ identify this neighborhood
with the maximal ideal. It is an open subgroup; the proper curve
$E(K_v)$ is compact, so its quotient is finite. This argument does not
assume that the equation is minimal at the place above two. For a finite nonidentity point $Q$, membership in $E_1$ is equivalent to $v(x(Q))<0$. Since
\[
 x(Q+T_0)=-D/x(Q),
\]
positive valuation of $x(Q)$ is equivalent to $Q+T_0\in E_1$. The image of $T_0$ in this quotient has order two.

If an odd multiple of $P$ belongs to $E_1$, the image of $P$ has odd order. None of its multiples can then equal the image of $T_0$, so all $x([n_i]P)$ have nonpositive valuation. If no odd multiple belongs to $E_1$, all these abscissae have nonnegative valuation. Thus the valuations have a common weak sign. This argument includes the place above two.

The only place with $v(D)>0$ is $\mathfrak d=(3-\sqrt2)$ above seven. The displayed integral model is minimal there, since $v_{\mathfrak d}(-D)=2<4$. Its reduction is $y^2=x^3$ over $\F_7$. Substituting $\sqrt2=3$ shows that $P$ reduces to $(1,6)$, a nonsingular point. By \cite[Proposition VII.2.1, pp. 188--190]{silverman}, all multiples therefore belong to the subgroup $E_0(K_{\mathfrak d})$ of points with nonsingular reduction. A positive valuation of $x$ would force a positive valuation of $y$, and hence reduction to the singular point $(0,0)$. It cannot occur. The valuations again have a common weak sign.

At either real embedding, each $[n_i]P$ is a double of a non-torsion point. The duplication formula makes its abscissa a nonzero square. Its ordinate is nonzero, so the equation of $E$ gives $x([n_i]P)>d$ at that embedding. Both real values $3\pm\sqrt2$ of $d$ exceed one.

At a finite place, numbers whose valuations have the same weak sign satisfy
\[
 \log\max\left(1,\left|\prod_i x([n_i]P)\right|_v\right)
 =\sum_i\log\max(1,|x([n_i]P)|_v).
\]
At both real places the same equality holds because all the factors exceed one. Summing with normalized local weights proves the assertion.
\end{proof}

The product here is an element of $K$. It is not its norm down to $\Q$. Norms can cancel a pole at one split place against a zero at the other, whereas the proof above works separately at every place.

\begin{corollary}\label{v9:cor:scales}
Suppose the four root points $Q_i=[n_i]P$ have odd indices, and write
\[
 C=\prod_{i=1}^4x(Q_i)=C_0+C_1\sqrt2.
\]
The original nine height-detector scales may be replaced by
\[
 1,\qquad C_0^2,\qquad (1+C_1^2)^2.
\]
The root-height estimate retains the original constant $7\mathcal H$, where $\mathcal H$ is the exponent-19 height constant of the source.
\end{corollary}
\begin{proof}
Both real values of $C$ are positive, so $C_0>0$. All three scales are positive nonzero rational squares. If a nonzero integer $N$ is divisible by the $k$th powers of the three contact radicals, the unchanged detector estimate gives
\[
 \operatorname{h}(\lambda),\ \operatorname{h}(C_0^2\lambda),\
 \operatorname{h}((1+C_1^2)^2\lambda)
 \le T:=(\mathcal H+\log3)|N|^{19/k}.
\]
The elementary rational height identities yield
\[
 \operatorname{h}(C_0)\le T,\qquad \operatorname{h}(C_1)\le T/2.
\]
Indeed, height subadditivity first bounds the heights of the two scales by $2T$; also $\operatorname{h}(C_0^2)=2\operatorname{h}(C_0)$ and $\operatorname{h}((1+C_1^2)^2)\ge4\operatorname{h}(C_1)$. Hence
\[
 \operatorname{h}(C)\le\operatorname{h}(C_0)+\operatorname{h}(C_1)+\tfrac32\log2
 \le\tfrac32T+\tfrac32\log2
 <7\mathcal H |N|^{19/k}.
\]
Here $\mathcal H\ge1$. Lemma~\ref{v9:lem:product} bounds each root height by $\operatorname{h}(C)$. In particular, the inherited lower bound, with $B=\max_i|n_i|$, gives
\begin{equation}\label{v9:eq:height-transfer}
 B^2-6<7\mathcal H |N|^{19/k}.
\end{equation}
The old completeness proof uses its scales only after the roots have been fixed, and requires fixed positive rational squares. Thus its prime selection and pole construction apply to this shorter list.
\end{proof}

\subsection{Enforcing odd indices by quadratic equations}\label{v9:sec:chart}
Put $\epsilon=1+\sqrt2$. The existing chart for a point $Q$ has $K$-coordinates
\[
 x=h^2,\qquad j^2=x^2-D,\qquad Q=(x,hj).
\]
To this chart add $t,r\in K$ and the equations
\begin{equation}\label{v9:eq:odd-chart}
 t^2=2(x+j),\qquad th+x+j=\epsilon r^2.
\end{equation}
These are four rational quadratic equations in four additional rational coordinates. No extra inverse witnesses are needed.

\begin{proposition}\label{v9:prop:odd-chart}
The projected points of the augmented chart are exactly $P+4E(K)$, that is, the odd multiples of $P$.
\end{proposition}
\begin{proof}
Use the squareclass homomorphism
\[
 \alpha:E(K)\longrightarrow K^*/K^{*2},\qquad
 \alpha(O)=1,\quad \alpha(T_0)=[-D]=[-1],\quad
 \alpha((a,b))=[a]\quad(a\ne0).
\]
For clarity, its homomorphism property follows from the line-intersection identity
\[
 x(R+S)x(R)x(S)=
 \left(\frac{x(R)y(S)-x(S)y(R)}{x(R)-x(S)}\right)^2
\]
when the displayed abscissae are nonzero and distinct. Doubling follows from
$x(2R)=((x(R)^2+D)/(2y(R)))^2$; addition by $T_0$ follows from $x(R+T_0)=-D/x(R)$. Inverse pairs and identity summands are immediate. If $R+S=T_0$, then $x(R)x(S)=-D$. These cases cover the exceptional pairs.

Since $2E(K)=\Z P$ and $2R_0=P$, every point of $E(K)$ has a unique expression $mR_0+T$, where $m\in\Z$ and $T\in E[2](K)$. The two-torsion is generated by $T_0$ and $T_d=(d,0)$, so
\[
 \alpha(mR_0+aT_0+bT_d)=[\epsilon]^m[-1]^a[d]^b.
\]
The classes $[\epsilon],[-1],[d]$ are independent. A square $\epsilon^c(-1)^a d^b$, with $a,b,c\in\{0,1\}$, would have rational square norm $(-1)^c7^b$. This forces $c=b=0$, and the real embeddings then force $a=0$. Consequently
\begin{equation}\label{v9:eq:squareclass}
 \alpha(R)=[\epsilon]\quad\Longleftrightarrow\quad R\in R_0+2E(K).
\end{equation}

Given the augmented chart, set $u=t^2=2(x+j)$ and
\[
 a=th+u/2,\qquad b=at.
\]
Expansion gives
\[
 u(x-u/4)=D,\qquad a^2-ua=D,\qquad b^2=a^3-Da.
\]
In particular $t,a,b$ are nonzero. The tangent slope at $R=(a,b)$ is $h+t$; duplication gives $2R=(x,hj)=Q$. Equation~\eqref{v9:eq:odd-chart} says $a=\epsilon r^2$, so $\alpha(R)=[\epsilon]$. By~\eqref{v9:eq:squareclass}, $Q\in P+4E(K)$.

Conversely, for $Q=(2m+1)P$, take $R=(2m+1)R_0=(a,b)$. Then $a,b\ne0$, and $a=\epsilon r^2$ for some $r\in K$. Set
\[
 t=b/a,\qquad h=(a^2+D)/(2b),\qquad x=h^2,
 \qquad j=t^2/2-x.
\]
The same duplication identities give $Q=(x,hj)$, $j^2=x^2-D$, and both equations~\eqref{v9:eq:odd-chart}.
\end{proof}

Only the four root points receive this augmentation. The shared point and the four comparison points retain their old charts: their indices need not be odd. The total cost is sixteen coordinates and sixteen quadratic residuals. The projected-point formulation matters, since the added equations can select a sign representation of the old $h,j$ chart.

\subsection{Odd labels, scalar bounds, and completeness}\label{v9:sec:labels}
Write the root points as
\[
 A=[n]P,\qquad S=[a]P,\qquad T=[t]P,\qquad Z=[b]P.
\]
Here $S$ denotes a point; the fixed coefficient remains $S_0$. Replace the old input expression by
\[
 \beta=2X^2+3.
\]
For rational $X$, integrality of $\beta$ is equivalent to integrality of $X$: if $X=u/v$ in lowest terms, then $\beta\in\Z$ forces $v^2\mid2$, and hence $v=1$. For integral $X$, $\beta$ is odd.

In the scalar norm block and the physical anchor, replace $S$ by
\[
 R=S-P=[a-1]P.
\]
The comparison and product-height blocks still use $S$. The retained subtraction guard excludes $a=\pm1$, which is harmless for completeness. The original scalar contact product is
\[
 \beta(\beta-1)(\beta-2401)(32\beta+343)(2\beta-1).
\]
It never vanishes for rational $X$: $\beta\ge3$, and $\beta=2401$ would give $X^2=1199$, which is not a rational square. At an odd prime, a negative valuation of $\beta$ is twice that of $X$, as required by the scalar argument.

\subsubsection{Soundness}\label{subsec:v9-soundness}
Put $c=|a-1|$. Applying the scalar and anchor lemmas to $R$ gives
\[
 \operatorname{h}(\beta)\le\mathcal H c^{19/32},\qquad
 \log_2c>25{,}165{,}823.5,
 \qquad 1\le\mathcal H<2^{2{,}093{,}966}.
\]
The anchor is still $S_0=3^{2^{26}}$. These explicit bounds imply
\[
 c>(80\mathcal H)^{10},\qquad
 |a|\ge c-1\ge c/2>(40\mathcal H)^{10},
\]
and
\begin{equation}\label{v9:eq:margins}
 \operatorname{h}(\beta)<c/80\le |a|/40,
 \qquad\log(40\mathcal H)<c/80\le |a|/40.
\end{equation}
For example $10(7+2{,}093{,}966)<25{,}165{,}823.5$, and $80<2^7$. The first estimate in~\eqref{v9:eq:margins} follows from $c^{13/32}>80\mathcal H$; the logarithmic estimate is weaker still.

The local comparison lemmas retain their original conclusions. Let
\[
 F=t^2-a(a-2)tb+(a-2)^2b^2-1,
 \qquad B=\max(|n|,|a|,|t|,|b|).
\]
For $B>4$, the triangle inequality gives $|F|\le2B^4+6B^3+5B^2+1<5B^4$. If $F\ne0$, the three radicals divide it to order 40, so~\eqref{v9:eq:height-transfer} gives
\[
 B^{1/10}<7\mathcal H\,5^{19/40}+6<24\mathcal H.
\]
This contradicts the anchor. Thus $F=0$. Since the point $Z$ is finite and nonidentity, $b\ne0$. Lemma~\ref{lem:o8-lucas} now gives $\log B>|a|-3$.

If $n\ne\beta$, the nonzero numerator $N$ of $n-\beta$, in reduced form, satisfies $|N|\le2e^{\operatorname{h}(\beta)}B$, and the radicals divide it to order 12. Equation~\eqref{v9:eq:height-transfer} and $B>4$ give
\[
 B^{5/12}<40\mathcal H e^{(19/12)\operatorname{h}(\beta)},
\]
so
\[
 \log B<\frac{12}{5}\log(40\mathcal H)+\frac{19}{5}\operatorname{h}(\beta)
 <\frac{31}{200}|a|<|a|-3.
\]
This is the other contradiction. Therefore $n=\beta\in\Z$, and $X\in\Z$.

\subsubsection{Completeness}\label{subsec:v9-completeness}
Fix $X\in\Z$, and set $n=\beta$. Perform the old scalar preparation: choose its finite correction and then a positive shared index $h_0$ divisible by 64 and sufficiently deep at the fixed small primes. Choose a positive scalar index $s$ sufficiently divisible to satisfy the scalar leaves and physical anchor for $[s]P$. We may require $6\mid s$ and take $s$ arbitrarily large. Set $a=s+1$. Then $a$ is odd, $a\equiv1\pmod3$, and $S-P=[s]P$ is exactly the prepared scalar point.

Define
\[
 U_0=0,\qquad U_1=1,\qquad U_{j+1}=aU_j-U_{j-1},
 \qquad k=2(a-2),
\]
and take $b=U_k/(a-2)$ and $t=U_{k+1}$. Modulo $a-2$, induction gives $U_j\equiv j$, so $b\in\Z$. Modulo two the sequence has period $0,1,1$; since $k\equiv1\pmod3$, both $U_k$ and $U_{k+1}$ are odd. The positive odd denominator $a-2$ shows that $b$ is odd. The Lucas identity gives
\[
 t^2-a(a-2)tb+(a-2)^2b^2=1.
\]
Thus all four root indices $n,a,t,b$ are positive and odd. Proposition~\ref{v9:prop:odd-chart} supplies their chart witnesses.

Fix these roots and the three scales of Corollary~\ref{v9:cor:scales}. Local continuity of the scalar and comparison functions at the prepared shared point gives the same fixed modulus $L$, with $64\mid L$. Only then choose the contact primes by the inherited simultaneous-prime theorem. The direct-target construction chooses a common index $h$ satisfying the primary and opposite reduction-order congruences and $h\equiv h_0\pmod L$, with the required real bounds. Its comparison points are $H=[h]P$ and $U_i=[hn_i]P$. Their shifted points have indices $h-1$ and $(h-1)n_i$.

All quotient-domain arguments remain valid. The product scales are fixed positive squares, so the original pole witnesses apply unchanged. The scalar block is preserved by the congruence modulo $L$. This gives a rational solution for every integral $X$ and completes the arithmetic proof of the primary construction.

\subsection{Replacing five logarithms by asymmetric formal depth}\label{v9:sec:depth}
This section is used only for the 523-variable construction. Let
\[
 U_0=H=[h]P,\quad U_i=[u_i]P,\quad
 V_0=H-P,\quad V_i=U_i-Q_i,
\]
where $Q_i=[n_i]P$ are the four roots. Put $f=1/x$ and $z=-x/y$. Keep the alternate truncated logarithms $\ell_{40}(V_i)=z(V_i)G_9(Df(V_i)^2)$, but replace every primary truncated logarithm by $z(U_i)$. Thus
\[
 L_i=z(U_i)/z(H),\qquad
 L_i'=\overline{\ell_{40}(V_i)/\ell_{40}(V_0)}.
\]
The bar denotes conjugation in $K$. The sixteen tests now have the following orders:
\[
\begin{array}{c|l}
1 & f(U_1),\ldots,f(U_4),\ \overline{f(V_0)},\ldots,\overline{f(V_4)}\\[2pt]
40 & f(H),\ L_1-L_1',\ldots,L_4-L_4',\ L_1-\beta,\
       \text{the direct Lucas residual}.
\end{array}
\]
The scalar tests retain their original order 32. The stronger beta tests imply all former order-12 conclusions.

\subsubsection{Local soundness}\label{subsec:v9-local-depth}
At a consumed unramified place above $p>76176$, a formal point $Q$ with $r=v(z(Q))\ge1$ satisfies
\begin{equation}\label{v9:eq:zerror}
 \log_E(Q)/z(Q)\in1+p^{4r}\mathcal O.
\end{equation}
This is the same formal-logarithm expansion used in the source, starting its tail at the first nonconstant term. Namely $v(f(Q))=2r$, and for its coefficients $a_j$ the factorial estimate is $v(a_j)\ge-(4j+1)/(p-1)$. If $4j+1<p$, the coefficient is integral. For the remaining $j\ge19044$, the displayed bound still makes the $j$th term have valuation at least $4r$. Both $z(Q)$ and its full logarithm have valuation $r$.

The new test on $f(H)$ gives $v(z(H))\ge20$. The order-one tests put the other points on the formal branch. The full logarithm ratios at the primary and conjugate places are
\[
 r_i=u_i/h,\qquad r_i'=(u_i-n_i)/(h-1).
\]
At least one of $h,h-1$ is a $p$-unit, so at least one of these ratios is integral. Even before knowing the depth of $U_i$, the weak unit estimate~\eqref{v9:eq:zerror} and the alternate truncation estimate show that $L_i,L_i'$ have the same valuations as $r_i,r_i'$. The order-40 cross test transfers integrality to both sides. Now $L_i$ integral implies $v(z(U_i))\ge v(z(H))\ge20$. The primary logarithmic corrections are therefore $1$ modulo $p^{80}$. The alternate corrections are $1$ modulo $p^{40}$. Since the ratios are integral, these relative estimates become absolute:
\[
 L_i\equiv r_i\equiv r_i'\pmod {p^{40}}.
\]
Multiplying by $h$ and $h-1$ and subtracting gives $L_i\equiv n_i\pmod {p^{40}}$. At an inert prime conjugation preserves the local valuation, so the same proof applies. The height contradictions of Section~\ref{v9:sec:labels} remain unchanged.

\subsubsection{A stronger choice within the existing prime-pattern argument}
For completeness, the common index must acquire extra primary divisibility. The following observation ensures compatibility with the opposite reduction orders.

\begin{lemma}\label{v9:lem:separation}
The contact primes in the source completeness argument may be chosen so that no selected prime divides any opposite reduced group order.
\end{lemma}
\begin{proof}
After merging repeated contact forms and fixing all scale and congruence data, the prime-pattern proof uses distinct forms $p_j=u+c_jw^2$, with fixed coefficients $c_j$. Integer dilations and ray congruences change only fixed constants. It counts $u\in[N/2,N]$ and $1\le w\le T$, where $T=\lfloor\tau(N)\sqrt N\rfloor$ and $\tau(N)\to0$ can be chosen arbitrarily slowly in the final diagonal argument. Its weighted count, after the prescribed obstruction removal, retains a fixed positive proportion of $NT$. These are the fixed-field weighted count in Proposition~\ref{prop:v5-artin-patterns} and the final diagonal argument of Section~\ref{subsec:v5-diagonal}, including its obstruction removal.

If there are $k$ forms, impose also
\[
 \frac{\tau(N)N^{1/6}}{(\log N)^k}\longrightarrow\infty.
\]
This is compatible with the same diagonal choice: for each fixed positive value of $\tau$, its threshold can be added to the existing finite list of thresholds. The strip $w<N^{1/3}$ has at most $O(N^{4/3})$ integer pairs and weight at most $O((\log N)^k)$ per pair. Its total weight is $o(NT)$, so an admissible pattern remains outside it.

For these patterns $w^2\ge N^{2/3}$ and $w^2=o(N)$. All $p_j$ are asymptotic to the same $u$. Hasse's bound \cite[Theorem V.1.1, p. 138]{silverman} gives $d_+(p_j)=p_j+O(\sqrt N)$ for the opposite reduced group order. Every ratio $d_+(p_j)/p_i$ is eventually strictly between $1/2$ and $2$. Hence $p_i\mid d_+(p_j)$ would force equality. If $i\ne j$, this contradicts
\[
 |p_i-p_j|\ge c_*N^{2/3},\qquad
 c_*:=\min_{i\ne j}|c_i-c_j|>0.
\]
If $i=j$, equality is impossible because $4\mid d_+(p_i)$ whereas $p_i$ is odd.
\end{proof}

\subsubsection{Completeness with deeper primary points}\label{subsec:v9-deep-completeness}
Keep the choice order in Section~\ref{v9:sec:labels}: the roots, scales and continuity modulus $L$ precede the prime selection. Choose primes using Lemma~\ref{v9:lem:separation}, beyond all prime factors of $L$ and the fixed scale data. Put
\[
 M_\pm=\mathop{\mathrm{lcm}}_p(d_\pm(p)/4),\qquad \Pi=\prod_p p.
\]
The direct-target prime theorem gives
\[
 \gcd(M_-,M_+)=1,\quad
 \gcd(M_-,L)=16,\quad\gcd(M_+,L)=1.
\]
Lemma~\ref{v9:lem:separation} gives $\gcd(\Pi,M_+)=1$, and the cutoff gives $\gcd(\Pi,L)=1$. Thus
\[
 h\equiv0\pmod{\Pi^{19}M_-},\qquad
 h\equiv1\pmod{M_+},\qquad h\equiv h_0\pmod L
\]
are compatible, since $64\mid h_0$. The inherited real-progression argument supplies a sufficiently large $h$ in this progression with the same real bounds.

The point $P$ is globally a double, so multiplication by the primary doubling-image order first makes it formal, with $z$-valuation at least one. Multiplication by $p^{19}$ raises this valuation by nineteen, using the formal logarithm isometry of \cite[Theorem IV.6.4, p. 132]{silverman}; these completions are unramified, $v(p)=1$, and $1>1/(p-1)$. Consequently $H=[h]P$ and $U_i=[hn_i]P$ have primary depth at least twenty, and their parameter ratios equal $n_i$ to order eighty. At the opposite places, $V_0=[h-1]P$ and $V_i=[(h-1)n_i]P$ are formal and their retained logarithm ratios equal $n_i$ to order forty.

The pole-completeness estimates still apply to these fixed test orders. In the source notation their large-real-parameter decay is $\lambda^{-3}$ for each order $m$, and the small-prime lower bound is $3\nu+(10m-8)\eta-2c$, which still tends to infinity with the preparatory divisibility. Scalar conditions and the anchor are preserved by $h\equiv h_0\pmod L$. This proves completeness for the optional refinement.

\subsection{Polynomial construction and exact verification}\label{subsec:v9-finite}
The generators apply three changes to the earlier point-chart source: the four odd-coset charts, the product and three scales, and the shifted scalar point with $\beta=2X^2+3$. The depth variant first changes the logarithms and test orders as in Section~\ref{v9:sec:depth}. Each field equation is represented by two rational component equations. Intermediate products are assigned coordinates as needed to keep every residual quadratic.

The full-log source has 577 raw coordinates; the pinned upstream compiler reduces this to 561. Seven further certified eliminations give 554. The depth source has 540 raw coordinates; compilation gives 532 and nine certified eliminations give 523. Every elimination is justified by exact polynomial identities and reversible substitutions preserving the input. Clearing residual denominators and taking the sum of their squares gives the displayed integer quartics. Over $\Q$, the sum vanishes exactly when every residual does.

The constant-chain versions add $c_0,\ldots,c_{16}$ with
\[
 c_0=3^{1024},\qquad c_{i+1}=c_i^2\quad(0\le i<16).
\]
These witnesses have unique rational values and $c_{16}=S_0$. Substitution for $S_0$ leaves every residual quadratic. Direct expansion gives the 3,247-bit maximum coefficient bound.

The accompanying archive contains the complete polynomial tables, residual systems, coordinate maps, elimination certificates, pinned source generators, arithmetic notes and independent checkers. Each final package has a standard-library-only \texttt{verify\_all.py}. It replays all eliminations, verifies the unique constant lift, independently expands the sums of squares, and compares every coefficient. Both full packages have passed fresh independent runs. A separate source replay rebuilds the two compiled precursors and checks the match to their frozen residual systems.

These finite checks certify the counts, degrees and stated algebraic transformations. The arithmetic assertions in Sections~\ref{v9:sec:height}--\ref{v9:sec:depth} received separate adversarial AI review; they have not received independent human peer review or a complete formal proof check. The inherited arithmetic inputs are the results of this manuscript identified in Section~\ref{subsec:v9-inputs}; no theorem from the contribution archive is assumed as an additional arithmetic input.



\begin{proof}[Proof of Theorem~\ref{thm:v9-quartics}]
The chart identities and Lemma~\ref{lem:v9-bridge} identify the exact
source used by the generators with the arithmetic system described here.
Proposition~\ref{v9:prop:odd-chart} enforces the odd root indices, so
Corollary~\ref{v9:cor:scales} supplies the required height bound.
Section~\ref{subsec:v9-local-depth} verifies the local implications for
the final test orders. Section~\ref{subsec:v9-soundness} then forces
the input to be an integer. Conversely, for an integer input,
Sections~\ref{subsec:v9-completeness} and
\ref{subsec:v9-deep-completeness}, using Lemma~\ref{v9:lem:separation},
construct witnesses for every source equation and retained guard.
The certified substitutions preserve this projection in both directions.
The sum of squares vanishes over $\Q$ precisely when all its residuals
vanish. Finally the 17-coordinate constant chain has exactly one rational
solution and ends at $S_0$, so it preserves the same projection. This
proves the statements for both polynomials.
\end{proof}

The complete tables are in the files
\begin{center}\small
\nolinkurl{quartics/quartic-523-symbolic-coefficients.json}\\
\nolinkurl{quartics/quartic-540-fully-expanded.txt}.
\end{center}
Their SHA-256 digests, respectively, are
\begin{center}\small
\nolinkurl{6a3685980de4512bdc7f3a781d0c0d4b32cce5829ec83f03f318264931bdbbe3}\\
\nolinkurl{83e7a31e11f9ad2edc2f67935e50bff437d87ef4b77d37b6a4391e8117770680}.
\end{center}
The reproduction command is
\nolinkurl{python scripts/check_odd_product.py}; the complete source,
certificates and repaired, non-mutating verifiers are in
\nolinkurl{quartics/quartic-523-540/}.

% END VERSION 9 ADDITION

\section{A quartic in 498 variables}\label{sec:v12}

We now reduce the 523-variable construction of
Theorem~\ref{thm:v9-quartics} to 498 variables.  The input remains one
rational variable $X$; every other variable remains existentially
quantified over $\Q$.  The changes below use the same normalized curve,
the same integer labels, the same three height detectors, and the same
published and manuscript inputs as Section~\ref{sec:v9}.  In particular,
they do not replace the restricted converse, the exponent-nineteen
height theorem, or the simultaneous-prime argument by new assumptions.

\begin{theorem}\label{thm:v12-quartic}
There is a specified integer polynomial $P_{498}$ of total degree four
in 498 variables such that
\[
 X\in\Z\quad\Longleftrightarrow\quad
 \exists\mathbf v\in\Q^{497}:P_{498}(X,\mathbf v)=0.
\]
The polynomial has 24,669 nonzero monomials, content one, and largest
absolute coefficient $16\cdot3^{2^{26}}$.  Its symbolic coefficient
table uses the fixed integer $C=3^{2^{25}}$, not an additional variable.
There is also a specified integer quartic $P_{514}$ with
\[
 X\in\Z\quad\Longleftrightarrow\quad
 \exists\mathbf w\in\Q^{513}:P_{514}(X,\mathbf w)=0.
\]
It has 24,716 nonzero monomials, content one, and largest absolute
coefficient $3^{2048}+1005072314676622<2^{3247}$.
\end{theorem}

The proof distinguishes two sorts of changes.  Some are rational
changes of private witnesses with a rational inverse.  Others alter
the arithmetic certificate while preserving both soundness and the
existence of witnesses for every integer input.  For the latter we
prove those two assertions separately; we do not claim equality of
all intermediate witness assignments.

\subsection{Conventions and retained inputs}\label{subsec:v12-inputs}

Throughout put
\[
 K=\Q(\rho),\quad \rho^2=2,\quad
 d_*=3-\rho,\quad D=d_*^2=11-6\rho,\quad
 E:y^2=x^3-Dx.
\]
The point $P=P_*$ is the same normalized generator as in
Section~\ref{subsec:v5-point-source} and Section~\ref{sec:v9}:
\[
 P=(c_*,a_*b_*),\qquad
 a_*=-1+\tfrac74\rho,\quad
 b_*=\tfrac{39}{8}-\tfrac92\rho,\quad
 c_*=a_*^2=\tfrac{57}{8}-\tfrac72\rho.
\]
Thus $b_*^2=c_*^2-D$.  In particular $a_*,b_*,c_*,D$ are nonzero.
We use the inherited equality $2E(K)=\Z P$, the nonzero point domains
of Lemma~\ref{lem:v5-point-domains}, and the inherited estimate
$\operatorname h(x([n]P))\ge n^2-6$.  Here $\operatorname h$ is the
absolute logarithmic Weil height.  At every unramified place above a
rational prime $p$ we normalize $v(p)=1$.  The coefficient basis is
always $(1,\rho)$ and conjugation sends $\rho$ to $-\rho$.

Retain the four odd-coset root charts of
Proposition~\ref{v9:prop:odd-chart}, with labels
\[
 Q_1=[n]P,\quad Q_2=S=[a]P,\quad
 Q_3=[t]P,\quad Q_4=[b]P,
 \qquad n,a,t,b\text{ odd}.
\]
The letter $b$ in this display is an integer point label.  We use
$\mathfrak b$ below for a rational input to a pole predicate.
Put $Q_0=P$, $U_0=H=[h]P$, $U_i=[u_i]P$, and
$V_i=U_i-Q_i$.  The scalar point is
\[
 R=S-P=[s]P,\qquad s=a-1,
\]
and the input label is $\beta=2X^2+3$.

We retain the pole obstruction and pole-completeness statements
identified in Section~\ref{subsec:v9-inputs}, specifically
Lemmas~\ref{hyp:o8-pole}, \ref{lem:original-seven-simultaneous-realization},
and \ref{lem:v5-one-place}, including their existing source dependencies.
We also retain the three scales of Corollary~\ref{v9:cor:scales} and
its product-height estimate.  Write $\mathcal H=H_{19}$, so that
\begin{equation}\label{eq:v12-height-input}
 1\le\mathcal H<2^{2093966},\qquad
 B^2-6<7\mathcal H|N|^{19/k}
\end{equation}
whenever the three contact radicals divide a nonzero integer $N$
to order $k$, and $B=\max(|n|,|a|,|t|,|b|)$.
This is exactly \eqref{v9:eq:height-transfer}; none of its scales is
replaced by a field norm in the present construction.

The local formal-group inputs are the integral formal group in
$z=-x/y$, its multiplication series, and its logarithm.  The integral
construction is the one in \cite[IV\S1, Example~IV.2.2.3 and
Proposition~IV.2.3(a)]{silverman}.  At the unramified places used below,
$p>76176$, and the logarithm is an isometry on depth at least one by
\cite[Theorem IV.6.4, p.~132]{silverman}; its hypothesis is
$1>1/(p-1)$.  The coefficient and truncation estimates are those of
Lemma~\ref{lem:quartic-reciprocal-log} and its unramified extension in
Lemma~\ref{lem:v5-unramified-log}.  We spell out their new applications
below.  No minimal-model assertion at the prime above two is needed.

\subsection{Mixed logarithmic precision}\label{subsec:v12-mixed}

Let $f=1/x$ and $z=-x/y$, as before.  Define
\[
 a_0=1,\quad a_1=-\tfrac25,\quad
 a_j=\frac{4j-5}{4j+1}a_{j-1}\quad(j\ge2),\qquad
 G_J(T)=\sum_{j=0}^J a_jT^j.
\]
Thus the source truncations are
$\ell_{12}=zG_2(Df^2)$ and $\ell_{40}=zG_9(Df^2)$.
Keep every primary logarithm equal to the uncorrected parameter $z$.
Only the alternate numerator for $V_1$ is shortened:
\begin{equation}\label{eq:v12-ratios}
 L_i=\frac{z(U_i)}{z(H)},\qquad
 L'_1=\overline{\frac{\ell_{12}(V_1)}{\ell_{40}(V_0)}},\qquad
 L'_i=\overline{\frac{\ell_{40}(V_i)}{\ell_{40}(V_0)}}\quad(2\le i\le4).
\end{equation}
The only denominator logarithm remains $\ell_{40}(V_0)$, with its
original nonzero domain.  In particular this change introduces no
new logarithm denominator.

The sixteen field tests and their orders are now
\begin{equation}\label{eq:v12-tests}
\begin{array}{c|l}
1&f(U_1),\ldots,f(U_4),\ \overline{f(V_0)},\ldots,\overline{f(V_4)}\\[2pt]
12&L_1-L'_1,\ L_1-\beta\\[2pt]
40&f(H),\ L_2-L'_2,\ L_3-L'_3,\ L_4-L'_4,\ F(L_2,L_3,L_4),
\end{array}
\end{equation}
where
\[
 F(A,T,B)=T^2-A(A-2)TB+(A-2)^2B^2-1.
\]
This changes the two order-forty input tests in the 523-variable
source to order twelve.  The three other label comparisons and the
Lucas test remain of order forty.

\begin{lemma}\label{lem:v12-mixed-local}
At every consumed unramified place above $p>76176$, the tests
\eqref{eq:v12-tests} imply
\[
 L_1\equiv n\pmod {p^{12}},\qquad
 L_i\equiv n_i\pmod {p^{40}}\quad(2\le i\le4),
 \qquad(n_2,n_3,n_4)=(a,t,b).
\]
The primary points $U_i$ have depth at least twenty after integrality
of the ratios has been established.
\end{lemma}
\begin{proof}
The order-one tests put the indicated points in the formal group.
The test on $f(H)$ gives $v(z(H))\ge20$.  For a formal point $Q$ of
depth $r\ge1$, the unchanged logarithm expansion gives
\[
 \log_E(Q)/z(Q)\in1+p^{4r}\mathcal O.
\]
Its truncation after $G_2$ has relative error of order at least twelve,
and its truncation after $G_9$ has relative error of order at least
forty.  To check these particular applications, the first omitted
indices are respectively $j=3$ and $j=10$ in the source recurrence.
For $4j+1<p$ the coefficients are integral.  For the remaining indices
the source factorial bound
$v_p(a_j)\ge-(4j+1)/(p-1)$ is dominated by $4jr$ and gives the same
stated lower bounds.  Both truncating factors are units.

The full ratios are the rational numbers
\[
 r_i=\frac{u_i}{h},\qquad r'_i=\frac{u_i-n_i}{h-1}.
\]
At least one of $h,h-1$ is a $p$-unit, so at least one of these two
ratios is integral.  Before any improved depth is inferred for $U_i$,
the unit factors just described show
$v(L_i)=v(r_i)$ and $v(L'_i)=v(r'_i)$.  Each cross-test has positive
order, so integrality transfers to both approximate ratios and both
full ratios.  This first step is essential.

Now $L_i$ integral gives $v(z(U_i))\ge v(z(H))\ge20$.
The primary relative errors have order at least eighty.  The alternate
relative errors have order twelve for $i=1$ and forty for $i\ge2$:
in the first case the denominator still has order forty, so it does
not reduce the numerator's order twelve.  Since the ratios are
integral, these relative errors give absolute congruences.  Together
with the cross-tests they yield $L_i\equiv r_i\equiv r'_i$ to the
asserted order.  Multiply by $h$ and $h-1$ and subtract.  The result
is $L_i\equiv n_i$, without division by either integer.  At an inert
prime conjugation preserves the single local valuation, so this
argument has the same conclusion there.
\end{proof}

For clarity, the detector's product is unchanged as a mathematical
operation.  For a test $D_j$ of order $m_j$, a scale $\mathcal A$, and
$t=\mathcal A\lambda$, $\delta=t-1$, its factors are
\[
 1+\mathcal R\left(\frac{D_j}{\delta^{m_j+1}\mathcal R^{m_j}}\right)^2,
 \qquad
 \mathcal R=\mathcal A^{-5}F_5(t)/\delta^6,
 \qquad F_5(t)=\prod_{\theta\in\Theta}(t-\theta),
\]
with the same $\Theta$ as \eqref{eq:v5-detector-data}.  On writing
$e=(\mathcal A\delta)^5/F_5(t)$, the multiplier of $D_j^2$ is
$q_{m_j}=q_1e^{2m_j-2}$, where $q_1=e/\delta^3$.
For each $m\in\{1,12,40\}$, form the coefficients of
\[
 \prod_{j:m_j=m}(1+TD_j^2)
\]
and evaluate at $T=q_m$.  Their degrees are respectively $9,2,5$.
Their product is exactly the same sixteen-factor detector expression.
The smaller $G_2$ numerator and these three grouped products explain
the changed quadratic circuit.  Multiplying in $K$ means the two
rational bilinear identities in the fixed basis.  Every reassociation
or grouping here follows from distributivity, with no division or
new witness domain.  The retained coefficient tables and exact
quadratic elimination identities specify the resulting circuit.

\subsection{Five rational half-trace guards}\label{subsec:v12-guards}

For each $0\le i\le4$ write the subtraction denominator as
\[
 d_i=x(U_i)-x(Q_i)=A_i+B_i\rho.
\]
The old field inverse used two private rational coordinates.  Replace
those two equations by the single equation
\begin{equation}\label{eq:v12-half-trace}
 A_i\gamma_i=1,\qquad\gamma_i\in\Q.
\end{equation}
No other equation used the removed inverse coordinates.  This saves
one coordinate in each of five blocks.  It strengthens the domain;
there is no claim that every old trace-zero assignment is retained.

Soundness has an explicit rational map.  From \eqref{eq:v12-half-trace}
we have $A_i\ne0$, hence $d_i\ne0$ and
$A_i^2-2B_i^2\ne0$.  Restore the old field inverse as
\[
 d_i^{-1}=\frac{A_i-B_i\rho}{A_i^2-2B_i^2}.
\]
All old subtraction outputs and all their consumers are unchanged.

For completeness fix the four root points and the scalar preparation
before choosing contact primes, as in Section~\ref{subsec:v9-completeness}.
Increase the allowed finite prime cutoff to exclude every prime at
which any of $P,Q_1,\ldots,Q_4$ has nonintegral abscissa.  This is a
finite restriction on already fixed points.  Choose any one of the
split contact primes supplied by Lemma~\ref{lem:v5-two-reductions},
and denote its primary and opposite places by $v,\bar v$.
The unchanged CRT construction makes every $U_i$ formal at $v$ and
makes it reduce to $Q_i$ at $\bar v$.  All five $Q_i$ have finite
integral reduction.  Consequently
\[
 v(d_i)<0,\qquad \bar v(d_i)\ge0.
\]
Conjugation exchanges the two places, so
$v(\overline{d_i})=\bar v(d_i)\ge0$.  The two terms of
$d_i+\overline{d_i}=2A_i$ have unequal valuations; the sum is nonzero.
Thus all five inverses $\gamma_i=1/A_i$ exist for the same completeness
witnesses.  This includes $i=0$.  It imposes no unit condition on the
root labels and no condition on the permitted zero pole witnesses.

\subsection{Normalizing the three detector poles}\label{subsec:v12-poles}

Here $\mathfrak b$ is one detector's rational pole input.  In the
compiler's rational coordinates the seven relevant equations are
\begin{align*}
 p^2&=p_2,&17\mathfrak b p_2+3e^2-34e\mathfrak b-289\mathfrak b^2&=0,\\
 u_2&=\mathfrak b^2,&17e\mathfrak b-eu_2-17g&=0,\\
 s_0^2&=s_2,&q_0^2&=q_2,\\
 &&17\mathfrak b q_2+867g^2-34gs_2-s_2^2&=0.
\end{align*}
In particular the stored $g$ is the source's quantity $Bh/17$,
not $Bh$.  The following substitution respects that factor:
\[
 e=\mathfrak b E_0,\quad p=\mathfrak b P_0,\quad
 s_0=\mathfrak b T_0,\quad q_0=\mathfrak b^2Q_0,
\]
\[
 p_2=\mathfrak b^2P_2,\quad s_2=\mathfrak b^2T_2,\quad
 q_2=\mathfrak b^4Q_2,\quad
 u_2=\mathfrak b^2,\quad g=\mathfrak b^2G_0.
\]
The replacement consists of six quadratic equations:
\begin{align}\label{eq:v12-pole-normalized}
 P_0^2&=P_2,&17\mathfrak b P_2+3E_0^2-34E_0-289&=0,\notag\\
 17E_0-E_0\mathfrak b-17G_0&=0,&T_0^2&=T_2,\\
 Q_0^2&=Q_2,&17\mathfrak b Q_2+867G_0^2-34G_0T_2-T_2^2&=0.\notag
\end{align}

The division by $\mathfrak b$ in the reverse coordinate change is
legitimate in both systems.  Indeed the unchanged detector equations,
with its actual product $W=w_0+w_1\rho$, include
\[
 w_0\iota=1,\qquad
 2w_0\nu=w_0^2-2w_1^2,
\]
and an equation whose specialization at $\mathfrak b=0$ is $\nu^2=0$.
Thus $\mathfrak b=0$ would give $\nu=0$ and $w_0^2=2w_1^2$.
Rationality then forces $w_0=w_1=0$, contradicting $w_0\iota=1$.

Under the displayed substitution the old $u_2$ equation is an identity.
The old $p_2$, first-conic, $g$, and $s_2$ equations are
$\mathfrak b^2$ times their replacements; the old $q_2$ and second-conic
equations are $\mathfrak b^4$ times their replacements.  These identities
prove mutually inverse rational maps fixing every outside coordinate.
They remain valid when any of $e,p,s_0,q_0,g$ is zero.  Applying the
change to three detectors saves three coordinates.  The scalar pole
with its preparation variable is not altered by this normalization.

\subsection{One coefficient in the physical anchor}\label{subsec:v12-anchor}

Put $C=3^{2^{25}}$ and write the actual scalar value as
$f(R)=f_0+f_1\rho$.  The previous fixed parameter was $S_0=C^2$.
The old three-witness anchor encoded
$f_0/C,f_1/C\in\Z_{(2)}\cap\Z_{(3)}$.  On the present integer orbit it
is enough to impose that condition on $f_1/C$.

\begin{lemma}\label{lem:v12-anchor-coefficient}
For every nonzero even integer $s$, write
$f([s]P)=f_0+f_1\rho$.  Both coefficients are $2$-integral, and
\[
 v_3(f_1)=2+2v_3(s),\qquad
 v_3(f_0)\ge3+2v_3(s).
\]
In particular $f_1\ne0$, and
\[
 f_1/C\in\Z_{(2)}\cap\Z_{(3)}
 \quad\Longleftrightarrow\quad
 f_0/C,f_1/C\in\Z_{(2)}\cap\Z_{(3)}.
\]
\end{lemma}
\begin{proof}
At $3$ the field is unramified quadratic with residue field
$\F_9=\F_3[\rho]/(\rho^2-2)$.  The exact starting calculation in
Section~\ref{subsec:o8-anchor-subgroup} is
\[
 z(2P)=
 \frac{3396356159088-4135204317756\rho}{3287811625151},
 \qquad z(2P)/3\equiv2+\rho\pmod3.
\]
The second congruence follows by reducing the displayed rational
coefficients after division by three.  Its square is $\rho$ in
$\F_9$.  Write $s=2\cdot3^k u$ with $3\nmid u$, allowing negative $u$.
The integral formal group has inverse $-z$, and hence odd
multiplication series.  Thus
\[
 [u]z=uz+O(z^3),\qquad [3]z=3z+O(z^3)
\]
with integral coefficients.  At depth $r\ge1$, division of the second
identity by $3^{r+1}$ leaves an error of valuation at least $2r-1>0$.
Multiplication by three therefore preserves the leading residue.
Multiplication by $u$ multiplies it by the nonzero residue of $u$.
Consequently
\[
 z([s]P)/3^{k+1}\equiv u(2+\rho)\pmod3.
\]
The exact curve relation is $f=z^2(1-Df^2)$.  On the formal branch its
second factor is one modulo three, and $u^2=1$ in $\F_3$.  Therefore
\[
 f([s]P)/3^{2k+2}\equiv\rho\pmod3.
\]
Independence of $1,\rho$ over $\F_3$ proves the two asserted
coefficient valuations, including the possible zero value of $f_0$.

At the place above two normalize $v(2)=1$.  The literal point has
$v(x(P))=-3$, $v(y(P))=-9/2$, and $v(z(P))=3/2$.  It lies in the
integral formal subgroup, which contains all its integral multiples.
One may construct that subgroup directly from
$w=z^3-Dzw^2$, where $w=-1/y$, as in
Section~\ref{subsec:o8-anchor-subgroup}; this uses no assertion that
the displayed equation is minimal at two.  Thus $f([s]P)$ is integral.
The equality
\[
 v(a+b\rho)=\min\{v_2(a),v_2(b)+\tfrac12\}\quad(a,b\in\Q)
\]
shows that both rational coefficients are $2$-integral.
Since $C$ is a $2$-unit, the equivalence now follows from the
coefficient valuations at three.
\end{proof}

The three-witness criterion in the proof of
Lemma~\ref{lem:v5-three-witness-anchor} applies to every rational $t$:
\begin{equation}\label{eq:v12-quaternion-criterion}
 t\in\Z_{(2)}\cap\Z_{(3)}
 \quad\Longleftrightarrow\quad
 \exists u,v,w\in\Q:\ u^2+6v^2-2w^2-3t^2+2=0.
\end{equation}
To recall its exact input, $(2,3)_\Q$ ramifies precisely at two and
three.  The embedding criterion
\cite[Proposition 14.6.7, pp.~223--224]{voight-quaternion} tests the nonsquareness of
$1-2/(3t^2)$ at those places when $t\ne0$.  At two, positive valuation
of $t$ gives odd negative valuation of the expression, valuation zero
gives residue three modulo eight, and negative valuation gives a
square in $1+8\Z_2$.  At three, nonnegative valuation gives odd negative
valuation of the expression, whereas negative valuation gives a square
in $1+3\Z_3$.  The expression is nonzero because $t^2=2/3$ has no
rational solution.  Multiplication by $3t^2$ gives
\eqref{eq:v12-quaternion-criterion}; for $t=0$ the witnesses are
$(0,0,1)$.  These are exactly the field, local conditions, and
exceptional case already present in that source proof.

Apply this criterion to $t=f_1/C$, multiply its witnesses by $C$,
and then replace the old $w$ by $w+C$.  The resulting anchor is
\begin{equation}\label{eq:v12-anchor}
 u^2+6v^2-2w^2-4Cw-3f_1^2=0.
\end{equation}
It uses the same three private rational witnesses and removes the
old coordinate for $f_0^2+f_0f_1-f_1^2$.  Since $a$ is odd and $R$ is
nonidentity, $s=a-1$ is a nonzero even integer.  By
Lemma~\ref{lem:v12-anchor-coefficient}, old and new anchors have the
same projection onto all their nonprivate coordinates.  In each
direction use \eqref{eq:v12-quaternion-criterion} to choose the private
witnesses; they need not be the same witnesses in both anchors.
The affine translation of $w$ is invertible.

The exact allowed subgroup is unchanged:
\begin{equation}\label{eq:v12-anchor-subgroup}
 s\in2\,3^{2^{24}-1}\Z\setminus\{0\},\qquad
 \log_2|s|\ge1+(2^{24}-1)\log_2 3>25165823.5.
\end{equation}
Indeed $v_3(f_1)\ge v_3(C)=2^{25}$ is precisely
$2+2v_3(s)\ge2^{25}$.  The subgroup permits arbitrarily divisible
positive $s$, as required by completeness.

\subsection{A shorter scalar norm tree}\label{subsec:v12-scalar}

Keep the exact contact polynomial and finite correction
\[
 r=\beta(\beta-1)(\beta-2401)(32\beta+343)(2\beta-1),\qquad
 c\in\left\{\pm\prod_{p\in I}p:
 I\subseteq\{2,3,5,7,11,13,17\}\right\}.
\]
Here $r\ne0$ for rational $X$.  The only positive contact not excluded
by $\beta\ge3$ is $\beta=2401$, which would imply $X^2=1199$,
impossible over $\Q$.  The old inverse equation for $r$ and its
power chain remain in place.

Write
\[
 f(H)=h_0+h_1\rho,\qquad z(R)/z(H)=z_0+z_1\rho,
 \qquad q=h_0^2-2h_1^2=\operatorname{Norm}_{K/\Q}(f(H)).
\]
Use $N_r(u,v)=u^2+rv^2$ and replace the old six-leaf scalar tree by
\begin{equation}\label{eq:v12-scalar-tree}
 W=N_r\left(N_r(1,q/r),\ N_r(z_0/r^{32},z_1/r^{32})\right).
\end{equation}
The scalar pole input is still exactly $\mathfrak b=crW^2$, with the
original preparation and nonzero condition.  This tree changes that
input's value, so completeness must supply new pole witnesses.

For an explicit quadratic encoding, let $i_{32}=r^{-32}$ be the
retained power-chain output and introduce eight rational coordinates
$q,A,U,V,V_2,B,B_2,W$ with equations
\begin{equation}\label{eq:v12-scalar-eight}
\begin{aligned}
 q&=h_0^2-2h_1^2,& r(A-1)&=q^2,\\
 U&=z_0i_{32},&V&=z_1i_{32},\\
 V_2&=V^2,&B&=U^2+rV_2,\\
 B_2&=B^2,&W&=A^2+rB_2.
\end{aligned}
\end{equation}
The existing $W^2$ coordinate is reattached to this $W$ with its
actual fixed rational normalization.  Thirteen old leaf and norm
coordinates disappear, giving a net saving of five at this stage.
The final elimination of $q$ is given below.

\begin{lemma}\label{lem:v12-scalar-sound}
The new scalar certificate implies
$M_5(\beta)^{32}\mid s$ and
$\operatorname h(\beta)\le\mathcal H|s|^{19/32}$.
\end{lemma}
\begin{proof}
Let $p>76176$ be a prime of the same five-contact radical as in
Lemma~\ref{v5:scalar32-bridge}, and put $\eta=v_p(r)$.
The unchanged contact calculation makes $\eta$ odd; the correction
$c$ is a $p$-unit.  The original pole obstruction, applied to the
retained nonzero input $crW^2$, gives
\[
 \eta+2v_p(W)>0.
\]
At a node $u^2+rv^2$, nonzero summands have opposite valuation parity,
so they cannot cancel.  Zero summands are harmless with valuation
$+\infty$.  If $\eta<0$, the constant-left path in
\eqref{eq:v12-scalar-tree} gives $v_p(A)\le0$ and $v_p(W)\le0$, a
contradiction.  Thus $\eta>0$ and $v_p(W)>-\eta/2$.
Descending the two levels gives
\[
 v_p(A)>-\eta/4,\quad v_p(B)>-3\eta/4,
\]
\[
 v_p(q/r)>-5\eta/8,\quad
 v_p(U)>-3\eta/8,\quad v_p(V)>-7\eta/8.
\]
Therefore
\begin{equation}\label{eq:v12-scalar-valuations}
 v_p(q)>3\eta/8>0,\qquad
 v_p(z_0)>(32-3/8)\eta,\qquad
 v_p(z_1)>(32-7/8)\eta.
\end{equation}
Since $\eta$ is a positive integer, rational integral valuations give
$v_p(z_0),v_p(z_1)\ge32$.

At a split prime $v_p(q)$ is the sum of the two valuations of $f(H)$;
at an inert prime it is twice the single valuation.  Positivity of
$v_p(q)$ therefore supplies at least one place $v$ above $p$ with
$v(f(H))>0$.  We claim only existence of one such place.  At that
place $H$ is formal and $v(z(H))\ge1$.  By the two coefficient bounds,
\[
 v(z(R))\ge v(z(H))+32>0.
\]
This need not make $R$ itself formal: it may lie near $(0,0)$.
We handle both branches explicitly.

Put $m=v(x(R))$.  The retained point domains give $x(R)y(R)\ne0$.
Since $D$ is a unit and
\[
 z(R)^2=\frac{x(R)}{x(R)^2-D},
\]
positive valuation of $z(R)$ excludes $m=0$.  For $m>0$ its valuation
is $m/2$, and for $m<0$ it is $-m/2$.  The exact duplication identity
\begin{equation}\label{eq:v12-duplication}
 x([2]R)=\frac{(x(R)^2+D)^2}{4x(R)(x(R)^2-D)}
\end{equation}
then gives, for either sign of $m$,
\[
 v(x([2]R))=-2v(z(R))<0,\qquad v(z([2]R))=v(z(R)).
\]
Its numerator and denominator are nonzero: with $m\ne0$ their
relevant terms have unequal valuations, and $p$ is odd.
Also $[2]R$ is nonidentity because $s\ne0$ and $P$ has infinite order.

Both $H$ and $[2]R$ are now in the formal subgroup where the logarithm
is a homomorphism and isometry.  Apply it to the exact global identity
\[
 [h]([2]R)=[2s]H.
\]
The points are nontorsion, so their nonzero formal logarithms permit
comparison of valuations.  Since $2$ is a $p$-unit,
\[
 v_p(h)+v(z([2]R))=v_p(s)+v(z(H)),
 \qquad v_p(s/h)\ge32.
\]
The nonzero integer $h$ has nonnegative $p$-valuation, hence
$v_p(s)\ge32$.  This proves the claimed divisibility at every prime
of $M_5(\beta)$.  The same exponent-nineteen height theorem as in
Lemma~\ref{lem:v9-bridge} gives the asserted height bound.
\end{proof}

For completeness use the old scalar preparation with its stronger
six-leaf conditions.  For integral $X$, $r$ is a nonzero integer.
Choose $c$, a sufficiently divisible positive $h_0$, and then a still
deeper positive $s$ satisfying the anchor, exactly as in the proof of
Lemma~\ref{v5:scalar32-bridge}.  At primes dividing $r$, both
coefficients of $f(H)$ are divisible by $r$ locally, so $q/r$ is
integral; the two retained ratio leaves are integral.  At primes
dividing $c$ but not $r$, the same leaves are divisible by the prime.
The constant-left tree is consequently one modulo every prime on
the support of $cr$.  It is nonzero and a unit at those primes.
Thus $crW^2$ is positive, has square reduced denominator, and has the
same local squareclasses as $cr$ at two, three, and seventeen.
The original pole-completeness lemma supplies fresh scalar pole
witnesses.  The new norm and ratio functions are continuous at the
chosen points, so the same finite-neighborhood argument provides a
congruence class of $h_0$ preserving all required conditions.

\subsection{A reciprocal chart for the scalar difference}
\label{subsec:v12-reciprocal}

This replacement is used only for $R=S-P$.  The five comparison
subtractions retain their strengthened guards
\eqref{eq:v12-half-trace}.  Write the incoming scalar point as
$S=(x,hj)$ with $x=h^2$ and $j^2=x^2-D$.  Put
\[
 d=c_*x+D,\quad N=a_*j-b_*h,\quad
 M=-2a_*Dh(x+c_*)-b_*j(c_*x-D).
\]
Both real embeddings give $c_*>0$, $D>0$, and $x\ge0$; hence $d\ne0$.
Introduce $g,k\in K$, retain the existing outputs $f,z\in K$, and
introduce a rational guard $\sigma$.  The final equations are
\begin{equation}\label{eq:v12-reciprocal}
 dg=N,\qquad
 a_*dk+2Dh+a_*b_*j-2Db_*g=0,\qquad
 f=g^2,\qquad f_0\sigma=1,\qquad kz=-g.
\end{equation}
Each field equation means its two rational component equations.
Every component equation is quadratic.  Since
$f_0=g_0^2+2g_1^2$, the guard is precisely $g\ne0$ over rational
coordinates, and imposes no unmentioned sign choice.

We give the identities and all exceptional cases needed to identify
these outputs with the actual group difference.  If
\[
 G=dg-N,\quad K_0=d^2k-M,\quad
 L=a_*dk+2Dh+a_*b_*j-2Db_*g,
\]
then expansion using $a_*^2=c_*$ and $b_*^2=c_*^2-D$ gives
\begin{equation}\label{eq:v12-dcancel}
 dL=a_*K_0-2Db_*G.
\end{equation}
Thus the first two equations of \eqref{eq:v12-reciprocal} are
exactly equivalent to $g=N/d$, $k=M/d^2$, because $a_*d\ne0$.
The latter expressions satisfy
\begin{equation}\label{eq:v12-reciprocal-identities}
 k^2+Dg^4=1,\qquad
 (a_*j-b_*h)(a_*j+b_*h)=(x-c_*)(c_*x+D).
\end{equation}
For example the first is the polynomial identity
$M^2+DN^4=d^4$ modulo $h^2=x$, $j^2=x^2-D$ and the two fixed-constant
identities.  This also specifies an exact denominator-cleared check.

The guard gives $g\ne0$.  If $k=0$, the first identity in
\eqref{eq:v12-reciprocal-identities} gives $d_*g^2=1$ or $-1$.
Taking rational norms gives $7\operatorname{Norm}(g)^2=1$, impossible
for a rational norm.  Thus $k\ne0$, and all displayed outputs are
well defined.  If $x\ne c_*$, use the ordinary subtraction slope
$m=(hj+a_*b_*)/(x-c_*)$.  Substitution and clearing denominators give
\[
 m^2-x-c_*=1/g^2,\qquad
 m(x-1/g^2)-hj=k/g^3.
\]
These are the abscissa and ordinate of $S-P$.  Therefore the retained
outputs are exactly $f(R)=g^2$ and $z(R)=-g/k$.

If $x=c_*$, the incoming chart forces
$h=\pm a_*$ and $j=\pm b_*$.  Equal choices of signs give $S=P$
and $g=0$, excluded by the guard.  The other choices give $S=-P$.
For $h=a_*,j=-b_*$ the explicit values are
\[
 g=\frac{-2a_*b_*}{c_*^2+D},\qquad
 k=\frac{c_*^4-6c_*^2D+D^2}{(c_*^2+D)^2}.
\]
They give the abscissa and ordinate of $-2P$ by the ordinary duplication
formula.  Simultaneously negating $h,j$ negates $g,k$ and leaves
$1/g^2$, $k/g^3$, and $-g/k$ unchanged.  Thus this branch also represents
the true difference, not an arbitrary value of an old zero-over-zero
slope.

In the complete present system that extra branch is impossible.
It would give $R=-2P$ and $v_3(f_1(R))=2$ by
Lemma~\ref{lem:v12-anchor-coefficient}, whereas
\eqref{eq:v12-anchor} requires $v_3(f_1(R))\ge2^{25}$.
It follows that every final solution has $x\ne c_*$.  Restore the old
inverse of $x-c_*$, slope, slope square, and ordinate products by their
ordinary rational formulas.  All retained outputs coincide.  Conversely
every old solution extends to \eqref{eq:v12-reciprocal} by the displayed
formulas, with $\sigma=1/f_0$; its actual nonzero $f(R)=g^2$ ensures
$f_0>0$.  This is an equivalence on all nonprivate coordinates of the
complete system, with the anchor explicitly used to exclude the
additional branch.

For the coordinate count one may first use $d_2=d^2$ and $d_2k=M$.
That intermediate reciprocal chart replaces eight old scalar
difference coordinates by the six components of $g,d_2,k$ and one
rational guard, saving one.  Identity \eqref{eq:v12-dcancel} then removes
the two components of $d_2$.  Forward, divide $dL=0$ by the already
nonzero $d$; backward, reinsert the unique $d_2=d^2$ and use
$a_*K_0=0$.  This second step is a direct rational equivalence requiring
no new exceptional-point argument.  The combined saving is three.

\subsection{A shorter weight circuit with both outputs retained}
\label{subsec:v12-toric}

In each detector let $e$ be its weight $e$, $q=q_1$, and keep both
outputs $q_{12},q_{40}$.  The grouped evaluations already contain
$Q=q^3$ and $T=q_{12}^2$.  Replace the six old intermediate powers of
$e$ by five rational coordinates $\alpha,\gamma,u,v,\tau$ and impose
\begin{equation}\label{eq:v12-toric}
 e\alpha=1,\quad e\gamma=\alpha^2,\quad
 \gamma u=eq,\quad \gamma v=u^2,\quad
 Qq_{12}=v^2,\quad q\tau=eT,\quad vq_{40}=T\tau.
\end{equation}
These are seven quadratic equations for seven coordinates when the
two retained outputs are counted.  The old circuit had eight.

The first equation gives $e\ne0$.  The unchanged definition of $q$
is $q\delta^3=e$, with the powers of $\delta$ represented by their
retained coordinates, so $q\ne0$ as well.  The equations now determine,
in order,
\[
 \alpha=e^{-1},\quad\gamma=e^{-3},\quad
 u=qe^4,\quad v=q^2e^{11},\quad q_{12}=qe^{22}.
\]
The division by $Q=q^3$ is legitimate.  The retained equation for
$T$ then gives $T=q^2e^{44}$, and the last two equations give
\[
 \tau=qe^{45},\qquad q_{40}=qe^{78}.
\]
In particular $v\ne0$ whenever it is used as a divisor.  This is a
successive computation, not a circular specification involving $T$.

The forward map from the old circuit assigns these five Laurent
monomials and retains $q_{12},q_{40}$, $Q,T$, and every other coordinate.
For the reverse map restore the old six powers as
$e^{-1},e^{-3},e^{-7},e^{-14},e^{-29},e^{-55}$.  The old output equations
then give the same $qe^{22}$ and $qe^{78}$.  The replaced internal
powers have no outside consumers.  Thus these maps preserve every
outside coordinate, not merely the input.  Applying this equivalence
to all three detectors saves three coordinates.

\subsection{Removing the scalar norm coordinate}\label{subsec:v12-norm-elim}

One further elimination uses coordinates already present in the
order-forty test $f(H)$.  With $f(H)=h_0+h_1\rho$, its square is stored
through
\[
 t=h_0^2+2h_1^2,\qquad u=h_0h_1.
\]
Here $u$ is half the $\rho$-coefficient, a normalization that matters.
The scalar coordinate $q=h_0^2-2h_1^2$ from
\eqref{eq:v12-scalar-eight} has only one other consumer, the equation
$r(A-1)=q^2$.  Use the identity
\begin{equation}\label{eq:v12-norm-square}
 (h_0^2-2h_1^2)^2=(h_0^2+2h_1^2)^2-8(h_0h_1)^2=t^2-8u^2.
\end{equation}
Delete $q$ and its defining row and replace the consumer by
\[
 r(A-1)=t^2-8u^2.
\]
This is still quadratic.  Forward, forget $q$.  Backward, reinsert
$q=h_0^2-2h_1^2$.  Identity \eqref{eq:v12-norm-square} proves both
rows in the old system and makes these maps mutually inverse.
No division, nonzero hypothesis, or extra arithmetic input is used.
The proof of Lemma~\ref{lem:v12-scalar-sound} applies by this unique
reinsertion; it does not require retaining $q$ as a quantified variable.

\subsection{Soundness of the combined construction}
\label{subsec:v12-soundness}

Let a rational solution of the final system be given.  The inherited
odd charts identify the four root points with odd integer labels.
The reciprocal chart supplies the actual $R=S-P$, is nonidentity,
and, together with the anchor, excludes $S=-P$.  The explicit maps
above therefore restore all unchanged point domains and quotient
guards needed by the earlier arguments.  The removed scalar norm
coordinate also has its unique rational value.

Put $\kappa=|s|=|a-1|$.  By
\eqref{eq:v12-anchor-subgroup} and Lemma~\ref{lem:v12-scalar-sound},
\[
 \log_2\kappa>25165823.5,\qquad
 \operatorname h(\beta)\le\mathcal H\kappa^{19/32}.
\]
These are the exact bounds consumed in the Version 9 height argument.
For completeness we repeat the numerical deductions.  Since
\[
 10(7+2093966)<25165823.5,\qquad80<2^7,
\]
we have
\[
 \kappa>(80\mathcal H)^{10},\qquad
 |a|\ge\kappa-1\ge\kappa/2>(40\mathcal H)^{10},
\]
and
\begin{equation}\label{eq:v12-margins}
 \operatorname h(\beta)<\kappa/80\le|a|/40,\qquad
 \log(40\mathcal H)<\kappa/80\le|a|/40.
\end{equation}
The first inequality uses
$\kappa^{13/32}>80\mathcal H$; the logarithmic inequality is weaker.
In particular all labels are in the large range needed below.

Set
\[
 F=t^2-a(a-2)tb+(a-2)^2b^2-1,\qquad
 B=\max(|n|,|a|,|t|,|b|)>4.
\]
The triangle inequality gives $|F|<5B^4$.  If $F\ne0$,
Lemma~\ref{lem:v12-mixed-local} and the unchanged detector implication
make the three contact radicals divide $F$ to order forty.
Equation~\eqref{eq:v12-height-input} then gives
\[
 B^{1/10}<7\mathcal H\,5^{19/40}+6<24\mathcal H,
\]
contradicting $B\ge|a|>(40\mathcal H)^{10}$.  Hence $F=0$.
The point $Q_4$ is nonidentity, so its integer label $b\ne0$.
Lemma~\ref{lem:o8-lucas} therefore gives $\log B>|a|-3$.

If $n\ne\beta$, let $N$ be the nonzero reduced numerator of
$n-\beta$.  Then $|N|\le2e^{\operatorname h(\beta)}B$, and the same
three radicals divide $N$ to order twelve.  Using
$B^2-6>(5/8)B^2$ for $B>4$, equation~\eqref{eq:v12-height-input}
gives
\[
 B^{5/12}<40\mathcal H e^{(19/12)\operatorname h(\beta)}.
\]
Consequently
\[
 \log B<\frac{12}{5}\log(40\mathcal H)
            +\frac{19}{5}\operatorname h(\beta)
       <\frac{31}{200}|a|<|a|-3,
\]
a contradiction.  Thus $n=\beta\in\Z$.  If $X=u/v$ in lowest terms
with $v>0$, integrality of $2X^2+3$ forces $v^2\mid2$, so $v=1$.
Every rational solution therefore projects to an integer input.

\subsection{Completeness with all replacements at once}
\label{subsec:v12-completeness}

Fix an integer $X$ and put $n=\beta=2X^2+3$, a positive odd integer.
First perform the scalar preparation with the original six-leaf
bounds, which, as shown above, suffice for the new scalar tree.
Choose its finite correction and a sufficiently divisible positive
shared index $h_0$, with $64\mid h_0$.  Choose a positive scalar index
$s$ still more divisible, large enough for the physical anchor, and
with $6\mid s$.  Put $a=s+1$.  The actual prepared scalar point is
$R=[s]P=S-P$.  The new anchor has exactly the old allowed subgroup,
and the new scalar tree has the required unit root.  Fresh scalar
pole witnesses therefore exist.

To choose the two other root labels, define
\[
 U_0=0,\quad U_1=1,\quad U_{j+1}=aU_j-U_{j-1},\qquad
 k=2(a-2),\quad b=U_k/(a-2),\quad t=U_{k+1}.
\]
Induction gives $U_j\equiv j\pmod{a-2}$, so $b$ is an integer.
Modulo two the sequence has period $0,1,1$.  Since $a\equiv1\pmod3$,
$k\equiv1\pmod3$, and both $U_k,U_{k+1}$ are odd.  The positive odd
denominator $a-2$ makes $b$ odd.  The Lucas identity gives
$F(a,t,b)=0$.  Thus all four root indices are positive and odd, and
Proposition~\ref{v9:prop:odd-chart} supplies their chart witnesses.

Fix these root points and the three positive rational square scales
of Corollary~\ref{v9:cor:scales}.  The new scalar functions and the
mixed alternate numerator are rational functions defined at the
prepared points.  Their required finite local conditions persist in
a sufficiently small neighborhood.  As in the original argument,
choose a common continuity modulus $L$, with $64\mid L$, so that
$h\equiv h_0\pmod L$ preserves them.  All these choices precede the
contact-prime selection.  Increase the finite prime cutoff to include
the fixed root denominators needed for the five half-trace guards.

Choose the contact primes using the same simultaneous-prime theorem
and Lemma~\ref{v9:lem:separation}.  Put
\[
 M_\pm=\operatorname{lcm}_p(d_\pm(p)/4),\qquad\Pi=\prod_p p.
\]
The retained direct-target construction supplies
\[
 \gcd(M_-,M_+)=1,\quad\gcd(M_-,L)=16,\quad\gcd(M_+,L)=1,
 \quad\gcd(\Pi,M_+)=\gcd(\Pi,L)=1.
\]
Thus the same three congruences are compatible:
\[
 h\equiv0\pmod{\Pi^{19}M_-},\qquad
 h\equiv1\pmod{M_+},\qquad h\equiv h_0\pmod L.
\]
Use the inherited real-progression argument to choose a sufficiently
large positive $h$ in this progression, avoiding the finitely many
forbidden point and quotient loci.  Set
$U_i=[hn_i]P$ and $V_i=[(h-1)n_i]P$.

At every primary place the points $H,U_i$ have depth at least twenty:
the doubling-image order first makes them formal, and multiplication
by $p^{19}$ raises depth by nineteen under the same formal logarithm
isometry.  At opposite places the points $V_i$ are formal.  Hence
$L_i$ agrees with its integer label to order at least eighty; the
alternate ratio agrees to order twelve for $i=1$ and forty for the
other three labels.  All tests \eqref{eq:v12-tests} hold.  The same
primary/opposite reductions prove that all five half-traces are
nonzero, by the argument following \eqref{eq:v12-half-trace}.
The chosen scalar label has $S\ne\pm P$, so its reciprocal chart
has the ordinary rational lift described above.

The pole-realization estimates are valid for these fixed orders
$1,12,40$.  In the source notation their large-real-parameter decay
is still $\lambda^{-3}$; the small-prime lower bound is
$3\nu+(10m-8)\eta-2c$, which tends to infinity under the same
preparatory divisibility.  No new prime-pattern assertion is used.
Obtain the detector pole witnesses and apply the rational
normalization \eqref{eq:v12-pole-normalized}.  Apply the weight-circuit
maps \eqref{eq:v12-toric}, the reciprocal denominator cancellation,
and the norm-square elimination.  These exact maps preserve all
outside coordinates.  They therefore give one simultaneous rational
solution of the final system for the chosen integer $X$.

This proves completeness in the required order of choices.  In
particular, it does not combine separately chosen primes or
incompatible continuity moduli from different optimizations.

\subsection{Exact polynomial, counts, and coefficient lift}
\label{subsec:v12-finite}

Each equation above is represented by its rational component rows;
intermediate products in the retained straight-line program are
represented by their specified coordinates.  Normalize each nonzero
rational residual by clearing denominators and dividing its integer
content, and let $P_{498}$ be the sum of the squares of those 467
normalized residuals.  This is an integer polynomial.  A sum of
rational squares vanishes exactly when every residual vanishes.
The soundness and completeness just proved therefore establish
Theorem~\ref{thm:v12-quartic}.

The following count records the actual transported, fully reduced
systems; it does not count unreduced formulas in isolation.
\begin{center}\small
\begin{tabular}{lrr}
\toprule
System or replacement & Total variables & Quadratic residuals\\
\midrule
The retained Version 9 baseline &523&492\\
Mixed input precision and its checked compilation &519&488\\
Three normalized detector poles &516&485\\
Five half-trace guards &511&480\\
Single-coefficient translated anchor &510&479\\
Shorter scalar norm tree &505&474\\
Scalar reciprocal chart with $d^2$ retained &504&473\\
Three weight circuits retaining $q_{12},q_{40}$ &501&470\\
Cancellation of the scalar $d^2$ coordinates &499&468\\
Norm-square coordinate elimination &498&467\\
\bottomrule
\end{tabular}
\end{center}

The mixed-depth compilation uses only exact quadratic identities,
constant-pivot substitutions, and invertible shears in shared
coefficient coordinates.  A constant-pivot substitution inserts its
unique polynomial value in the reverse direction; a shear inserts
its inverse linear transformation.  Both maps fix the input.
The final norm-square elimination is the explicit identity
\eqref{eq:v12-norm-square}.  Thus no compiler reduction introduces a
variable denominator or alters an arithmetic guard.  The exact
certificates also check that private coordinates deleted in the
other replacements have no undeclared consumers.  In particular,
the half-trace witnesses occur only in their guard equations; the
removed weight powers occur only in their old weight circuit; the
removed scalar $d^2$ coordinates occur only in their definitions and
$k$ equations; and the removed norm occurs only in its definition and
constant-node equation.

The frozen residual file is \texttt{candidate498-ir.json}, with
SHA-256
\begin{center}\small
\nolinkurl{7682a1229149c75e27322e86dec424449b55321b06c48e468f90fa78bcebb61a}.
\end{center}
The complete sparse coefficient table \texttt{polynomial498.json}
has SHA-256
\begin{center}\small
\nolinkurl{817fec936dee3b358fc8b5b0b9a5570694adf009b134432467ef0be914aeedc4}.
\end{center}
Its exact expansion has 24,669 monomials and degree four; all 498
variables occur.  A coefficient of absolute value one proves content
one.  The only quadratic-in-$C$ coefficient is $16C^2$, at the square
of the translated anchor witness $w$, with no lower-$C$ term at that
monomial.  All other coefficients are independent of $C$ or linear
in $C$ and are smaller in absolute value.  The coefficient comparison
is carried out symbolically; no expansion of $3^{2^{25}}$ is needed.
It gives the exact maximum $16C^2=16\cdot3^{2^{26}}$.

There is also an ordinary-coefficient lift with sixteen additional
coordinates:
\[
 c_0=3^{1024},\qquad c_{j+1}=c_j^2\quad(0\le j<15),\qquad C=c_{15}.
\]
The recurrence has exactly one rational solution and ends at
$3^{1024\cdot2^{15}}=3^{2^{25}}$.  Replace $-4Cw$ in
\eqref{eq:v12-anchor} by $-4c_{15}w$.  All residuals remain quadratic.
This gives a specified quartic in 514 variables with the same input
projection and 513 existential witnesses.  It has 483 quadratic
residuals and 24,716 monomials, with content one and maximum coefficient
$3^{2048}+1005072314676622<2^{3247}$.  Its coefficient table,
\texttt{polynomial514.json}, has SHA-256
\begin{center}\small
\nolinkurl{d247a27709afda5f2be626390cffa1a175c3dfcbee1a5c832a1c457b32629bd5}.
\end{center}
The lift is unique, so it requires no additional arithmetic argument.

The exact checks establish the residual identities, the witness
counts, and every coefficient of the displayed sums of squares.
They are separate from the arithmetic arguments above.  The
inherited source qualifications of Section~\ref{subsec:v9-inputs}
remain in force; a finite polynomial check is not a new proof of the
restricted converse, the height theorem, or the simultaneous-prime
input.

